% Lesson 3 — Vocabulary: Words and expressions related to undergoing a job interview — Anglais Tle A — Cours signé OnBuch+
% Programme officiel MINESEC. Compiler avec : tectonic EN03-vocabulary-words-and-expressions-related-to-undergoing.tex
\newcommand{\DOCMATIERE}{Anglais}
\newcommand{\DOCNIVEAU}{Tle A}
\newcommand{\DOCMODULE}{Chapter 1 : Family and Social Life}
\newcommand{\DOCLECON}{EN03}
\newcommand{\DOCTITRE}{Lesson 3 — Vocabulary: Words and expressions related to undergoing a job interview}
% OnBuch+ — préambule « cours signés » ENRICHI (compilé avec tectonic / XeLaTeX)
% Système visuel « Pop » d'OnBuch+ : fond crème (jamais blanc), bordures encre
% épaisses, ombres « dures » décalées, titres Archivo Black en capitales.
% Version MATHÉMATIQUES : graphes (pgfplots/tikz), boîtes pédagogiques
% (définition, propriété, méthode, attention, le savais-tu, exercice résolu,
% exercice, TP, bilan), page de garde et signature OnBuch+ ancrée en bas.
%
% Macros attendues dans main.tex AVANT \input{preamble} :
%   \DOCMATIERE  \DOCNIVEAU  \DOCMODULE  \DOCLECON (numéro)  \DOCTITRE
\documentclass[11pt]{article}

\usepackage{fontspec}
\usepackage[english,french]{babel} % français = langue principale ; l'anglais via \eng{..} / textebox
\usepackage[a4paper,margin=2cm,top=3.4cm,bottom=2.8cm,headheight=1.4cm,footskip=1.3cm]{geometry}
\usepackage{amsmath,amssymb,amsfonts}
\usepackage{mathtools} % \xmapsto, \coloneqq... souvent utilisés par les modèles
\usepackage{cancel} % simplifications barrées (bilans de forces)
% \abs{z} (module, valeur absolue) et \norm{u} (norme), utilisés par les modèles
\providecommand{\abs}{}\let\abs\relax\DeclarePairedDelimiter{\abs}{\lvert}{\rvert}
\providecommand{\norm}{}\let\norm\relax\DeclarePairedDelimiter{\norm}{\lVert}{\rVert}
\usepackage[table]{xcolor} % [table] : fournit aussi \rowcolors (alternance de lignes)
\usepackage{pagecolor}
\usepackage{tikz}
\usetikzlibrary{arrows.meta,positioning,calc,shapes.geometric,shapes.misc,shapes.arrows,decorations.pathmorphing,decorations.pathreplacing,decorations.markings,patterns,shadows,mindmap,trees,fit,backgrounds,matrix,intersections,angles,quotes}
\usepackage{pgfplots}
\usepackage[version=4]{mhchem} % équations nucléaires \ce{^{235}_{92}U}
\usepackage[european,siunitx]{circuitikz} % schémas électriques
\pgfplotsset{compat=1.18}
\usepgfplotslibrary{fillbetween,groupplots} % aires entre courbes (calcul intégral)
\usepackage{enumitem}
\usepackage{fancyhdr}
% [most] charge toutes les bibliothèques tcolorbox (listings, minted, raster,
% documentation...) et gonfle inutilement la mémoire TeX sur un document long
% (~300 boîtes) : on ne charge que ce qui est réellement utilisé.
\usepackage[skins,breakable]{tcolorbox}
\usepackage{graphicx}
\usepackage{colortbl}
\usepackage{booktabs}
\usepackage{tabularx}
\usepackage{diagbox} % case d'en-tête diagonale des tableaux à double entrée
% Colonnes extensibles centrée / gauche / droite (C, L, R), souvent utilisées par les modèles
\newcolumntype{C}{>{\centering\arraybackslash}X}
\newcolumntype{L}{>{\raggedright\arraybackslash}X}
\newcolumntype{R}{>{\raggedleft\arraybackslash}X}
\usepackage{array}
\usepackage{multicol}
\usepackage{multirow}
\usepackage{siunitx}
% parse-numbers=false : accepte aussi \SI{2,5 \times 10^{-3}}{...}
\sisetup{locale=FR,output-decimal-marker={,},group-minimum-digits=4,parse-numbers=false}
\DeclareSIUnit\ohm{\ensuremath{\symup{\Omega}}} % U+2126 absent de la police
\DeclareSIUnit\division{div} % réglages d'oscilloscope (ms/div, V/div)
\DeclareSIUnit\tr{tr} % tours (vitesse de rotation, ex. \SI{1500}{\tr\per\minute})
\DeclareSIUnit\molar{M} % concentration molaire (ex. \SI{3}{\molar}, \SI{1}{\micro\molar})
\DeclareSIUnit\an{an} % année (le modèle confond parfois avec \year, un compteur TeX) — ex. \SI{5}{\kilo\meter\per\an}
\DeclareSIUnit\FCFA{FCFA} % franc CFA (ex. \SI{500}{\FCFA\per\kilo\gram})
\DeclareSIUnit\degreeBrix{\textdegree Brix} % teneur en sucre d'un sirop/jus (réfractométrie)
\DeclareSIUnit\month{mois} % le modèle utilise parfois \month, un compteur TeX, comme unité de durée
\usepackage{qrcode}
% \degree : commande fréquemment utilisée par les modèles pour un angle ou une
% température en dehors de \SI{}{} (non fournie par les paquets chargés).
\providecommand{\degree}{^{\circ}}
% \qed (fin de démonstration), utilisé par les modèles sans amsthm
\providecommand{\qed}{\hfill$\square$}
% \barre : double barre des valeurs interdites dans un tableau de variations (tkz-tab)
\providecommand{\barre}{\ensuremath{\|}}
% environnement proof (amsthm non chargé) utilisé par les modèles
\ifdefined\proof\else\newenvironment{proof}[1][Démonstration]{\par\noindent\textit{#1.}\ }{\qed\par}\fi
\usepackage{lastpage}
\usepackage{hyperref}
\hypersetup{hidelinks}

% ============================================================
% Palette « Pop » OnBuch+ — mêmes hex que la classe P (plus_ui.dart)
% ============================================================
\definecolor{poporange}{HTML}{F59321}
\definecolor{popdark}{HTML}{E07A0C}
\definecolor{popcream}{HTML}{FFF6E8}
\definecolor{popink}{HTML}{1C1714}
\definecolor{poppurple}{HTML}{7A5AE0}
\definecolor{popgreen}{HTML}{2E9E6A}
\definecolor{poppink}{HTML}{E0457B}
\definecolor{popblue}{HTML}{2F7DE1}
\definecolor{popgold}{HTML}{C98A12}
\definecolor{popmuted}{HTML}{8A7F76}
\definecolor{popline}{HTML}{EADFD2}
\definecolor{popgray}{HTML}{8A7F76}
\colorlet{popgrey}{popgray}
% Alias tolérants (le modèle nomme parfois les couleurs différemment)
\colorlet{popwhite}{white}
\colorlet{popblack}{popink}
\colorlet{popred}{poppink}
\colorlet{popyellow}{popgold}
% Teintes claires (fonds de tableaux/encadrés) — le modèle nomme parfois
% les couleurs avec un suffixe L (ex. popblueL) sans les avoir définies.
\colorlet{poporangeL}{poporange!20}
\colorlet{popdarkL}{popdark!20}
\colorlet{poppurpleL}{poppurple!20}
\colorlet{popgreenL}{popgreen!20}
\colorlet{poppinkL}{poppink!20}
\colorlet{popblueL}{popblue!20}
\colorlet{popgoldL}{popgold!20}
\colorlet{popredL}{popred!20}
\colorlet{popgrayL}{popgray!20}
% Teintes claires pour fonds de tableaux / zones
\colorlet{popblueL}{popblue!10!white}
\colorlet{poporangeL}{poporange!14!white}
\colorlet{popgreenL}{popgreen!12!white}
\colorlet{poppurpleL}{poppurple!12!white}
\colorlet{poppinkL}{poppink!12!white}
\colorlet{popgoldL}{popgold!14!white}

\pagecolor{popcream}

% ============================================================
% Typographie
% ============================================================
\defaultfontfeatures{Ligatures=TeX}
\setmainfont{PlusJakartaSans-Medium.ttf}[Path=../../../fonts/,BoldFont=PlusJakartaSans-Bold.ttf]
% Symboles, API et emojis absents de la police principale : repli sur DejaVu Sans (caractères actifs)
\newfontface\symfont{DejaVuSans.ttf}[Path=../../../fonts/]
\makeatletter
\newcommand\symactive[1]{\catcode#1=13 \begingroup\lccode`\~=#1 \lowercase{\endgroup\def~}{{\symfont\char#1}}}
\newcommand\symblank[1]{\catcode#1=13 \begingroup\lccode`\~=#1 \lowercase{\endgroup\def~}{}}
\newcommand\symhyph[1]{\catcode#1=13 \begingroup\lccode`\~=#1 \lowercase{\endgroup\def~}{\mbox{-}}}
\newcount\symc
\newcommand\symrange[2]{\symc=#1 \@whilenum\symc<#2 \do{\expandafter\symactive\expandafter{\the\symc}\advance\symc by 1 }}
\newcommand\symblankrange[2]{\symc=#1 \@whilenum\symc<#2 \do{\expandafter\symblank\expandafter{\the\symc}\advance\symc by 1 }}
\makeatother
\symrange{"0250}{"02FF}\symactive{"2713}\symactive{"2714}\symactive{"2717}\symactive{"2718}\symactive{"2610}\symactive{"2611}\symactive{"2612}
\symhyph{"2011}\symhyph{"2E1A}\symblankrange{"1F300}{"1FAFF}\symblank{"FE0F}
% Maths en sans-serif (Fira Math) avec chiffres et lettres droites de Plus
% Jakarta, pour que les formules se fondent dans le texte.
\usepackage[math-style=ISO,bold-style=ISO]{unicode-math}
\setmathfont{FiraMath-Regular.otf}[Path=../../../fonts/]
\setmathfont{PlusJakartaSans-Medium.ttf}[Path=../../../fonts/,range={up/num,up/{Latin,latin},\mathup}]
\usepackage{icomma} % virgule décimale sans espace en mode math
% Caractères mathématiques Unicode absents des polices de texte (Plus Jakarta,
% Archivo Black) : rendus par leur équivalent mathématique, en texte comme en
% formule — sinon ils s'affichent en carré vide (ex. « Arithmétique dans ℤ »).
\usepackage{newunicodechar}
\newunicodechar{ℝ}{\ensuremath{\mathbb{R}}}
\newunicodechar{ℤ}{\ensuremath{\mathbb{Z}}}
\newunicodechar{ℂ}{\ensuremath{\mathbb{C}}}
\newunicodechar{ℕ}{\ensuremath{\mathbb{N}}}
\newunicodechar{ℚ}{\ensuremath{\mathbb{Q}}}
\newunicodechar{𝔻}{\ensuremath{\mathbb{D}}}
\newunicodechar{α}{\ensuremath{\alpha}}
\newunicodechar{β}{\ensuremath{\beta}}
\newunicodechar{γ}{\ensuremath{\gamma}}
\newunicodechar{δ}{\ensuremath{\delta}}
\newunicodechar{ε}{\ensuremath{\varepsilon}}
\newunicodechar{θ}{\ensuremath{\theta}}
\newunicodechar{λ}{\ensuremath{\lambda}}
\newunicodechar{μ}{\ensuremath{\mu}}
\newunicodechar{σ}{\ensuremath{\sigma}}
\newunicodechar{τ}{\ensuremath{\tau}}
\newunicodechar{φ}{\ensuremath{\varphi}}
\newunicodechar{ψ}{\ensuremath{\psi}}
\newunicodechar{ω}{\ensuremath{\omega}}
\newunicodechar{Σ}{\ensuremath{\Sigma}}
% majuscules produites par \MakeUppercase dans les titres (α-aminés -> Α)
\newunicodechar{Α}{\ensuremath{\alpha}}
\newunicodechar{Β}{\ensuremath{\beta}}
\newunicodechar{Γ}{\ensuremath{\Gamma}}
\newunicodechar{Δ}{\ensuremath{\Delta}}
\newunicodechar{Ε}{\ensuremath{\varepsilon}}
\newunicodechar{Θ}{\ensuremath{\Theta}}
\newunicodechar{Λ}{\ensuremath{\Lambda}}
\newunicodechar{Μ}{\ensuremath{\mu}}
\newunicodechar{Τ}{\ensuremath{\tau}}
\newunicodechar{Φ}{\ensuremath{\Phi}}
\newunicodechar{Ψ}{\ensuremath{\Psi}}
\newunicodechar{Ω}{\ensuremath{\Omega}}
\newunicodechar{↦}{\ensuremath{\mapsto}}
\newunicodechar{∈}{\ensuremath{\in}}
\newunicodechar{∉}{\ensuremath{\notin}}
\newunicodechar{∧}{\ensuremath{\wedge}}
\newunicodechar{⇔}{\ensuremath{\Leftrightarrow}}
\newunicodechar{⟺}{\ensuremath{\Longleftrightarrow}}
\newunicodechar{⇒}{\ensuremath{\Rightarrow}}
\newunicodechar{⟹}{\ensuremath{\Longrightarrow}}
\newunicodechar{∘}{\ensuremath{\circ}}
\newunicodechar{∪}{\ensuremath{\cup}}
\newunicodechar{∩}{\ensuremath{\cap}}
\newunicodechar{≡}{\ensuremath{\equiv}}
\newunicodechar{⊂}{\ensuremath{\subset}}
\newunicodechar{⊥}{\ensuremath{\perp}}
\newunicodechar{∃}{\ensuremath{\exists}}
\newunicodechar{∀}{\ensuremath{\forall}}
\newunicodechar{∛}{\ensuremath{\sqrt[3]{\,}}}
\newunicodechar{∜}{\ensuremath{\sqrt[4]{\,}}}
\newunicodechar{⃗}{\ensuremath{{}^{\to}}}
\newunicodechar{ⁿ}{\ifmmode^{n}\else\textsuperscript{\ensuremath{n}}\fi}
\newunicodechar{ˣ}{\ifmmode^{x}\else\textsuperscript{\ensuremath{x}}\fi}
\newunicodechar{ᵇ}{\ifmmode^{b}\else\textsuperscript{\ensuremath{b}}\fi}
\newunicodechar{ʳ}{\ifmmode^{r}\else\textsuperscript{\ensuremath{r}}\fi}
\newunicodechar{ᵘ}{\ifmmode^{u}\else\textsuperscript{\ensuremath{u}}\fi}
\newunicodechar{ʸ}{\ifmmode^{y}\else\textsuperscript{\ensuremath{y}}\fi}
\newunicodechar{ᵃ}{\ifmmode^{a}\else\textsuperscript{\ensuremath{a}}\fi}
\newunicodechar{ᵉ}{\ifmmode^{e}\else\textsuperscript{\ensuremath{e}}\fi}
\newunicodechar{ᵏ}{\ifmmode^{k}\else\textsuperscript{\ensuremath{k}}\fi}
\newunicodechar{ᵖ}{\ifmmode^{p}\else\textsuperscript{\ensuremath{p}}\fi}
\newunicodechar{ᴺ}{\ifmmode^{N}\else\textsuperscript{\ensuremath{N}}\fi}
\newunicodechar{ᶜ}{\ifmmode^{c}\else\textsuperscript{\ensuremath{c}}\fi}
\newunicodechar{ʰ}{\ifmmode^{h}\else\textsuperscript{\ensuremath{h}}\fi}
\newunicodechar{ᵠ}{\ifmmode^{q}\else\textsuperscript{\ensuremath{q}}\fi}
\newunicodechar{ₙ}{\ifmmode_{n}\else\textsubscript{\ensuremath{n}}\fi}
\newunicodechar{ₐ}{\ifmmode_{a}\else\textsubscript{\ensuremath{a}}\fi}
\newunicodechar{ᵢ}{\ifmmode_{i}\else\textsubscript{\ensuremath{i}}\fi}
\newunicodechar{ᵣ}{\ifmmode_{r}\else\textsubscript{\ensuremath{r}}\fi}
\newunicodechar{ₖ}{\ifmmode_{k}\else\textsubscript{\ensuremath{k}}\fi}
\newunicodechar{ᵦ}{\ifmmode_{\beta}\else\textsubscript{\ensuremath{\beta}}\fi}
\newunicodechar{ₓ}{\ifmmode_{x}\else\textsubscript{\ensuremath{x}}\fi}
\newfontfamily\popdisplay{ArchivoBlack-Regular.ttf}[Path=../../../fonts/]
\newfontfamily\poplogo{SpaceGrotesk-Bold.ttf}[Path=../../../fonts/]
\newfontfamily\popsemi{PlusJakartaSans-SemiBold.ttf}[Path=../../../fonts/]
\newfontfamily\popxbold{PlusJakartaSans-ExtraBold.ttf}[Path=../../../fonts/]
\newfontfamily\popxboldls{PlusJakartaSans-ExtraBold.ttf}[Path=../../../fonts/,LetterSpace=3.0]

\setlist[enumerate]{leftmargin=1.7em,itemsep=5pt,topsep=4pt}
\setlist[itemize]{leftmargin=1.7em,itemsep=4pt,topsep=4pt,label={\color{poporange}\textbullet}}
\setlength{\parindent}{0pt}
\setlength{\parskip}{5pt}
\renewcommand{\arraystretch}{1.25}

% pgfplots : style Pop par défaut
\pgfplotsset{
  every axis/.append style={axis line style={popink,line width=1pt},tick style={popink},
    grid style={popline,line width=0.5pt},label style={font=\small},tick label style={font=\footnotesize}},
  popcurve/.style={poporange,line width=1.8pt,no markers,smooth},
}
% Même style disponible dans un \draw TikZ simple (hors pgfplots)
\tikzset{popcurve/.style={poporange,line width=1.8pt}}
\tikzset{popgrid/.style={popline,very thin}}
\colorlet{popcurve}{poporange} % le nom du style est parfois utilisé comme couleur

% ============================================================
% Pill / squiggle
% ============================================================
\newcommand{\poppill}[3]{%
  \tikz[baseline=-0.55ex]\node[draw=popink,line width=1.2pt,rounded corners=2.6pt,%
    fill=#2,inner xsep=6.5pt,inner ysep=3pt]{%
    \popxboldls\fontsize{7}{7}\selectfont\color{#3}\MakeUppercase{#1}};%
}
\newcommand{\popsquiggle}[1]{%
  \tikz[baseline=-0.4ex]{
    \draw[#1,line width=2.6pt,line cap=round]
      plot [smooth] coordinates {(0,0.09) (0.34,0.01) (0.68,0.21) (1.02,0.01) (1.36,0.10)};
  }%
}
% Mot-clé surligné (marqueur orange sous le texte)
\newcommand{\cle}[1]{{\popxbold\color{popdark}#1}}

% ============================================================
% En-tête / pied
% ============================================================
\pagestyle{fancy}
\fancyhf{}
\renewcommand{\headrulewidth}{0pt}
\renewcommand{\footrulewidth}{0pt}
\lhead{\raisebox{-0.15ex}{\poplogo\fontsize{13}{13}\selectfont\color{popink}OnBuch\color{poporange}+}}
\rhead{\poppill{\DOCMATIERE\ \textbullet\ \DOCNIVEAU\ \textbullet\ Leçon \DOCLECON}{white}{popink}}
\lfoot{\popxbold\fontsize{7.6}{7.6}\selectfont\color{popmuted}\MakeUppercase{Cours signé OnBuch+}\ \textbullet\ {\popsemi\DOCTITRE}}
\rfoot{\tikz[baseline=-0.6ex]\node[rounded corners=8.5pt,fill=popink,minimum height=17pt,minimum width=17pt,inner xsep=5pt,inner sep=0]{\popxbold\fontsize{8}{8}\selectfont\color{popcream}\thepage\,/\,\pageref*{LastPage}};}

% ============================================================
% Titres
% ============================================================
\newcommand{\chaptitre}[2]{%
  \begin{tcolorbox}[enhanced,colback=poporange,colframe=popink,boxrule=2.6pt,arc=11pt,
      drop shadow=popink,left=15pt,right=15pt,top=12pt,bottom=11pt,before skip=3pt,after skip=18pt]
    \poppill{#2}{popink}{popcream}\par\vspace{9pt}
    {\raggedright\hyphenpenalty=10000\exhyphenpenalty=10000\popdisplay\fontsize{22}{24}\selectfont\color{popcream}\MakeUppercase{#1}\par}
  \end{tcolorbox}
}
% Section numérotée : pastille orange avec le numéro + titre Archivo
\newcounter{popsec}
\newcommand{\coursec}[1]{%
  \stepcounter{popsec}\setcounter{popsub}{0}%
  \par\vspace{16pt}\noindent
  \begin{minipage}{\linewidth}
  \tikz[baseline=-0.7ex]\node[draw=popink,line width=1.6pt,fill=poporange,rounded corners=4pt,minimum size=22pt,inner sep=0]{\popdisplay\fontsize{12}{12}\selectfont\color{popcream}\thepopsec};%
  \hspace{8pt}{\popdisplay\fontsize{15}{17}\selectfont\color{popink}\MakeUppercase{#1}}\par
  \vspace{4pt}\noindent{\color{poporange}\rule{36pt}{3.6pt}}
  \end{minipage}\par\nopagebreak\vspace{8pt}
}
\newcounter{popsub}
\newcommand{\courssub}[1]{%
  \stepcounter{popsub}%
  \par\vspace{9pt}\noindent{\popxbold\fontsize{11.5}{13}\selectfont\color{popink}{\color{poporange}\thepopsec.\thepopsub}\hspace{6pt}#1}\par\nopagebreak\vspace{4pt}
}

% ============================================================
% Boîtes pédagogiques — même recette : fond blanc, bordure épaisse colorée,
% ombre dure encre, badge plein. Titre optionnel [#1].
% ============================================================
\tcbset{popbase/.style={breakable,enhanced,colback=white,boxrule=2.3pt,arc=9pt,
  shadow={3pt}{-3pt}{0pt}{popink},left=14pt,right=14pt,top=11pt,bottom=11pt,before skip=11pt,after skip=15pt,
  fontupper=\popsemi\fontsize{10.6}{14.5}\selectfont\color{popink},
  fontlower=\popsemi\fontsize{10.6}{14.5}\selectfont\color{popink}}}
% Coche dessinée (le glyphe ✓ n'existe pas dans les polices du document)
\newcommand{\popcheck}[1]{\tikz[baseline=-0.5ex]\draw[#1,line width=1.6pt,line cap=round,line join=round](-0.13,0.0)--(-0.04,-0.09)--(0.13,0.1);}
\providecommand{\coche}{\popcheck{popgreen}}
\providecommand{\resultat}[1]{\par\smallskip\noindent\fcolorbox{popgreen}{popgreenL}{\parbox{\dimexpr\linewidth-2\fboxsep-2\fboxrule\relax}{\textbf{Résultat.} #1}}\par\smallskip}
\newcommand{\popboxhead}[3]{\poppill{#1}{#2}{white}\ifx\relax#3\relax\else\hspace{6pt}{\popxbold\fontsize{10.6}{12}\selectfont\color{popink}#3}\fi\par\vspace{7pt}}

\newtcolorbox{aretenir}[1][]{popbase,colframe=popgreen,before upper={\popboxhead{À retenir}{popgreen}{#1}}}
% Texte en anglais (document de lecture, dialogue, extrait) : cadre
% d'étude. Un titre optionnel (source, auteur) s'affiche dans l'en-tête.
\newtcolorbox{textebox}[1][]{popbase,colframe=popblue,before upper={\popboxhead{Text}{popblue}{#1}\begin{otherlanguage}{english}},after upper={\end{otherlanguage}}}
% Marques ✓ / ✗ (forme correcte / fautive) souvent utilisées par les modèles
\providecommand{\cross}{\textcolor{poppink}{\ensuremath{\times}}}
\providecommand{\xmark}{\cross}\providecommand{\crossmark}{\cross}\providecommand{\cmark}{\ensuremath{\checkmark}}
% Ligne de réponse / justification à compléter (inventée par les modèles)
\providecommand{\justification}[1]{\par\noindent\textit{Justification~:} \dotfill\par}
% \textipa (package tipa non chargé) : la police ne couvre pas l'API -> simple passage
\providecommand{\textipa}[1]{#1}
% Soulignement (ulem non chargé) : \uline, \ul
\providecommand{\uline}[1]{\underline{#1}}\providecommand{\ul}[1]{\underline{#1}}
% \glossaire{\item ...} inventée par les modèles : liste de glossaire
\providecommand{\glossaire}[1]{\begin{itemize}#1\end{itemize}}
% \alert (beamer) inventée par les modèles : texte mis en évidence
\providecommand{\alert}[1]{\textbf{\textcolor{poppink}{#1}}}
% variantes ASCII d'ordinaux écrites en macro
\providecommand{\iere}{\textsuperscript{re}}\providecommand{\ieme}{\textsuperscript{e}}\providecommand{\ere}{\textsuperscript{re}}\providecommand{\eme}{\textsuperscript{e}}\providecommand{\er}{\textsuperscript{er}}\providecommand{\ie}{\textsuperscript{e}}
% Ordinaux français écrits en macro par les modèles : 1\ière, 3\ième
\newcommand{\ière}{\textsuperscript{re}}\newcommand{\ième}{\textsuperscript{e}}
% \exemple (inventée par les modèles) : étiquette « Exemple : »
\providecommand{\exemple}{\textbf{Exemple~:}\ }
% \square utilisable aussi hors mode math (cases à cocher)
\AtBeginDocument{\let\squareMath\square\renewcommand{\square}{\ensuremath{\squareMath}}}
% Mot ou phrase en anglais à l'intérieur d'un paragraphe français
\newcommand{\eng}[1]{\ifmmode\text{\textit{#1}}\else\begingroup\selectlanguage{english}\itshape #1\endgroup\fi}
\let\esp\eng % alias toléré
% flèches d'intonation inventées par les modèles
\providecommand{\textsearrow}{}\renewcommand{\textsearrow}{\ensuremath{\searrow}}\providecommand{\textnearrow}{}\renewcommand{\textnearrow}{\ensuremath{\nearrow}}
% \fr{..} : mot français (inventé par les modèles)
\providecommand{\fr}[1]{\textit{#1}}
\newtcolorbox{exemplebox}[1][]{popbase,colframe=poporange,before upper={\popboxhead{Exemple}{poporange}{#1}}}
\newtcolorbox{definition}[1][]{popbase,colframe=popblue,before upper={\popboxhead{Définition}{popblue}{#1}}}
\newtcolorbox{propriete}[1][]{popbase,colframe=poppurple,before upper={\popboxhead{Propriété}{poppurple}{#1}}}
\newtcolorbox{methode}[1][]{popbase,colframe=popgold,before upper={\popboxhead{Méthode}{popgold}{#1}}}
\newtcolorbox{attention}[1][]{popbase,colframe=poppink,before upper={\popboxhead{Attention}{poppink}{#1}}}
\newtcolorbox{experience}[1][]{popbase,colframe=popblue,colback=popblueL,before upper={\popboxhead{Expérience}{popblue}{#1}}}
% « Le savais-tu ? » : carte violette pleine (contexte, culture, Cameroun)
\newtcolorbox{savaistu}[1][]{popbase,colback=poppurple,colframe=popink,
  fontupper=\popsemi\fontsize{10.4}{14.2}\selectfont\color{white},
  before upper={\poppill{Le savais-tu\,?}{white}{poppurple}\ifx\relax#1\relax\else\hspace{6pt}{\popxbold\color{white}#1}\fi\par\vspace{7pt}}}
% Exercice résolu : énoncé en haut, \tcblower puis la solution sur fond crème
\newcounter{popexo}\newcounter{popexr}
\newtcolorbox{exoresolu}[1][]{popbase,colframe=poporange,colbacklower=popcream,
  segmentation style={popink,line width=1.2pt,dashed},
  before upper={\stepcounter{popexr}\popboxhead{Exercice résolu \thepopexr}{poporange}{#1}},
  before lower={\poppill{Solution}{popink}{popcream}\par\vspace{6pt}}}
% Exercice (non résolu en place) — la correction est dans la partie Corrigés
\newtcolorbox{exercice}[1][]{popbase,colframe=popink,
  before upper={\stepcounter{popexo}\popboxhead{Exercice \thepopexo}{popink}{#1}}}
% Corrigé
\newtcolorbox{corrige}[1][]{popbase,colframe=popgreen,colback=popgreenL,
  before upper={\popboxhead{Corrigé}{popgreen}{#1}}}
% Objectifs / prérequis / bilan
\newtcolorbox{objectifs}{popbase,colframe=popink,colback=poporangeL,
  before upper={\popboxhead{Objectifs de la leçon}{popink}{}}}
\newtcolorbox{prerequis}{popbase,colframe=popink,colback=popblueL,
  before upper={\popboxhead{Prérequis}{popblue}{}}}
\newtcolorbox{bilan}{popbase,colframe=popink,colback=white,boxrule=2.8pt,
  before upper={\popboxhead{Fiche bilan}{poporange}{L'essentiel à connaître pour l'examen}}}
% Figure encadrée avec légende
\newtcolorbox{popfigure}[1][]{enhanced,breakable=false,colback=white,colframe=popline,boxrule=1.4pt,arc=8pt,
  left=10pt,right=10pt,top=10pt,bottom=8pt,before skip=10pt,after skip=12pt,halign=center,
  lower separated=false,halign lower=center,
  fontlower=\popsemi\fontsize{9}{11.5}\selectfont\color{popmuted},
  #1}
\newcommand{\legende}[1]{\par\vspace{6pt}{\popsemi\fontsize{9}{11.5}\selectfont\color{popmuted}#1}}

% ============================================================
% Signature OnBuch+
% ============================================================
\newcommand{\poplogoinline}[1]{{\poplogo\fontsize{#1}{#1}\selectfont\color{popink}OnBuch\color{poporange}+}}
\newcommand{\signatureonbuch}{%
  \begin{tcolorbox}[enhanced,colback=popink,colframe=popink,boxrule=2.2pt,arc=10pt,
      drop shadow=poporange,left=14pt,right=14pt,top=10pt,bottom=10pt,before skip=0pt,after skip=0pt]
    \tikz[baseline=-0.7ex]\node[circle,fill=popgreen,draw=popcream,line width=1.3pt,minimum size=22pt,inner sep=0]{\popcheck{white}};%
    \hspace{9pt}\begin{minipage}[c]{0.62\linewidth}
      {\popxbold\fontsize{12}{13}\selectfont\color{popcream}Signé\ \textbullet\ {\poplogo OnBuch\color{poporange}+}}\par\vspace{2pt}
      {\raggedright\popsemi\fontsize{8}{10.5}\selectfont\color{popline}\MakeUppercase{Équipe pédagogique OnBuch+}\\\MakeUppercase{Conforme au programme officiel MINESEC}\par}
    \end{minipage}\hfill
    {\poplogo\fontsize{22}{22}\selectfont\color{popcream}OnBuch\color{poporange}+}
  \end{tcolorbox}%
}

% ============================================================
% Page de garde
% #1 titre · #2 matière · #3 niveau · #4 module · #5 n° leçon · #6 durée ·
% #7 description / objectifs (texte ou itemize)
% ============================================================
\newcommand{\pagegarde}[7]{%
  \thispagestyle{empty}
  \newgeometry{a4paper,margin=1.7cm,top=1.6cm,bottom=1.4cm}%
  \begin{tikzpicture}[remember picture,overlay]
    \fill[poporange!18] ([xshift=-3.2cm,yshift=-3.6cm]current page.north east) circle (4.6cm);
    \fill[poppurple!14] ([xshift=2.4cm,yshift=7.6cm]current page.south west) circle (3.8cm);
  \end{tikzpicture}%
  {\poplogo\fontsize{30}{30}\selectfont\color{popink}OnBuch\color{poporange}+}\hfill
  \raisebox{0.6ex}{\poppill{Cours signé \textbullet\ Premium}{popink}{popcream}}\par
  \vspace{1pt}\popsquiggle{poppurple}\par
  \vspace{8pt}
  \begin{tcolorbox}[enhanced,colback=poporange,colframe=popink,boxrule=2.8pt,arc=13pt,
      drop shadow=popink,left=18pt,right=18pt,top=12pt,bottom=13pt,after skip=10pt]
    \poppill{#2 \textbullet\ #3}{popink}{popcream}\hspace{5pt}\poppill{#4}{white}{popink}\par\vspace{8pt}
    {\popdisplay\fontsize{14}{14}\selectfont\color{popink}LEÇON #5}\par\vspace{4pt}
    {\raggedright\hyphenpenalty=10000\exhyphenpenalty=10000\popdisplay\fontsize{25}{27}\selectfont\color{popcream}\MakeUppercase{#1}\par}
    \vspace{8pt}
    {\popxbold\fontsize{9.5}{9.5}\selectfont\color{popink}\MakeUppercase{Durée conseillée : #6}}
  \end{tcolorbox}
  \begin{tcolorbox}[enhanced,colback=white,colframe=popink,boxrule=2.2pt,arc=10pt,
      drop shadow=popink,left=16pt,right=16pt,top=10pt,bottom=10pt,after skip=8pt]
    \poppill{Dans cette leçon}{poppurple}{white}\par\vspace{5pt}
    \popsemi\fontsize{9.3}{12.6}\selectfont\color{popink}#7
  \end{tcolorbox}
  \noindent\begin{minipage}[t]{0.487\linewidth}
    \begin{tcolorbox}[enhanced,colback=popgreen,colframe=popink,boxrule=2.2pt,arc=10pt,
        drop shadow=popink,left=11pt,right=11pt,top=10pt,bottom=10pt,height=3.1cm,valign=center]
      \tcbox[colback=white,colframe=popink,boxrule=1.3pt,arc=5pt,boxsep=0pt,nobeforeafter,
        left=3pt,right=3pt,top=3pt,bottom=3pt,valign=center]{\qrcode[height=1.9cm]{https://play.google.com/store/apps/details?id=com.luvvix.onbuch&hl=fr}}%
      \hspace{8pt}\begin{minipage}[c]{0.5\linewidth}
        {\popdisplay\fontsize{11}{12.5}\selectfont\color{white}\MakeUppercase{Télécharge OnBuch}}\par\vspace{4pt}
        {\raggedright\popsemi\fontsize{8.4}{11}\selectfont\color{white}Gratuit \textbullet\ Google Play\\Tuteur IA, annales, concours, résultats\par}
      \end{minipage}
    \end{tcolorbox}
  \end{minipage}\hfill
  \begin{minipage}[t]{0.487\linewidth}
    \begin{tcolorbox}[enhanced,colback=poppurple,colframe=popink,boxrule=2.2pt,arc=10pt,
        drop shadow=popink,left=11pt,right=11pt,top=10pt,bottom=10pt,height=3.1cm,valign=center]
      \tikz[baseline=-0.7ex]\node[circle,fill=white,draw=popink,line width=1.3pt,minimum size=1.6cm,inner sep=0]{\popdisplay\fontsize{24}{24}\selectfont\color{poppurple}+};%
      \hspace{8pt}\begin{minipage}[c]{0.6\linewidth}
        {\popdisplay\fontsize{11}{12.5}\selectfont\color{white}\MakeUppercase{Passe à OnBuch+}}\par\vspace{4pt}
        {\raggedright\popsemi\fontsize{8.4}{11}\selectfont\color{white}Tous les cours signés, Tuteur IA illimité, fiches PDF — onglet OnBuch+ de l'app\par}
      \end{minipage}
    \end{tcolorbox}
  \end{minipage}
  \vspace{8pt}
  \popcontenu
  \vfill
  \signatureonbuch
  \restoregeometry
  \newpage
}

% Rangée « Ce cours contient » (page de garde)
\newcommand{\poptile}[3]{% #1 couleur #2 gros texte #3 légende
  \begin{tcolorbox}[enhanced,colback=#1,colframe=popink,boxrule=2pt,arc=9pt,shadow={2.5pt}{-2.5pt}{0pt}{popink},
      left=6pt,right=6pt,top=9pt,bottom=9pt,halign=center,valign=center,height=1.75cm,width=\linewidth,nobeforeafter]
    {\popdisplay\fontsize{12.5}{13}\selectfont\color{white}\MakeUppercase{#2}}\par\vspace{3pt}
    {\centering\popsemi\fontsize{7.8}{9.5}\selectfont\color{white}#3\par}
  \end{tcolorbox}}
\newcommand{\popcontenu}{%
  \par\noindent{\popxboldls\fontsize{7.5}{7.5}\selectfont\color{popmuted}CE COURS SIGNÉ CONTIENT}\par\vspace{7pt}
  \noindent\begin{minipage}[t]{0.235\linewidth}\poptile{popblue}{Cours}{complet, illustré, pas à pas}\end{minipage}\hfill
  \begin{minipage}[t]{0.235\linewidth}\poptile{poporange}{Méthodes}{et exercices résolus}\end{minipage}\hfill
  \begin{minipage}[t]{0.235\linewidth}\poptile{poppink}{Exercices}{type examen + corrigés}\end{minipage}\hfill
  \begin{minipage}[t]{0.235\linewidth}\poptile{popgreen}{Fiche bilan}{l'essentiel à retenir}\end{minipage}\par}

% Sommaire
\newtcolorbox{sommaire}{enhanced,colback=white,colframe=popink,boxrule=2.2pt,arc=10pt,
  drop shadow=popink,left=18pt,right=18pt,top=13pt,bottom=13pt,before skip=6pt,after skip=20pt,
  before upper={\poppill{Sommaire}{poppurple}{white}\par\vspace{9pt}\popsemi\fontsize{10.6}{14.5}\selectfont\color{popink}}}

% Signature de fin de cours — poussée tout en bas de la dernière page
\newcommand{\findecours}{%
  \par\vfill\nopagebreak
  \begin{center}\popsquiggle{poporange}\end{center}\vspace{4pt}
  \signatureonbuch
}
\AtBeginDocument{\def\llbracket{\mathopen{[\mkern-3.5mu[}}\def\rrbracket{\mathclose{]\mkern-3.5mu]}}} % intervalles d'entiers (glyphe ⟦ absent de la police)
% Glyphes absents de Fira Math : redéfinis à partir de glyphes disponibles
\AtBeginDocument{%
  \def\setminus{\mathbin{\backslash}}\let\smallsetminus\setminus
  \def\vdots{\mathord{\vbox{\baselineskip3.2pt\lineskiplimit0pt\kern4pt\hbox{.}\hbox{.}\hbox{.}}}}%
  \def\ddots{\mathinner{\mkern1mu\raise6.5pt\hbox{.}\mkern2mu\raise3.6pt\hbox{.}\mkern2mu\raise0.7pt\hbox{.}\mkern1mu}}%
  \def\triangle{\mathord{\scalebox{0.9}{$\symup{\Delta}$}}}%
  \def\nabla{\mathord{\raisebox{\depth}{\scalebox{1}[-1]{$\symup{\Delta}$}}}}%
  \def\checkmark{\popcheck{popdark}}%
  \def\star{\mathbin{\ast}}\def\bigstar{\mathord{\ast}}%
  \def\diamond{\mathbin{\scalebox{0.6}{$\Diamond$}}}%
  \def\bigcup{\mathop{\mathchoice{\raisebox{-0.3em}{\scalebox{1.7}{$\cup$}}}{\scalebox{1.25}{$\cup$}}{\cup}{\cup}}\displaylimits}%
  \def\bigcap{\mathop{\mathchoice{\raisebox{-0.3em}{\scalebox{1.7}{$\cap$}}}{\scalebox{1.25}{$\cap$}}{\cap}{\cap}}\displaylimits}%
  \def\lesseqqgtr{\mathrel{\lessgtr}}\def\bigoplus{\mathop{\oplus}\displaylimits}%
}
\providecommand{\ssi}{si et seulement si} % abréviation fréquente
\AtBeginDocument{\providecommand{\wideparen}{\overparen}} % arc de cercle

%% FIGFIX-BEGIN v3 — correctifs globaux des figures (docs/onbuch-generation/tools/figfix)
\usepackage{adjustbox}\usepackage{etoolbox}
% toute figure TikZ/pgfplots plus large que la ligne est réduite à la largeur disponible
\BeforeBeginEnvironment{tikzpicture}{\begin{adjustbox}{max width=\linewidth}}
\AfterEndEnvironment{tikzpicture}{\end{adjustbox}}
% légendes pgfplots lisibles par-dessus les courbes
\pgfplotsset{every axis/.append style={legend style={fill=white,fill opacity=0.92,text opacity=1,draw=black!15,inner sep=3pt}}}
% étiquettes d'arêtes (syntaxe edge["..."]) avec fond blanc
\tikzset{every edge quotes/.append style={fill=white,inner sep=1.2pt}}
% bulles de cartes mentales : texte à la ligne dans la bulle, bulle un peu plus grande
\tikzset{every concept/.append style={text width=2.0cm,align=center,minimum size=2.3cm,inner sep=2pt}}
%% FIGFIX-END
\begin{document}

\pagegarde{Lesson 3 — Vocabulary: Words and expressions related to undergoing a job interview}{Anglais}{Tle A}{Chapter 1 : Family and Social Life}{EN03}{2 h}{Cette leçon de vocabulaire présente les mots, collocations et expressions liés au déroulement d'un entretien d'embauche, de la convocation à la réponse finale. Elle entraîne les élèves à employer ce lexique en contexte, à l'oral comme à l'écrit, en évitant les faux amis et les calques du français.
\vspace{4pt}
\begin{multicols}{2}\small
\begin{itemize}
\item À la fin de cette leçon, je sais nommer les étapes d'un entretien d'embauche en anglais (shortlisting, interview, follow-up).
\item À la fin de cette leçon, je sais utiliser le vocabulaire des candidats, des recruteurs et des documents (applicant, interviewer, CV, cover letter).
\item À la fin de cette leçon, je sais employer les collocations essentielles (attend an interview, answer questions confidently, make a good impression).
\item À la fin de cette leçon, je sais utiliser les expressions idiomatiques propres à l'entretien (sell yourself, think on your feet, break the ice).
\item À la fin de cette leçon, je sais éviter les faux amis et les calques du français (to pass an interview, formation, candidature).
\item À la fin de cette leçon, je sais former les mots de la même famille (apply, applicant, application ; employ, employer, employee, employment).
\end{itemize}\end{multicols}}

\begin{sommaire}
\begin{multicols}{2}
\begin{enumerate}
\item The stages of a job interview
\item People and documents
\item Collocations and key expressions
\item False friends and French traps
\item Word families, pronunciation and spelling
\item Vocabulary in context: a model dialogue
\item Activité d'intégration
\item Méthodes et exercices résolus
\item Exercices
\item Corrigés des exercices
\item Fiche bilan
\end{enumerate}
\end{multicols}
\end{sommaire}

\begin{objectifs}
\begin{itemize}
\item À la fin de cette leçon, je sais nommer les étapes d'un entretien d'embauche en anglais (shortlisting, interview, follow-up).
\item À la fin de cette leçon, je sais utiliser le vocabulaire des candidats, des recruteurs et des documents (applicant, interviewer, CV, cover letter).
\item À la fin de cette leçon, je sais employer les collocations essentielles (attend an interview, answer questions confidently, make a good impression).
\item À la fin de cette leçon, je sais utiliser les expressions idiomatiques propres à l'entretien (sell yourself, think on your feet, break the ice).
\item À la fin de cette leçon, je sais éviter les faux amis et les calques du français (to pass an interview, formation, candidature).
\item À la fin de cette leçon, je sais former les mots de la même famille (apply, applicant, application ; employ, employer, employee, employment).
\item À la fin de cette leçon, je sais prononcer et orthographier correctement les mots clés (interview, candidate, appointment, nervous).
\item À la fin de cette leçon, je sais réemployer ce lexique dans des phrases, des dialogues et un court texte sur un entretien d'embauche.
\end{itemize}
\end{objectifs}

\begin{prerequis}
\begin{itemize}
\item Vocabulaire de base du monde du travail vu en Première (job, work, salary, employer).
\item Rédaction d'un CV et d'une lettre de motivation (cover letter) — leçon précédente.
\item Les questions en anglais (mots interrogatifs, inversion du sujet).
\item Le présent simple et le present perfect pour parler de son expérience.
\item Les registres de langue : langage formel vs informel.
\end{itemize}
\end{prerequis}

\newpage

\coursec{The stages of a job interview}

\textbf{Situation d'accroche}

\emph{Imagine: you have just graduated from the University of Yaoundé I and you see an advertisement — the Ministry of Public Service and Administrative Reform (MINFOPRA) is recruiting young bilingual administrators. You send your CV and cover letter, and two weeks later, you receive a phone call: “You are shortlisted for an interview next Tuesday at the Regional Delegation in Buea.” Your heart beats faster. What will the panel ask? How should you dress? What should you say about your strengths and weaknesses?}

En anglais comme en français, réussir un entretien d'embauche (\eng{a job interview}) ne dépend pas seulement de vos diplômes : cela dépend aussi de votre maîtrise des \cle{mots et expressions} propres à chaque étape. Au Cameroun, comme au Nigeria, au Ghana, au Royaume-Uni ou aux États-Unis, le candidat qui connaît le vocabulaire approprié fait immédiatement meilleure impression.

\textbf{Problématique :} quelles sont les étapes d'un entretien d'embauche, et quel vocabulaire anglais permet de décrire chacune d'elles avec précision, sans calquer le français ?

\courssub{Overview: the six stages}

Avant d'entrer dans le détail, observons le parcours complet d'un candidat, de la présélection à la réponse finale.

\begin{popfigure}
\begin{tikzpicture}[node distance=0.35cm]
\tikzset{etape/.style={rectangle, rounded corners=3pt, draw=popblue, fill=popblueL, text width=2.4cm, align=center, minimum height=1.4cm, font=\small}}
\node[etape, align=center] (s1){\textbf{1. Shortlisted}\\ \footnotesize convocation};
\node[etape, right=of s1, align=center] (s2){\textbf{2. Preparation}\\ \footnotesize research, outfit};
\node[etape, right=of s2, align=center] (s3){\textbf{3. Greeting}\\ \footnotesize arrival, handshake};
\node[etape, below=0.6cm of s3, align=center] (s4){\textbf{4. Answering}\\ \footnotesize qualifications, strengths};
\node[etape, left=of s4, align=center] (s5){\textbf{5. Asking}\\ \footnotesize salary, duties};
\node[etape, left=of s5, align=center] (s6){\textbf{6. Follow-up}\\ \footnotesize outcome, letter};
\draw[-{Stealth}, thick, poporange] (s1) -- (s2);
\draw[-{Stealth}, thick, poporange] (s2) -- (s3);
\draw[-{Stealth}, thick, poporange] (s3) -- (s4);
\draw[-{Stealth}, thick, poporange] (s4) -- (s5);
\draw[-{Stealth}, thick, poporange] (s5) -- (s6);
\end{tikzpicture}
\legende{The six stages of a job interview — les six étapes d'un entretien d'embauche.}
\end{popfigure}

\begin{center}
\begin{tabularx}{\linewidth}{|X|X|X|}
\hline
\rowcolor{popblueL} \textbf{English} & \textbf{Français} & \textbf{Remarque} \\
\hline
to be shortlisted & être présélectionné & verbe passif : \eng{be} + participe passé \\
\hline
an interview panel & un jury d'entretien & ne pas dire \eng{a jury} seul \\
\hline
the outcome & le résultat, l'issue & synonyme : \eng{the result} \\
\hline
an appointment letter & une lettre de nomination & contraire : \eng{a rejection letter} \\
\hline
a candidate / an applicant & un candidat & \eng{applicant} = celui qui postule \\
\hline
a post / a position & un poste & faux ami : \eng{a post} n'est pas « la poste » \\
\hline
a CV (curriculum vitae) & un CV & UK : \eng{CV} ; US : \eng{résumé} \\
\hline
a cover letter & une lettre de motivation & accompagne le CV \\
\hline
\end{tabularx}
\end{center}

\courssub{Stage 1: Being shortlisted}

Observons d'abord un document authentique : la convocation que reçoit un candidat présélectionné.

\begin{textebox}[A letter of invitation — une convocation type]
\textbf{CAMWATER Regional Office, Bamenda}

Dear Mr Ndi,

You are shortlisted for the post of accounts clerk. Your interview will take place on Tuesday 14th May at 9 a.m. at our Bamenda office, Commercial Avenue. Please bring your original certificates, your national identity card and a passport photograph. The interview panel will consist of three members. If you cannot attend, kindly inform us at least two days before the date.

We wish you the best of luck.

Yours sincerely,\\
Mrs Comfort Ashu\\
Human Resources Officer
\end{textebox}

\textbf{Glossaire :} \emph{accounts clerk} (nom) : employé de comptabilité ; \emph{to take place} (verbe) : avoir lieu ; \emph{original certificates} (nom) : diplômes originaux ; \emph{to consist of} (verbe) : se composer de ; \emph{kindly} (adverbe) : avoir l'obligeance de / veuillez.

\begin{definition}[To shortlist]
\cle{To shortlist} signifie « sélectionner un petit nombre de candidats pour l'étape finale d'un recrutement ». Le nom correspondant est \cle{the shortlist} : « la liste des présélectionnés ». On emploie presque toujours ce verbe au \textbf{passif} : \eng{She was shortlisted.} (Elle a été présélectionnée.)
\end{definition}

\begin{propriete}[Forme et emploi de to shortlist]
\begin{itemize}
\item \textbf{Forme active :} \eng{The company shortlisted five candidates.} (L'entreprise a présélectionné cinq candidats.)
\item \textbf{Forme passive (la plus fréquente) :} \eng{be} + \eng{shortlisted} : \eng{I was shortlisted for the post.} / \eng{We have been shortlisted.}
\item \textbf{Préposition :} on est présélectionné \eng{for a post} (pour un poste), jamais \eng{to a post}.
\item \textbf{Nom :} \eng{Twenty names are on the shortlist.} (Vingt noms figurent sur la liste des présélectionnés.)
\end{itemize}
\end{propriete}

\begin{exemplebox}[Exemples traduits]
\eng{She was shortlisted for the post.} = Elle a été présélectionnée pour le poste.

\eng{The panel interviewed ten candidates.} = Le jury a interrogé dix candidats.

\eng{Out of two hundred applicants, only six were shortlisted.} = Sur deux cents candidats, six seulement ont été présélectionnés.

\eng{He received a phone call inviting him for an interview.} = Il a reçu un appel téléphonique le convoquant à un entretien.
\end{exemplebox}

\courssub{Stage 2: Preparing for the interview}

Un bon candidat ne s'improvise pas. Trois verbes dominent cette étape :

\begin{itemize}
\item \cle{to research the company} : se renseigner sur l'entreprise (son histoire, ses produits, ses valeurs). \eng{Before the interview, she researched the company on the Internet.}
\item \cle{to rehearse answers} : répéter ses réponses (à voix haute, devant un miroir ou avec un ami). \eng{He rehearsed his answers the night before.}
\item \cle{to choose an outfit} : choisir une tenue (correcte, sobre). \eng{She chose a smart outfit: a dark suit and polished shoes.}
\end{itemize}

\begin{definition}[Vocabulaire de la préparation]
\cle{an outfit} : une tenue vestimentaire ; \cle{smart} : élégant, correct, soigné (attention : \eng{smart} signifie aussi « intelligent » en général, mais pour les vêtements, cela veut dire « bien habillé ») ; \cle{to rehearse} : répéter (prononciation : « ri-heurce ») ; \cle{a mock interview} : un entretien blanc, une simulation d'entretien ; \cle{a cover letter} : une lettre de motivation.
\end{definition}

\begin{exemplebox}[Exemples traduits]
\eng{She practised answering questions with her elder brother.} = Elle s'est entraînée à répondre aux questions avec son frère aîné.

\eng{Candidates should arrive well prepared and well dressed.} = Les candidats doivent arriver bien préparés et bien habillés.

\eng{He read the job advertisement again to understand the duties.} = Il a relu l'annonce d'emploi pour comprendre les fonctions du poste.
\end{exemplebox}

\courssub{Stage 3: Arriving and greeting}

Le jour J, la première impression est décisive. Le vocabulaire clé :

\begin{itemize}
\item \cle{to arrive on time} : arriver à l'heure (jamais en retard : \eng{late}). \eng{Arrive on time, or even ten minutes early.}
\item \cle{to greet the panel} : saluer le jury. \eng{“Good morning, members of the panel.”}
\item \cle{to shake hands} : serrer la main. \eng{He shook hands with each member of the panel.}
\item \cle{to sit down when invited} : s'asseoir quand on vous y invite. \eng{“Please, take a seat.” — “Thank you.”}
\end{itemize}

\begin{attention}[To shake hands : la bonne construction]
On dit \eng{to shake hands with somebody} (serrer la main de quelqu'un), avec \eng{hands} au \textbf{pluriel} et la préposition \eng{with}. Ne dites jamais \eng{to shake the hand of somebody} ni \eng{to shake somebody's hand} dans un contexte formel sans la préposition correcte. Exemple correct : \eng{He shook hands with the chairman.} (Il a serré la main du président.)
\end{attention}

\courssub{Stage 4: Answering questions}

C'est le cœur de l'entretien. Le jury interroge le candidat sur trois grands thèmes :

\begin{itemize}
\item \cle{qualifications} : les diplômes. \eng{“What are your qualifications?” — “I hold a Bachelor's degree in Accounting.”}
\item \cle{experience} : l'expérience professionnelle. \eng{“Do you have any work experience?” — “I did a three-month internship at a bank in Douala.”}
\item \cle{strengths and weaknesses} : les points forts et les points faibles. \eng{“My main strength is that I am hardworking and punctual. My weakness? I sometimes work too much!”}
\end{itemize}

\begin{definition}[Mots clés de cette étape]
\cle{a strength} : un point fort (contraire : \cle{a weakness}, un point faible) ; \cle{an internship} : un stage (américain ; britannique : \eng{a work placement}) ; \cle{to hold a degree} : être titulaire d'un diplôme ; \cle{hardworking} : travailleur ; \cle{punctual} : ponctuel ; \cle{reliable} : fiable.
\end{definition}

\begin{exemplebox}[Exemples traduits]
\eng{The panel asked her about her strengths and weaknesses.} = Le jury l'a interrogée sur ses points forts et ses points faibles.

\eng{He answered all the questions calmly and confidently.} = Il a répondu à toutes les questions calmement et avec assurance.

\eng{She has five years of experience in teaching.} = Elle a cinq ans d'expérience dans l'enseignement.
\end{exemplebox}

\courssub{Stage 5: Asking your own questions}

À la fin, le jury dit souvent : \eng{“Do you have any questions for us?”} Un bon candidat pose une ou deux questions pertinentes sur :

\begin{itemize}
\item \cle{the salary} : le salaire (mensuel, pour employé de bureau). \eng{“May I ask what the starting salary is?”}
\item \cle{the duties} : les fonctions, les tâches. \eng{“What would my main duties be?”}
\item \cle{the working conditions} : les conditions de travail (horaires, congés). \eng{“What are the working hours?”}
\item \cle{promotion prospects} : les perspectives d'avancement. \eng{“Are there opportunities for promotion?”}
\end{itemize}

\begin{attention}[Salary, wages : ne confondez pas]
\eng{Salary} (salaire mensuel fixe, versé aux employés de bureau/cadres) se prononce « sa-la-ri » ; \eng{wages} (salaire hebdomadaire ou journalier, souvent pour le travail manuel/horaire) se prononce « oué-djiz ». Un candidat à un poste de bureau parle de son \eng{salary}, pas de ses \eng{wages}.
\end{attention}

\courssub{Stage 6: The follow-up}

Après l'entretien, le candidat attend \cle{the outcome} (le résultat, l'issue). Deux issues possibles :

\begin{itemize}
\item \cle{an appointment letter} : une lettre de nomination (succès). \eng{Two weeks later, she received her appointment letter.}
\item \cle{a rejection letter} : une lettre de refus (échec). \eng{Unfortunately, he received a rejection letter.}
\end{itemize}

\begin{exemplebox}[Exemples traduits]
\eng{After the interview, the candidates waited anxiously for the outcome.} = Après l'entretien, les candidats ont attendu le résultat avec anxiété.

\eng{She was appointed as a bilingual secretary.} = Elle a été nommée secrétaire bilingue.

\eng{If you are rejected, do not give up: apply again!} = Si vous êtes refusé, n'abandonnez pas : postulez à nouveau !
\end{exemplebox}

\courssub{People and documents}

Voici les acteurs et documents essentiels, avec leurs équivalents précis.

\begin{center}
\begin{tabularx}{\linewidth}{|X|X|X|}
\hline
\rowcolor{popblueL} \textbf{English} & \textbf{Français} & \textbf{Note} \\
\hline
the interviewer / the panel & l'intervieweur / le jury & \eng{panel} = plusieurs personnes \\
\hline
the interviewee / the candidate & l'interviewé / le candidat & \eng{interviewee} = celui qui subit l'entretien \\
\hline
the HR manager / HR officer & le responsable RH / l'agent RH & \eng{HR} = Human Resources \\
\hline
the chairman / the chairperson & le président (du jury) & \eng{chairperson} neutre \\
\hline
a CV / a résumé & un CV & UK : \eng{CV} ; US : \eng{résumé} \\
\hline
a cover letter & une lettre de motivation & \eng{covering letter} aussi \\
\hline
a reference / a referee & une recommandation / un référent & \eng{referee} = personne qui recommande \\
\hline
a testimonial & un certificat de travail / attestation & souvent écrit par un ancien employeur \\
\hline
original certificates & diplômes originaux & \eng{certificates} au pluriel \\
\hline
\end{tabularx}
\end{center}

\courssub{Collocations and key expressions}

Les \cle{collocations} (associations de mots naturelles) sont la clé de la fluidité. Apprenez-les par blocs.

\begin{center}
\begin{tabularx}{\linewidth}{|X|X|X|}
\hline
\rowcolor{popblueL} \textbf{Collocation (English)} & \textbf{Traduction} & \textbf{Exemple} \\
\hline
to apply for a post & postuler à un poste & \eng{He applied for the post of clerk.} \\
\hline
to fill in an application form & remplir un formulaire & \eng{Please fill in the form in block capitals.} \\
\hline
to be shortlisted for an interview & être convoqué à un entretien & \eng{She was shortlisted yesterday.} \\
\hline
to attend an interview & se présenter à un entretien & \eng{Twenty candidates attended.} \\
\hline
to conduct an interview & mener / passer un entretien (recruteur) & \eng{The HR manager conducted the interview.} \\
\hline
to make a good impression & faire bonne impression & \eng{He made a good impression.} \\
\hline
to shake hands with & serrer la main à & \eng{She shook hands with the panel.} \\
\hline
to answer confidently & répondre avec assurance & \eng{Answer confidently and clearly.} \\
\hline
to ask about salary & demander le salaire & \eng{He asked about the salary.} \\
\hline
to receive an appointment letter & recevoir une lettre de nomination & \eng{She received her appointment letter.} \\
\hline
to turn down an offer & refuser une offre & \eng{He turned down the offer.} \\
\hline
\end{tabularx}
\end{center}

\courssub{False friends and French traps}

\begin{attention}[Pièges majeurs pour francophones]
\begin{itemize}
\item \textbf{Assister à} $\neq$ \eng{to assist}. \eng{To assist} = aider. \rightarrow \eng{to attend an interview} / \eng{to go for an interview}.
\item \textbf{Un stage} $\neq$ \eng{a stage}. \eng{A stage} = une scène. \rightarrow \eng{an internship} (US) / \eng{a work placement} (UK).
\item \textbf{Un poste} $\neq$ \eng{a post office}. \eng{A post} = un poste (emploi) ; \eng{the post office} = la poste.
\item \textbf{Une journée} $\neq$ \eng{a journey}. \eng{A journey} = un voyage. \rightarrow \eng{a day} / \eng{a working day}.
\item \textbf{Un salaire} $\neq$ \eng{a salary} uniquement. \eng{Wages} = salaire horaire/hebdomadaire (ouvrier). Distinction de registre.
\item \textbf{Une formation} $\neq$ \eng{a formation} (rare, géologie/militaire). \rightarrow \eng{training} / \eng{a course} / \eng{a degree}.
\item \textbf{Actuellement} $\neq$ \eng{actually}. \eng{Actually} = en fait / en réalité. \rightarrow \eng{currently} / \eng{at present}.
\item \textbf{Prétendre} $\neq$ \eng{to pretend}. \eng{To pretend} = faire semblant. \rightarrow \eng{to claim} / \eng{to apply for}.
\item \textbf{Une opportunité} $\neq$ \eng{an opportunism}. \rightarrow \eng{an opportunity}. \eng{Opportunism} = opportunisme (péjoratif).
\end{itemize}
\end{attention}

\courssub{Word families, pronunciation and spelling}

Travailler les familles de mots permet d'enrichir son lexique rapidement.

\begin{center}
\begin{tabular}{|l|l|l|l|}
\hline
\rowcolor{popblueL} \textbf{Verbe} & \textbf{Nom (action/état)} & \textbf{Nom (agent)} & \textbf{Adjectif} \\
\hline
to apply & an application & an applicant & — \\
\hline
to appoint & an appointment & — & — \\
\hline
to employ & employment / unemployment & an employer / an employee & employed / unemployed \\
\hline
to interview & an interview & an interviewer / an interviewee & — \\
\hline
to qualify & a qualification & — & qualified / unqualified \\
\hline
to recruit & recruitment & a recruiter / a recruit & — \\
\hline
to reject & rejection & — & rejected \\
\hline
to succeed & success & — & successful \\
\hline
to promote & promotion & — & promotional \\
\hline
to train & training & a trainer / a trainee & trained / untrained \\
\hline
\end{tabular}
\end{center}

\begin{definition}[Points de prononciation et d'orthographe]
\begin{itemize}
\item \textbf{Interview} : « in-tè-vyou » (accent sur \textbf{in} et \textbf{view}), pas « in-tèr-vyou ».
\item \textbf{CV} : se prononce lettre par lettre « si-vi » ; pluriel \eng{CVs}.
\item \textbf{Salary} : « sa-la-ri » (3 syllabes, \emph{a} comme dans \eng{cat}).
\item \textbf{Qualification} : « qua-li-fi-kei-chn » (accent sur \textbf{kei}).
\item \textbf{Appointment} : « a-point-ment » (le \emph{oi} se prononce \emph{oï}).
\item \textbf{Strength} : « stren(g)th » (le \emph{g} est souvent muet, \emph{th} = \eng{θ}).
\item \textbf{Weakness} : « wiik-niss » (\emph{ea} = \eng{i:}, \emph{ss} = \eng{s}).
\item \textbf{Rehearse} : « ri-heurce » (\eng{ea} = \eng{eə}, \eng{rse} = \eng{rs}).
\item \textbf{Certificate} : « ser-ti-fi-kit » (accent sur \textbf{ti}).
\end{itemize}
\end{definition}

\courssub{Vocabulary in context: a model dialogue}

Voici un dialogue complet simulant les étapes 3 à 5. Observez le registre formel (\eng{formal register}) et les collocations vues plus haut.

\begin{textebox}[Model dialogue: the interview]
\textbf{Setting:} \emph{Regional Delegation of Secondary Education, Buea. Panel of three.}

\textbf{Chairperson:} Good morning. Please, take a seat.

\textbf{Candidate:} Good morning, Sir. Good morning, Madam. Good morning, Sir. \textit{(Shakes hands with each member)}

\textbf{Chairperson:} You are Mr Atem, correct? You applied for the post of English teacher.

\textbf{Candidate:} Yes, Sir. I hold a Bachelor's degree in English Modern Letters from the University of Buea.

\textbf{Member 1:} What are your main strengths for this position?

\textbf{Candidate:} I am hardworking, punctual and adaptable. During my teaching practice at Government High School Muea, I managed large classes successfully.

\textbf{Member 2:} And your weaknesses?

\textbf{Candidate:} I sometimes spend too much time preparing lessons, but I am learning to balance thoroughness and efficiency.

\textbf{Chairperson:} Do you have any questions for us?

\textbf{Candidate:} Yes, Sir. May I ask what the starting salary is? And what are the opportunities for professional development?

\textbf{Chairperson:} The salary follows the civil service grid. In-service training is organised regularly. We will contact you within two weeks.

\textbf{Candidate:} Thank you very much for your time. Goodbye.

\textbf{Panel:} Goodbye.
\end{textebox}

\begin{exemplebox}[Expressions utiles du dialogue]
\eng{“Please, take a seat.”} = « Asseyez-vous, je vous en prie. » (politesse standard)
\eng{“You applied for the post of...”} = « Vous avez postulé au poste de... »
\eng{“I hold a Bachelor's degree in...”} = « Je suis titulaire d'une licence en... » (plus formel que \eng{I have})
\eng{“During my teaching practice...”} = « Pendant mon stage pédagogique... »
\eng{“May I ask what the starting salary is?”} = « Puis-je demander quel est le salaire de débutant ? » (formule polie)
\eng{“We will contact you within two weeks.”} = « Nous vous contacterons dans les deux semaines. »
\end{exemplebox}

\courssub{Method: telling an interview story in the past}

\begin{methode}[Raconter un entretien au passé]
Pour raconter un entretien (à l'oral ou à l'écrit), utilisez les six étapes de la frise comme plan, et reliez-les avec des marqueurs chronologiques :
\begin{enumerate}
\item \eng{First, I was shortlisted for the post of...} (présélection)
\item \eng{Then, I prepared carefully: I researched the company and rehearsed my answers.} (préparation)
\item \eng{On the day of the interview, I arrived on time and greeted the panel.} (arrivée)
\item \eng{After that, I answered questions about my qualifications and experience.} (réponses)
\item \eng{Before leaving, I asked about the salary and the duties.} (questions)
\item \eng{Finally, two weeks later, I received my appointment letter!} (résultat)
\end{enumerate}
Employez le \textbf{prétérit simple} (\eng{past simple}) pour les actions datées et les marqueurs \eng{first}, \eng{then}, \eng{after that}, \eng{finally}, \eng{two weeks later}.
\end{methode}

\begin{attention}[Le piège du calque : “assister à”]
En français, on dit « assister à un entretien ». En anglais, \eng{to assist} signifie « aider » ! On dit :
\begin{itemize}
\item \eng{to attend an interview} : assister à / se présenter à un entretien ;
\item \eng{to go for an interview} : passer un entretien.
\end{itemize}
Jamais \eng{to assist at an interview}. Exemple : \eng{Twenty candidates attended the interview.} (Vingt candidats se sont présentés à l'entretien.)
\end{attention}

\courssub{Solved exercise}

\begin{exoresolu}[Reorder the stages and translate]
\textbf{Énoncé.} Put the following stages of a job interview in the correct chronological order, then translate each sentence into French.

a) She received her appointment letter.\\
b) She was shortlisted for the post.\\
c) She greeted the panel and shook hands.\\
d) She researched the company and chose an outfit.\\
e) She answered questions about her strengths and weaknesses.\\
f) She asked about the salary and the working conditions.

\tcblower
\textbf{Solution détaillée.} L'ordre chronologique suit la frise : présélection → préparation → salutations → réponses → questions du candidat → résultat.

\begin{enumerate}
\item \textbf{b)} \eng{She was shortlisted for the post.} = Elle a été présélectionnée pour le poste.
\item \textbf{d)} \eng{She researched the company and chose an outfit.} = Elle s'est renseignée sur l'entreprise et a choisi une tenue.
\item \textbf{c)} \eng{She greeted the panel and shook hands.} = Elle a salué le jury et serré les mains.
\item \textbf{e)} \eng{She answered questions about her strengths and weaknesses.} = Elle a répondu aux questions sur ses points forts et ses points faibles.
\item \textbf{f)} \eng{She asked about the salary and the working conditions.} = Elle a posé des questions sur le salaire et les conditions de travail.
\item \textbf{a)} \eng{She received her appointment letter.} = Elle a reçu sa lettre de nomination.
\end{enumerate}
Ordre correct : \textbf{b – d – c – e – f – a}.
\end{exoresolu}

\courssub{Practice exercises}

\begin{exercice}[Vocabulary consolidation]
\textbf{A. Complete the sentences with the correct word from the box:}
\eng{shortlisted — panel — outcome — outfit — appointment letter — internship — strengths — wages — rehearsed — attended}

1. Only eight candidates were \ldots\ for the final interview.
2. The \ldots\ asked him difficult questions about his experience.
3. She bought a new \ldots\ for the big day.
4. Everybody is waiting for the \ldots\ of the interviews.
5. He was so happy to receive his \ldots\ from the Ministry.
6. Before the interview, he \ldots\ his answers in front of a mirror.
7. During his studies, he did a three-month \ldots\ at a bank in Douala.
8. The candidate spoke confidently about his \ldots\ and weaknesses.
9. Manual workers are often paid weekly \ldots\, not a monthly salary.
10. Over fifty people \ldots\ the interview, but only five were selected.

\textbf{B. Match the English expressions (1--5) with their French equivalents (a--e).}
1. to fill in an application form \hfill a. faire bonne impression
2. to make a good impression \hfill b. refuser une offre
3. to turn down an offer \hfill c. mener un entretien
4. to conduct an interview \hfill d. remplir un formulaire
5. to shake hands with \hfill e. serrer la main à

\textbf{C. Spot the error (false friend or wrong collocation) and correct the sentence.}
1. He assisted at the interview yesterday.
2. She is actually working as a teacher.
3. He pretends to the post of director.
4. I did a stage in a hotel last summer.
5. The journey at the office starts at 7:30 am.
\end{exercice}

\begin{corrige}[Exercices « Practice exercises »]
\textbf{A.} 1. \eng{shortlisted} 2. \eng{panel} 3. \eng{outfit} 4. \eng{outcome} 5. \eng{appointment letter} 6. \eng{rehearsed} 7. \eng{internship} 8. \eng{strengths} 9. \eng{wages} 10. \eng{attended}

\textbf{B.} 1--d, 2--a, 3--b, 4--c, 5--e.

\textbf{C.}
1. Erreur : \eng{assisted} $\rightarrow$ \eng{attended}. \eng{He attended the interview yesterday.}
2. Erreur : \eng{actually} $\rightarrow$ \eng{currently}. \eng{She is currently working as a teacher.}
3. Erreur : \eng{pretends} $\rightarrow$ \eng{applied for} / \eng{claims}. \eng{He applied for the post of director.}
4. Erreur : \eng{a stage} $\rightarrow$ \eng{an internship} / \eng{a work placement}. \eng{I did an internship in a hotel last summer.}
5. Erreur : \eng{journey} $\rightarrow$ \eng{working day} / \eng{day}. \eng{The working day at the office starts at 7:30 am.}
\end{corrige}

\courssub{Speaking activity: Mock interview}

\begin{experience}[Entretien blanc — Mock interview]
\textbf{Objectif :} Réemployer le vocabulaire de la leçon en situation orale.
\textbf{Organisation :} Par groupes de 3 (1 candidat, 2 membres du jury). Rotation des rôles.
\textbf{Consignes :}
\begin{enumerate}
\item Le candidat tire un « poste » au sort (ex. : \eng{accounts clerk at CAMWATER}, \eng{English teacher at GHS}, \eng{receptionist at a hotel}, \eng{IT technician at MTN}).
\item Le jury dispose de 2 min pour préparer 3 questions (qualifications, experience, strengths/weaknesses).
\item Le candidat a 1 min pour « préparer » (notes rapides).
\item Déroulement : Greeting $\rightarrow$ Questions/Réponses $\rightarrow$ Candidate's questions $\rightarrow$ Closing.
\item Le jury note le vocabulaire utilisé (\eng{strengths, duties, salary, appointment, shortlisted, rehearsed, outfit...}).
\item Débriefing en classe : quel candidat a utilisé le plus de collocations cibles ?
\end{enumerate}
\end{experience}

\begin{savaistu}[Recrutement au Cameroun et dans le Commonwealth]
Au Cameroun, la fonction publique recrute principalement par \textbf{concours} (\eng{competitive examination}) organisés par le MINFOPRA. Les lauréats reçoivent une \eng{decree of appointment} (\eng{décret de nomination}). Dans le secteur privé (banques, ONG, entreprises comme MTN, Orange, Guinness, CDC), le processus ressemble au modèle anglo-saxon : \eng{CV + cover letter $\rightarrow$ shortlisting $\rightarrow$ interview $\rightarrow$ appointment letter}.
Au Royaume-Uni et au Nigeria, le \eng{CV} (2 pages max) est roi ; aux États-Unis, on utilise un \eng{résumé} (1 page). Au Cameroun anglophone (régions Nord-Ouest et Sud-Ouest), les termes \eng{shortlisted}, \eng{panel}, \eng{appointment letter} sont le standard administratif.
\end{savaistu}

\begin{aretenir}[L'essentiel de la leçon]
\begin{itemize}
\item Un entretien suit six étapes : \eng{shortlisted $\rightarrow$ preparation $\rightarrow$ greeting $\rightarrow$ answering $\rightarrow$ asking $\rightarrow$ follow-up}.
\item \eng{To be shortlisted for a post} = être présélectionné pour un poste ; \eng{the panel} = le jury ; \eng{the outcome} = le résultat.
\item Succès : \eng{an appointment letter} ; échec : \eng{a rejection letter}.
\item Collocations clés : \eng{apply for a post, attend an interview, shake hands with, make a good impression, ask about salary}.
\item Pièges majeurs : \eng{attend} (pas \eng{assist}), \eng{internship/work placement} (pas \eng{stage}), \eng{salary} (pas \eng{wages} pour un bureau), \eng{currently} (pas \eng{actually}).
\item Familles de mots : \eng{apply/application/applicant}, \eng{employ/employer/employee/employment}, \eng{qualify/qualification/qualified}.
\item Pour raconter : \textbf{past simple} + marqueurs \eng{first, then, after that, finally}.
\end{itemize}
\end{aretenir}

\coursec{People and documents}

Dans la section précédente, nous avons suivi les \cle{stages of a job interview} (les étapes d'un entretien d'embauche), de l'annonce du poste jusqu'à la décision finale. Mais qui sont les \emph{personnes} qui participent à cet entretien, et quels \emph{documents} le candidat doit-il préparer ? C'est précisément l'objet de cette deuxième section : connaître les acteurs et les pièces du dossier, c'est déjà maîtriser la moitié du vocabulaire de l'entretien d'embauche.

\courssub{Observation : a short text first}

Lisons d'abord ce court texte original. Soulignez mentalement tous les mots qui désignent des \emph{personnes} et des \emph{documents}.

\begin{textebox}[Recruitment at a Douala shipping company]
Last month, the Atlantic Shipping Company in Douala advertised a \cle{vacancy} for an accounts clerk. More than eighty \cle{applicants} sent their files to the \cle{employer}. Each \cle{candidate} had to \cle{submit a CV} and a \cle{cover letter}, and to \cle{attach} photocopies of their \cle{certificates}. Those who applied online were asked to \cle{fill in an application form} and to \cle{provide references} — that is, the names and telephone numbers of two \cle{referees} who could recommend them.

After studying the files, the human resources officer shortlisted six \cle{job seekers}. On the day of the interview, each \cle{interviewee} was welcomed at the reception desk and asked to show an \cle{identity card}. The \cle{interview} was conducted by a \cle{panel} of three \cle{interviewers}: the managing director, the head of accounts and a representative of a \cle{recruitment agency} based in Bonanjo.

At the end of the day, the panel selected a young woman from Buea. Her \cle{résumé} — she had studied in the United States, so she used the American word — was clear and well organised, and her referees gave excellent \cle{testimonials} about her honesty and hard work. She was offered the \cle{post} and signed her contract a week later.
\end{textebox}

\textbf{Glossaire}\par

\begin{center}
\begin{tabularx}{\linewidth}{|l|l|X|}
\hline
\rowcolor{popblueL} \textbf{Mot anglais} & \textbf{Nature} & \textbf{Traduction française} \\
\hline
vacancy & nom & poste vacant \\
\hline
applicant & nom & candidat (celui qui postule) \\
\hline
to submit & verbe & soumettre, remettre \\
\hline
cover letter & nom & lettre de motivation \\
\hline
to attach & verbe & joindre (à un dossier) \\
\hline
referee & nom & répondant, personne qui recommande \\
\hline
to shortlist & verbe & présélectionner \\
\hline
panel & nom & jury, comité d'entretien \\
\hline
testimonial & nom & attestation, témoignage écrit \\
\hline
\end{tabularx}
\end{center}

\textbf{Comprehension questions}\par

\begin{enumerate}
\item How many people applied for the job?
\item What documents did each candidate have to submit?
\item Who interviewed the candidates?
\item Why did the successful candidate use the word \eng{résumé} instead of \eng{CV}?
\end{enumerate}

\courssub{The people involved in a job interview}

\begin{definition}[The people]
\begin{itemize}
\item \cle{The interviewer} : la personne qui pose les questions pendant l'entretien (le recruteur). \eng{“The interviewer asked me about my experience.”} « Le recruteur m'a interrogé sur mon expérience. »
\item \cle{The panel} : le jury, c'est-à-dire le groupe de plusieurs interviewers qui interrogent ensemble le candidat. \eng{“The panel was made up of three people.”} « Le jury était composé de trois personnes. »
\item \cle{The interviewee} : la personne qui \emph{répond} aux questions, c'est-à-dire le candidat pendant l'entretien. \eng{“The interviewee answered confidently.”} « La personne interrogée a répondu avec assurance. »
\item \cle{The candidate} / \cle{the applicant} : le candidat. \emph{Applicant} insiste sur le fait d'avoir \emph{postulé} (to apply) ; \emph{candidate} est le terme général.
\item \cle{The referee} : le répondant, la personne qui recommande le candidat.
\item \cle{The employer} : l'employeur, la personne ou l'entreprise qui embauche. Attention au mot opposé : \cle{the employee} = l'employé(e), le salarié.
\item \cle{A job seeker} : un chercheur d'emploi, une personne qui cherche du travail.
\end{itemize}
\end{definition}

\begin{propriete}[La logique des suffixes -er / -or et -ee]
L'anglais forme souvent des paires de noms de personnes à partir d'un verbe :
\begin{itemize}
\item le suffixe \textbf{-er / -or} désigne celui qui \emph{fait} l'action : \eng{interview} $\rightarrow$ \eng{interviewer} (celui qui interroge) ; \eng{employ} $\rightarrow$ \eng{employer} (celui qui emploie) ;
\item le suffixe \textbf{-ee} désigne celui qui \emph{subit} l'action : \eng{interview} $\rightarrow$ \eng{interviewee} (celui qui est interrogé) ; \eng{employ} $\rightarrow$ \eng{employee} (celui qui est employé).
\end{itemize}
Autres exemples utiles : \eng{trainer / trainee} (formateur / stagiaire), \eng{examiner / examinee} (examinateur / candidat à un examen).
\end{propriete}

\begin{definition}[Referee]
Un \cle{referee} est une personne qui connaît bien le candidat (un ancien employeur, un professeur, un chef de village) et que l'employeur peut \emph{contacter} pour vérifier ses qualités. On dit aussi \eng{a reference} pour la lettre ou l'attestation écrite par cette personne. \textbf{Attention} : ne pas confondre avec \eng{the referee} qui signifie aussi « l'arbitre » dans un match de football ! Le contexte décide : dans un dossier de candidature, il s'agit toujours du \emph{répondant}.
\end{definition}

\begin{attention}[Prononciation et orthographe de « candidate »]
\begin{itemize}
\item \eng{A candidate} se prononce « \textbf{KAN}-di-date » : le premier \emph{a} est fort (accent tonique), le second \emph{a} est faible (comme un « e » français bref). Ne dites pas « kan-di-dat ».
\item En anglais, le mot prend toujours un \textbf{-e final} : \eng{candidate}. N'écrivez jamais \emph{candidat} comme en français.
\item Même vigilance pour \eng{applicant} (prononcé « a-pli-kant »), qui n'a pas de \emph{-e} final, contrairement au français « applicant » que l'on entend parfois par erreur.
\end{itemize}
\end{attention}

\courssub{The documents}

\begin{definition}[The documents you need]
\begin{itemize}
\item \cle{A CV} (curriculum vitae) : le document qui résume vos études, vos diplômes et votre expérience professionnelle. C'est le terme \textbf{britannique}.
\item \cle{A résumé} : le même document, terme \textbf{américain}. Prononcé « ré-zu-mé ».
\item \cle{A covering letter} (UK) / \cle{a cover letter} (US) / \cle{an application letter} : la lettre de motivation qui accompagne le CV et explique pourquoi vous postulez.
\item \cle{Certificates} : les diplômes et attestations (GCE certificates, university degrees).
\item \cle{Testimonials} : les attestations écrites par d'anciens employeurs ou professeurs.
\item \cle{An identity card} (ID card) : la carte nationale d'identité, souvent exigée le jour de l'entretien.
\item \cle{An application form} : le formulaire de candidature à remplir, sur papier ou en ligne.
\end{itemize}
\end{definition}

\begin{propriete}[UK / USA : deux mots pour une même réalité]
\begin{center}
\begin{tabularx}{\linewidth}{|X|X|X|}
\hline
\rowcolor{popblueL} \textbf{British English (UK)} & \textbf{American English (US)} & \textbf{Français} \\
\hline
a CV & a résumé & un curriculum vitae \\
\hline
a post & a position & un poste \\
\hline
a covering letter & a cover letter & une lettre de motivation \\
\hline
to fill in a form & to fill out a form & remplir un formulaire \\
\hline
\end{tabularx}
\end{center}
Au Cameroun, les deux variétés se rencontrent, mais dans un même document il faut rester \emph{cohérent} : ne mélangez pas \eng{CV} et \eng{résumé} dans la même lettre.
\end{propriete}

\begin{savaistu}[Pas de photo sur un résumé américain]
Aux États-Unis, on ne met presque \emph{jamais} de photo, ni d'âge, ni d'état civil (married, single) sur un \eng{résumé} : la loi lutte contre les discriminations à l'embauche. C'est l'inverse de l'usage courant au Cameroun francophone, où le CV comporte souvent une photo d'identité et la mention « célibataire » ou « marié(e) ». En Grande-Bretagne aussi, la photo est facultative et de plus en plus rare. Si vous postulez dans une entreprise américaine, supprimez la photo !
\end{savaistu}

\courssub{Tableau récapitulatif : people, documents and key words}\par

\begin{center}
\begin{tabularx}{\linewidth}{|X|X|X|}
\hline
\rowcolor{popblueL} \textbf{English} & \textbf{Français} & \textbf{Remarque} \\
\hline
a job seeker & un chercheur d'emploi & aussi : a job hunter \\
\hline
an applicant / a candidate & un candidat & \emph{candidate} avec -e final \\
\hline
an interviewee & la personne interrogée & suffixe -ee = qui subit \\
\hline
an interviewer & le recruteur, l'intervieweur & suffixe -er = qui fait \\
\hline
a panel & le jury & plusieurs interviewers \\
\hline
a referee & le répondant & pas l'arbitre ici ! \\
\hline
an employer / an employee & l'employeur / l'employé & paire -er / -ee \\
\hline
a recruitment agency & un cabinet de recrutement & agency = agence \\
\hline
a vacancy & un poste vacant & adjectif : vacant \\
\hline
a post (UK) / a position (US) & un poste & synonymes \\
\hline
a CV (UK) / a résumé (US) & un CV & rester cohérent \\
\hline
a cover letter & une lettre de motivation & to cover = accompagner \\
\hline
an application form & un formulaire de candidature & to fill in (UK) \\
\hline
\end{tabularx}
\end{center}

\courssub{Collocations with documents}

En anglais, chaque nom « appelle » un verbe précis : ce sont les \cle{collocations}. Les apprendre par paires évite les calques du français.

\begin{propriete}[Verbe + document : les collocations essentielles]
\begin{itemize}
\item \cle{to fill in an application form} : remplir un formulaire (US : \eng{to fill out}).
\item \cle{to submit a CV} : soumettre, remettre un CV (plus formel que \eng{to send}).
\item \cle{to attach certificates} : joindre des diplômes (à un dossier ou à un courriel).
\item \cle{to provide references} : fournir des références (les coordonnées des referees).
\item \cle{to write a cover letter} : rédiger une lettre de motivation.
\item \cle{to show an identity card} : présenter une pièce d'identité.
\end{itemize}
\end{propriete}

\begin{exemplebox}[Exemples traduits]
\begin{itemize}
\item \eng{“Applicants must submit a CV and a cover letter.”} « Les candidats doivent soumettre un CV et une lettre de motivation. »
\item \eng{“Please fill in this form.”} « Veuillez remplir ce formulaire. »
\item \eng{“She attached her certificates to the email.”} « Elle a joint ses diplômes au courriel. »
\item \eng{“You must provide two references.”} « Vous devez fournir deux références. »
\item \eng{“The interviewee showed his identity card at the entrance.”} « Le candidat a présenté sa carte d'identité à l'entrée. »
\end{itemize}
\end{exemplebox}

\begin{attention}[Pièges fréquents des francophones]
\begin{itemize}
\item Ne dites pas \emph{to fill a form} : il faut la préposition \eng{in} (UK) ou \eng{out} (US) : \eng{to fill in a form}.
\item \eng{A certificate} = un diplôme, une attestation ; ne confondez pas avec \eng{a medical certificate} (certificat médical).
\item \eng{A formation} n'existe pas au sens de « formation professionnelle » : dites \eng{training} ou \eng{qualifications}.
\item \eng{Actually} ne veut pas dire « actuellement » mais « en fait » ; pour « actuellement », dites \eng{currently} : \eng{“I am currently looking for a job.”} « Je cherche actuellement un emploi. »
\end{itemize}
\end{attention}

\courssub{Mind map : people and documents}\par

\begin{popfigure}
\begin{center}
\begin{tikzpicture}[
  mindmap,
  grow cyclic,
  every node/.style={concept, align=center, font=\small},
  root concept/.append style={concept color=popblue, font=\bfseries\small},
  level 1/.append style={level distance=3.2cm, sibling angle=180},
  level 2/.append style={level distance=2.6cm, sibling angle=45}
]
\node[root concept, align=center]{JOB\\INTERVIEW}
  child[concept color=popgreen!70] { node {PEOPLE}
    child[concept color=popgreen!40] { node[align=center]{interviewer\\panel} }
    child[concept color=popgreen!40] { node[align=center]{candidate\\interviewee\\applicant} }
    child[concept color=popgreen!40] { node[align=center]{referee\\employer} }
  }
  child[concept color=poporange!70] { node {DOCUMENTS}
    child[concept color=poporange!40] { node {CV / résumé} }
    child[concept color=poporange!40] { node[align=center]{cover letter\\certificates} }
    child[concept color=poporange!40] { node[align=center]{application\\form / ID card} }
  };
\end{tikzpicture}
\end{center}
\legende{Carte mentale du vocabulaire : les personnes et les documents de l'entretien d'embauche.}
\end{popfigure}

\courssub{Exercice résolu et entraînement}

\begin{exoresolu}[Complétez avec le mot qui convient]
Énoncé — Complétez chaque phrase avec l'un des mots suivants : \eng{referee}, \eng{panel}, \eng{vacancy}, \eng{application form}, \eng{résumé}.
\begin{enumerate}
\item The company advertised a \dots\ for a secretary.
\item All the candidates were questioned by a \dots\ of four interviewers.
\item My former headmaster agreed to be my \dots .
\item Please fill in this \dots\ and sign at the bottom.
\item In the United States, you send a \dots , not a CV.
\end{enumerate}
\tcblower
Solution détaillée :
\begin{enumerate}
\item \textbf{vacancy} : on « annonce un poste vacant » (\eng{to advertise a vacancy}).
\item \textbf{panel} : un groupe de plusieurs interviewers forme un jury (\eng{a panel of four}).
\item \textbf{referee} : l'ancien proviseur est un répondant qui peut recommander le candidat.
\item \textbf{application form} : on remplit un formulaire (\eng{to fill in}) et on signe en bas.
\item \textbf{résumé} : c'est le terme américain pour \eng{CV}.
\end{enumerate}
\end{exoresolu}

\begin{exercice}[À vous de jouer]
\begin{enumerate}
\item Traduisez en anglais : « Les candidats doivent joindre leurs diplômes et fournir deux références. »
\item Complétez : \eng{The \dots\ asked the questions and the \dots\ answered them politely.} (interviewer / interviewee)
\item UK ou US ? Classez : \eng{résumé}, \eng{CV}, \eng{post}, \eng{position}, \eng{to fill out}, \eng{to fill in}.
\item Rédigez trois phrases avec : \eng{to submit a CV}, \eng{to attach certificates}, \eng{to provide references}.
\end{enumerate}
\end{exercice}

\begin{corrige}[Exercice « À vous de jouer »]
\begin{enumerate}
\item \eng{“Applicants must attach their certificates and provide two references.”}
\item \eng{“The \textbf{interviewer} asked the questions and the \textbf{interviewee} answered them politely.”} (celui qui fait l'action : -er ; celui qui la subit : -ee).
\item UK : \eng{CV}, \eng{post}, \eng{to fill in} — US : \eng{résumé}, \eng{position}, \eng{to fill out}.
\item Exemples de réponses : \eng{“I submitted my CV online yesterday.”} / \eng{“She attached her GCE certificates to the file.”} / \eng{“The company asked me to provide two references.”}
\end{enumerate}
\end{corrige}

\begin{aretenir}[L'essentiel de la section]
\begin{itemize}
\item Les personnes : \eng{interviewer / panel} (qui interrogent), \eng{candidate / applicant / interviewee} (qui postulent et répondent), \eng{referee} (qui recommande), \eng{employer} (qui embauche).
\item Les documents : \eng{CV} (UK) / \eng{résumé} (US), \eng{cover letter}, \eng{certificates}, \eng{application form}, \eng{identity card}.
\item Les collocations à mémoriser par cœur : \eng{to fill in a form}, \eng{to submit a CV}, \eng{to attach certificates}, \eng{to provide references}.
\item Orthographe et prononciation : \eng{candidate} avec -e final, prononcé « \textbf{KAN}-di-date ».
\end{itemize}
\end{aretenir}

\coursec{Collocations and key expressions}

In the previous section, we met the people of the interview — the \cle{candidate}, the \cle{interviewer}, the \cle{panel} — and the documents they use, such as the \cle{CV} and the \cle{cover letter}. But knowing words in isolation is not enough. In real English, words travel in groups: we do not simply say \eng{interview}, we say \eng{to attend an interview}; we do not simply say \eng{impression}, we say \eng{to make a good impression}. These natural word partnerships are called \cle{collocations}, and they are exactly what makes a candidate sound fluent and professional — in English as in any language. This section teaches you the essential collocations and fixed expressions of the job interview, with the questions you will hear and the answers you can give.

\courssub{What is a collocation?}

\begin{definition}[Collocation]
A \cle{collocation} is a natural, fixed combination of two or more words that native speakers habitually use together: \eng{make an impression}, \eng{attend an interview}, \eng{answer a question}. You cannot guess a collocation by translating word for word from French; you must learn it as a single unit.
\end{definition}

\begin{propriete}[Collocations are fixed partnerships]
Les collocations sont des associations \textbf{fixes} de mots. On dit \eng{to make an impression}, jamais \eng{to do an impression} dans ce sens. On dit \eng{to attend an interview}, jamais \eng{to assist an interview}. Changer le verbe, même pour un synonyme, produit une phrase incorrecte ou étrange. Retenez toujours le couple \textbf{verbe + nom} ensemble.
\end{propriete}

\begin{attention}[Faux amis et calques du français]
\begin{itemize}
\item \eng{to assist} signifie « aider » ; « assister à un entretien » se dit \eng{to attend an interview}.
\item « passer un entretien » ne se traduit pas par \eng{pass an interview} : on dit \eng{to sit for an interview} ou \eng{to go for an interview}.
\item « faire une impression » au sens de « produire un effet » = \eng{to make an impression}. \eng{To do an impression of someone} signifie « imiter quelqu'un » (imitation comique) !
\end{itemize}
\end{attention}

\courssub{Collocations with \eng{interview}}

\begin{center}
\begin{tabularx}{\linewidth}{|X|X|X|}\hline
\rowcolor{popblueL} Collocation & Français & Exemple \\ \hline
to attend an interview & assister à / se présenter à un entretien & \eng{She attended three interviews last week.} \\ \hline
to go for an interview & aller à un entretien & \eng{I am going for an interview in Douala tomorrow.} \\ \hline
to sit for an interview & passer un entretien (formel) & \eng{Twenty candidates sat for the interview.} \\ \hline
a job interview & un entretien d'embauche & \eng{He has a job interview at the bank.} \\ \hline
a panel interview & un entretien devant un jury & \eng{A panel interview can be stressful.} \\ \hline
a phone interview & un entretien téléphonique & \eng{The first stage was a phone interview.} \\ \hline
a video interview & un entretien en visioconférence & \eng{We had a video interview via Zoom.} \\ \hline
\end{tabularx}
\end{center}

\courssub{Collocations of behaviour: before and during the interview}

Observez d'abord ce court texte, puis étudions les expressions en gras.

\begin{textebox}[Advice from a careers adviser in Yaoundé]
When you attend a job interview, small details matter enormously. First, dress smartly: a clean shirt, polished shoes, and tidy hair show respect for the panel. Second, arrive on time — in fact, arrive fifteen minutes early. When you enter the room, greet the interviewers politely and wait to be invited to sit down. During the interview, sit upright; do not cross your arms, because this looks defensive. Keep eye contact with the person who is speaking to you, and speak clearly, at a calm speed. If you do not understand a question, simply ask the interviewer to repeat it. Answer every question confidently, even when the question is difficult, and always give examples from your studies or your experience. Highlight your strengths, but be honest enough to admit one or two weaknesses — and explain how you are working on them. Finally, at the end, thank the panel for the opportunity. Remember: the interviewers are not your enemies. They want you to succeed, because they need a strong candidate. Your job is simply to sell yourself.
\end{textebox}

\textbf{Glossaire.} \emph{careers adviser} (n.) : conseiller d'orientation ; \emph{smartly} (adv.) : élégamment ; \emph{polished} (adj.) : cirées ; \emph{upright} (adv.) : bien droit ; \emph{defensive} (adj.) : sur la défensive ; \emph{strengths} (n.) : points forts ; \emph{weaknesses} (n.) : points faibles.

\textbf{Comprehension questions.}
\begin{enumerate}
\item Why should a candidate arrive early?
\item What does crossing your arms communicate?
\item What should you do if you do not understand a question?
\end{enumerate}

\begin{aretenir}[Les collocations de comportement]
\begin{itemize}
\item \eng{to dress smartly} = s'habiller élégamment
\item \eng{to arrive on time} = arriver à l'heure
\item \eng{to speak clearly} = parler clairement
\item \eng{to keep eye contact} = maintenir le contact visuel
\item \eng{to sit upright} = se tenir bien droit
\item \eng{to make a good impression} = faire bonne impression
\end{itemize}
\end{aretenir}

\begin{attention}[Le contact visuel : une question de culture]
\eng{To keep eye contact} (regarder l'interlocuteur dans les yeux) est fortement valorisé dans les cultures anglophones : il signale l'honnêteté et la confiance en soi. Éviter le regard peut être interprété comme un manque de confiance ou de sincérité. Attention cependant : il ne faut pas fixer (\eng{stare}) non plus — un contact visuel naturel et régulier suffit.
\end{attention}

\courssub{Collocations of answering}

\begin{center}
\begin{tabularx}{\linewidth}{|X|X|}\hline
\rowcolor{popgreenL} Collocation & Français \\ \hline
to answer questions confidently & répondre aux questions avec assurance \\ \hline
to give examples & donner des exemples \\ \hline
to highlight your strengths & mettre en avant vos points forts \\ \hline
to admit your weaknesses & admettre vos points faibles \\ \hline
to ask for clarification & demander une précision \\ \hline
to back up your answers & étayer vos réponses \\ \hline
\end{tabularx}
\end{center}

\begin{exemplebox}[Exemples traduits]
\eng{You must sell yourself during the interview.} « Tu dois te mettre en valeur pendant l'entretien. »

\eng{She answered all the questions confidently.} « Elle a répondu à toutes les questions avec assurance. »

\eng{He highlighted his strengths but also admitted one weakness.} « Il a mis en avant ses points forts mais a aussi admis un point faible. »

\eng{Always back up your answers with concrete examples.} « Étayez toujours vos réponses avec des exemples concrets. »
\end{exemplebox}

\courssub{Idiomatic expressions}

\begin{definition}[Idiom]
An \cle{idiom} is a fixed expression whose meaning cannot be understood from the individual words: \eng{to break the ice} has nothing to do with ice!
\end{definition}

\begin{itemize}
\item \eng{to sell yourself} = se mettre en valeur, se « vendre » (parler de ses qualités sans fausse modestie).
\item \eng{to think on your feet} = réagir vite, trouver une réponse immédiatement.
\item \eng{to break the ice} = briser la glace, détendre l'atmosphère.
\item \eng{to be a strong candidate} = être un candidat sérieux, bien placé pour obtenir le poste.
\item \eng{to land the job} = décrocher le poste (familier).
\end{itemize}

\begin{savaistu}[Small talk before the interview]
Au Royaume-Uni et aux États-Unis, l'entretien commence souvent par du \eng{small talk} — une conversation légère sur la météo, le voyage ou le bureau — pour \eng{break the ice} avant les questions sérieuses. Savoir participer à ce rituel en anglais simple et souriant fait déjà partie de l'épreuve !
\end{savaistu}

\courssub{Word families, pronunciation and spelling}

\begin{propriete}[Mots de la même famille]
De nombreux noms d'entretien forment des familles régulières. Apprenez le nom, le verbe, l'adjectif et l'adverbe ensemble.
\end{propriete}

\begin{center}
\begin{tabularx}{\linewidth}{|X|X|X|X|}\hline
\rowcolor{poppurpleL} Noun & Verb & Adjective & Adverb \\ \hline
interview & to interview & interview (adj.) & — \\ \hline
impression & to impress & impressive & impressively \\ \hline
preparation & to prepare & prepared / preparatory & — \\ \hline
qualification & to qualify & qualified & — \\ \hline
application & to apply & applicable & — \\ \hline
experience & to experience & experienced & — \\ \hline
confidence & — & confident & confidently \\ \hline
strength & to strengthen & strong & strongly \\ \hline
weakness & to weaken & weak & weakly \\ \hline
\end{tabularx}
\end{center}

\begin{aretenir}[Prononciation — pièges fréquents]
\begin{itemize}
\item \eng{interview} : « in-tè-vyou » (accent sur la 1re syllabe, pas « in-tar-vyou »).
\item \eng{CV} : se lit lettre par lettre « si-vi » (pas « curriculum vitae » à l'oral).
\item \eng{qualification} : « qua-li-fi-cay-chn » (son \eng{/ʃ/} pour \eng{ti}).
\item \eng{experience} : « ik-spi-ri-ence » (le \eng{x} se prononce \eng{/k/} + \eng{/s/}).
\item \eng{confident} : « kon-fi-dent » (le \eng{ent} final muet, pas « kon-fi-dant »).
\item \eng{panel} : « pa-nel » (accent sur la 1re syllabe, \eng{a} court).
\end{itemize}
\end{aretenir}

\begin{attention}[Orthographe : doubles consonnes et \eng{-ence}/\eng{-ance}]
\begin{itemize}
\item \eng{interview} (un seul \eng{r}, un seul \eng{v}) — ne pas écrire *\eng{interveiw} ni *\eng{interrview}.
\item \eng{impression} (double \eng{s}, double \eng{s}) — \eng{impressive} (un seul \eng{s}).
\item \eng{preparation} (un seul \eng{p}, un seul \eng{r}) — \eng{prepare} (un seul \eng{p}).
\item \eng{experience} / \eng{qualification} / \eng{application} : terminaison \eng{-ence} ou \eng{-ation}, jamais \eng{-ance} pour ces mots.
\item \eng{confident} (terminaison \eng{-ent}) vs \eng{confidence} (terminaison \eng{-ence}).
\end{itemize}
\end{attention}

\courssub{Questions of the panel and possible answers}

\begin{center}
\begin{tabularx}{\linewidth}{|X|X|}\hline
\rowcolor{poporangeL} Questions of the panel & Possible answers \\ \hline
\eng{Tell us about yourself.} & \eng{I am a hard-working and reliable person.} \\ \hline
\eng{Why should we hire you?} & \eng{I believe I have the qualifications you need.} \\ \hline
\eng{What are your strengths and weaknesses?} & \eng{I am organised, but I sometimes work too slowly because I check everything twice.} \\ \hline
\eng{Where do you see yourself in five years?} & \eng{I hope to grow with your company and take on more responsibility.} \\ \hline
\eng{Why do you want to work for us?} & \eng{I admire your company's reputation and values.} \\ \hline
\eng{Do you have any questions for us?} & \eng{Yes, could you tell me about the team I would be working with?} \\ \hline
\end{tabularx}
\end{center}

\begin{popfigure}
\begin{tikzpicture}[font=\small]
\node[draw=popblue, fill=popblueL, rounded corners, align=center, text width=4.6cm] (q1) at (0,5.5) {\textbf{Tell us about}\\ \textbf{yourself.}};
\node[draw=popblue, fill=popblueL, rounded corners, align=center, text width=4.6cm] (q2) at (0,4) {\textbf{Why should}\\ \textbf{we hire you?}};
\node[draw=popblue, fill=popblueL, rounded corners, align=center, text width=4.6cm] (q3) at (0,2.5) {\textbf{What are your strengths}\\ \textbf{and weaknesses?}};
\node[draw=popblue, fill=popblueL, rounded corners, align=center, text width=4.6cm] (q4) at (0,1) {\textbf{Where do you see yourself}\\ \textbf{in five years?}};
\node[draw=popblue, fill=popblueL, rounded corners, align=center, text width=4.6cm] (q5) at (0,-0.5) {\textbf{Why do you want to}\\ \textbf{work for us?}};
\node[draw=popgreen, fill=popgreenL, rounded corners, align=center, text width=5.4cm] (a1) at (8.6,5.5) {I am a hard-working\\ and reliable person.};
\node[draw=popgreen, fill=popgreenL, rounded corners, align=center, text width=5.4cm] (a2) at (8.6,4) {I believe I have the\\ qualifications you need.};
\node[draw=popgreen, fill=popgreenL, rounded corners, align=center, text width=5.4cm] (a3) at (8.6,2.5) {I am organised; I sometimes\\ check everything twice.};
\node[draw=popgreen, fill=popgreenL, rounded corners, align=center, text width=5.4cm] (a4) at (8.6,1) {I hope to grow with your\\ company.};
\node[draw=popgreen, fill=popgreenL, rounded corners, align=center, text width=5.4cm] (a5) at (8.6,-0.5) {I admire your company's\\ reputation and values.};
\draw[-{Stealth}, thick, poporange] (q1) -- (a1);
\draw[-{Stealth}, thick, poporange] (q2) -- (a2);
\draw[-{Stealth}, thick, poporange] (q3) -- (a3);
\draw[-{Stealth}, thick, poporange] (q4) -- (a4);
\draw[-{Stealth}, thick, poporange] (q5) -- (a5);
\node[text=popdark, align=center] at (0,6.6) {\textbf{QUESTIONS OF THE PANEL}};
\node[text=popdark, align=center] at (8.6,6.6) {\textbf{POSSIBLE ANSWERS}};
\end{tikzpicture}
\legende{Des questions typiques du jury aux réponses modèles du candidat.}
\end{popfigure}

\courssub{How to memorise collocations}

\begin{methode}[Apprendre les couples, pas les mots isolés]
\begin{enumerate}
\item N'apprenez jamais un nom seul : apprenez chaque \textbf{verbe avec son nom} — \eng{attend + interview}, \eng{make + impression}, \eng{answer + question}, \eng{keep + eye contact}.
\item Écrivez chaque collocation dans une phrase complète qui vous concerne : \eng{I will attend an interview in Buea next month.}
\item Regroupez les collocations par thème (avant / pendant / après l'entretien) dans une carte mentale.
\item Testez-vous à voix haute : le français d'un côté, l'anglais de l'autre.
\end{enumerate}
\end{methode}

\begin{exoresolu}[Choose the correct collocation]
Énoncé — Complétez avec le bon verbe : \eng{make}, \eng{attend}, \eng{keep}, \eng{highlight}, \eng{break}.
\begin{enumerate}
\item You should \_\_\_\_\_\_ the interview at least ten minutes early.
\item Try to \_\_\_\_\_\_ a good impression from the first minute.
\item During the interview, \_\_\_\_\_\_ eye contact with the panel.
\item In your answers, \_\_\_\_\_\_ your strengths with concrete examples.
\item A joke can \_\_\_\_\_\_ the ice at the beginning.
\end{enumerate}
\tcblower
Solution détaillée.
\begin{enumerate}
\item \eng{attend} — on « assiste à » un entretien : \eng{attend an interview} (jamais \eng{assist}).
\item \eng{make} — la collocation fixe est \eng{make a good impression} (jamais \eng{do}).
\item \eng{keep} — \eng{keep eye contact} : maintenir le contact visuel.
\item \eng{highlight} — \eng{highlight your strengths} : mettre en avant ses points forts.
\item \eng{break} — l'idiome est \eng{break the ice} : détendre l'atmosphère.
\end{enumerate}
\end{exoresolu}

\begin{exercice}[Your turn]
Traduisez en anglais : 1. « J'ai passé un entretien devant un jury hier. » 2. « Elle s'est habillée élégamment et est arrivée à l'heure. » 3. « Il faut te mettre en valeur et répondre avec assurance. » 4. « Merci pour cette opportunité. »
\end{exercice}

\begin{corrige}[Exercice « Your turn »]
1. \eng{I sat for a panel interview yesterday.} 2. \eng{She dressed smartly and arrived on time.} 3. \eng{You must sell yourself and answer confidently.} 4. \eng{Thank you for this opportunity.}
\end{corrige}

\begin{experience}[Role play: the first two minutes]
Par groupes de trois : un candidat et deux membres du jury. Le jury commence par du \eng{small talk} (\eng{Did you find our office easily?}), puis pose la question \eng{Tell us about yourself.} Le candidat doit utiliser au moins trois collocations de la leçon (\eng{make a good impression}, \eng{highlight my strengths}, \eng{keep eye contact}\ldots). Changez ensuite les rôles.
\end{experience}

You now master the word partnerships of the interview. In the next section, we will hunt down the false friends and French traps that can spoil an otherwise excellent performance.

\coursec{Les faux amis et pièges du français}

Dans la section précédente, nous avons appris les collocations et expressions clés qui rendent votre anglais naturel en contexte d'entretien, comme \eng{to apply for a job} ou \eng{to meet the requirements}. Mais connaître les bons mots ne suffit pas : il faut aussi éviter les \emph{mauvais} mots — ceux qui semblent familiers car ils ressemblent au français, mais ont un sens différent en anglais. Ces mots traîtres s'appellent des \cle{faux amis} (\eng{false friends}), et ils sont l'une des principales raisons pour lesquelles les candidats francophones perdent des points au Baccalauréat et de la crédibilité lors de vrais entretiens. Dans cette section, nous les observerons, les classerons et apprendrons une méthode fiable pour les éviter.

\begin{definition}[Qu'est-ce qu'un faux ami ?]
Un \cle{faux ami} (\eng{false friend}) est un mot anglais qui ressemble à un mot français mais a un \textbf{sens différent}. Par exemple, le mot anglais \eng{formation} existe, mais il ne signifie \emph{pas} « formation professionnelle » : il désigne la forme ou l'arrangement de quelque chose (a rock formation, a military formation). Les faux amis sont dangereux précisément parce qu'ils \emph{paraissent} corrects.
\end{definition}

\courssub{Observation : repérez les erreurs}

Lisez le texte suivant. Une candidate francophone, Marie de Douala, l'a rédigé pour son dossier de candidature. Il contient \textbf{six} pièges typiques du français. Saurez-vous les trouver avant de lire la correction ?

\begin{textebox}[Brouillon de Marie — avec erreurs cachées]
Dear Sir,

I am writing because I passed an interview in your company last month and I am waiting for the results. I have twenty-two years and I have a good formation in accounting. I assisted to several training sessions in Yaoundé, and I have three years of experience in a bank in Douala. I sent my candidature in January. I am agree that the job is demanding, but I am motivated and I can make many efforts.

Yours faithfully,

Marie Ngo Bell
\end{textebox}

\textbf{Glossaire :} \emph{draft} (n.) : brouillon ; \emph{accounting} (n.) : comptabilité ; \emph{demanding} (adj.) : exigeant.

Les six erreurs sont : \eng{passed an interview}, \eng{I have twenty-two years}, \eng{formation}, \eng{assisted to}, \eng{candidature}, \eng{I am agree}. Étudions maintenant chaque piège en détail.

\courssub{Faux ami 1 : « passer un entretien » n'est pas “to pass an interview”}

En français, « passer » signifie à la fois « se présenter à » (un examen, un entretien) et « réussir » (« j'ai passé mon bac » = je l'ai passé ; « je l'ai passé ! » peut aussi suggérer la réussite). L'anglais sépare clairement les deux idées.

\begin{propriete}[« Pass », « take », « succeed in » : la règle]
\begin{itemize}
\item \eng{to take an exam / an interview} ou \eng{to sit an exam} = \emph{se présenter à}, participer à (sans dire le résultat).
\item \eng{to pass an exam / a test} = \emph{réussir} un examen, un test (avec un résultat admis/recadré).
\item Pour un entretien, l'anglais préfère : \eng{to succeed in an interview}, \eng{to do well in an interview}, \eng{to get through the interview} = réussir un entretien.
\end{itemize}
\end{propriete}

\begin{exemplebox}[Exemples — traduits]
\eng{She did well in the interview and got the job.} « Elle a réussi son entretien et a obtenu le poste. »

\eng{He succeeded in his second interview in Bamenda.} « Il a réussi son deuxième entretien à Bamenda. »

\eng{I passed my GCE Advanced Level last year.} « J'ai réussi mon baccalauréat l'année dernière. » (examen → \eng{pass} correct)

\eng{Over fifty candidates took the interview; only five got through.} « Plus de cinquante candidats ont passé l'entretien ; seulement cinq ont été retenus. »
\end{exemplebox}

\begin{attention}[Piège]
✗ \eng{He passed the interview.} est un calque du français. Dites : \eng{He did well in the interview.} ou \eng{He succeeded in the interview.} = « Il a réussi son entretien. »
\end{attention}

\courssub{Faux ami 2 : “formation” ≠ « formation »}

Le mot anglais \eng{formation} existe — c'est pourquoi il est si dangereux. Mais il signifie « formation » seulement au sens géologique ou militaire : \eng{a rock formation} (une formation rocheuse), \eng{to fly in formation} (voler en formation).

\begin{propriete}[Comment dire « formation » en anglais]
\begin{itemize}
\item \eng{training} = formation professionnelle, stage de formation : \eng{vocational training} (formation professionnelle), \eng{a training course} (une formation, un stage).
\item \eng{education} = formation scolaire, études : \eng{higher education} (enseignement supérieur).
\item \eng{qualifications} = diplômes et titres : \eng{academic qualifications}.
\end{itemize}
\end{propriete}

\begin{exemplebox}[Exemples — traduits]
\eng{She received excellent training in nursing.} « Elle a reçu une excellente formation en soins infirmiers. »

\eng{The job requires a good education and computer skills.} « Le poste exige une bonne formation et des compétences en informatique. »

✗ \eng{I have a good formation.} → ✓ \eng{I have a good training background.} / \eng{I am well trained.}
\end{exemplebox}

\courssub{Faux ami 3 : « candidature » n'est pas un mot anglais}

Le mot \eng{candidature} n'\textbf{existe pas} en anglais courant. Le mot correct est \eng{application}.

\begin{propriete}[La famille « application »]
\begin{itemize}
\item \eng{an application} = une candidature ; \eng{an application letter / a cover letter} = une lettre de motivation ; \eng{an application form} = un formulaire de candidature.
\item \eng{to apply for a job} = poser sa candidature à un poste.
\item \eng{an applicant / a candidate} = un candidat.
\end{itemize}
\end{propriete}

\begin{exemplebox}[Exemples — traduits]
\eng{I sent my application two weeks ago.} « J'ai envoyé ma candidature il y a deux semaines. »

\eng{Over two hundred applicants applied for the post.} « Plus de deux cents candidats ont postulé pour le poste. »

✗ \eng{my candidature} → ✓ \eng{my application}.
\end{exemplebox}

\courssub{Faux ami 4 : « assister à » n'est pas “to assist”}

\begin{propriete}[« To assist » vs « to attend »]
\begin{itemize}
\item \eng{to assist (somebody)} = \textbf{aider} quelqu'un : \eng{The secretary assisted the manager.} « La secrétaire a aidé le directeur. »
\item \eng{to attend} = \textbf{assister à}, être présent à : \eng{to attend an interview / a meeting / a training session}. Attention : \eng{attend} est suivi \textbf{directement} du complément, sans préposition : ✗ \eng{attend to a meeting} → ✓ \eng{attend a meeting}.
\end{itemize}
\end{propriete}

\begin{exemplebox}[Exemples — traduits]
\eng{All the shortlisted candidates attended the interview.} « Tous les candidats présélectionnés ont assisté à l'entretien. »

\eng{She attends evening classes in Buea.} « Elle suit des cours du soir à Buea. »

✗ \eng{I assisted to the interview.} → ✓ \eng{I attended the interview.}
\end{exemplebox}

\courssub{Faux ami 5 : “experience” vs “an experiment”}

Bonne nouvelle : \eng{experience} est un vrai ami — \eng{work experience} = « expérience professionnelle ». Mais attention à son jumeau : \eng{an experiment} = une \emph{expérience scientifique} (de laboratoire).

\begin{exemplebox}[Exemples — traduits]
\eng{She has five years of experience as an accountant.} « Elle a cinq ans d'expérience comme comptable. »

\eng{The students carried out an experiment in the science laboratory.} « Les élèves ont réalisé une expérience au laboratoire. »

Notez aussi : \eng{experienced} (adj.) = expérimenté : \eng{an experienced driver} « un chauffeur expérimenté ».
\end{exemplebox}

\courssub{Calques grammaticaux : “I have twenty years” et “I am agree”}

Certains pièges ne sont pas lexicaux mais \emph{structurels}, calqués sur le français.

\begin{propriete}[Deux calques classiques à éliminer]
\begin{enumerate}
\item Âge : le français utilise « avoir », l'anglais utilise \textbf{to be} : ✗ \eng{I have twenty years.} → ✓ \eng{I am twenty years old.} « J'ai vingt ans. »
\item Accord : \eng{agree} est un \textbf{verbe}, pas un adjectif : ✗ \eng{I am agree.} → ✓ \eng{I agree.} « Je suis d'accord. »
\end{enumerate}
\end{propriete}

\begin{center}
\begin{tabularx}{\linewidth}{|X|X|X|}
\hline
\rowcolor{popblueL} Français & Anglais & Remarque \\
\hline
passer (réussir) un entretien & to pass an interview & to do well in / to succeed in an interview \\
\hline
une formation & a formation & a training (course), education \\
\hline
une candidature & a candidature & an application \\
\hline
assister à & to assist & to attend \\
\hline
une expérience (professionnelle) & an experiment & (work) experience \\
\hline
j'ai vingt ans & I have twenty years & I am twenty years old \\
\hline
je suis d'accord & I am agree & I agree \\
\hline
actuellement & actually (= en fait) & currently, at present \\
\hline
un stage & a stage (= une scène, une étape) & an internship, a work placement \\
\hline
\end{tabularx}
\end{center}

\begin{popfigure}
\begin{tikzpicture}[font=\small]
\draw[fill=popblueL, draw=popblue, thick] (0,0) rectangle (3.4,1.0);
\node[align=center] at (1.7,0.5) {\textbf{MOT FRANÇAIS}};
\draw[fill=poporangeL, draw=poporange, thick] (4.2,0) rectangle (7.6,1.0);
\node[align=center] at (5.9,0.5) {\textbf{FAUX AMI (✗)}};
\draw[fill=popgreenL, draw=popgreen, thick] (8.4,0) rectangle (11.8,1.0);
\node[align=center] at (10.1,0.5) {\textbf{MOT CORRECT (✓)}};
\foreach \y/\a/\b/\c in {
 -0.8/{formation}/{formation}/{training},
 -1.8/{candidature}/{candidature}/{application},
 -2.8/{assister à}/{to assist}/{to attend},
 -3.8/{passer un entretien}/{to pass an interview}/{to do well in an interview},
 -4.8/{j'ai 20 ans}/{I have 20 years}/{I am 20 years old}}{
 \node[align=center] at (1.7,\y) {\a};
 \node[align=center, text=poporange] at (5.9,\y) {\b};
 \node[align=center, text=popgreen] at (10.1,\y) {\c};
 \draw[-{Stealth}, thick, popmuted] (3.5,\y) -- (4.1,\y);
 \draw[-{Stealth}, thick, popmuted] (7.7,\y) -- (8.3,\y);
}
\end{tikzpicture}
\legende{Du mot français au mot anglais correct : toujours franchir l'étape de vérification.}
\end{popfigure}

\begin{methode}[Comment éviter les calques : pensez en anglais]
\begin{enumerate}
\item \textbf{Ne traduisez pas mot à mot.} Construisez la phrase directement en anglais, en partant d'un modèle anglais que vous connaissez (\eng{I am... years old}, \eng{I applied for...}).
\item \textbf{Vérifiez mentalement chaque verbe} que vous avez traduit du français : ce verbe anglais signifie-t-il vraiment ce que je pense ? (\eng{assist} ? \eng{pass} ?)
\item \textbf{Tenez un cahier personnel de faux amis} : deux colonnes, piège français / anglais correct, et ajoutez chaque nouveau piège rencontré.
\item \textbf{Relisez votre production écrite} une fois, en ne cherchant que les calques : âge avec \eng{have}, \eng{I am agree}, \eng{actually} pour « actuellement ».
\item \textbf{Apprenez les mots en collocations}, pas isolés : \eng{attend an interview}, \eng{gain experience}, \eng{submit an application}. Une collocation mémorisée comme un bloc ne peut pas être trahie par le français.
\end{enumerate}
\end{methode}

\begin{attention}[Au Bac]
Les correcteurs du Baccalauréat pénalisent systématiquement les calques du français, surtout dans les épreuves de dissertation et de rédaction de lettre. Une phrase comme ✗ \eng{I have twenty years and I passed the interview} signale un candidat faible même si les idées sont bonnes. Avant de rendre votre copie, prenez deux minutes pour vérifier uniquement les pièges du français : âge, \eng{agree}, \eng{actually}, \eng{assist}, \eng{formation}, \eng{experience/experiment}.
\end{attention}

\begin{exoresolu}[Corrigez les phrases de Marie]
Énoncé — Réécrivez chaque phrase en anglais correct. 1. \eng{I have twenty-two years.} 2. \eng{I passed the interview last week.} 3. \eng{I have a good formation in accounting.} 4. \eng{I assisted to a training session.} 5. \eng{I sent my candidature in January.} 6. \eng{I am agree with you.}
\tcblower
Solution détaillée —
1. \eng{I am twenty-two years old.} (l'âge se construit avec \eng{to be} + \eng{years old}).
2. \eng{I did well in the interview last week.} / \eng{I succeeded in the interview last week.} (\eng{pass} se réserve aux examens).
3. \eng{I have good training in accounting.} / \eng{I was well trained in accounting.} (\eng{formation} = formation rocheuse ou militaire).
4. \eng{I attended a training session.} (\eng{attend} + complément direct, sans \eng{to} ; \eng{assist} = aider).
5. \eng{I sent my application in January.} (\eng{candidature} n'existe pas en anglais courant).
6. \eng{I agree with you.} (\eng{agree} est un verbe ; pas de \eng{am} devant).
\end{exoresolu}

\begin{exercice}[À vous — trouvez et corrigez les pièges]
Chaque phrase contient un piège du français. Soulignez-le et réécrivez la phrase correctement.
\begin{enumerate}
\item \eng{Actually, I am unemployed, but I am looking for a job in Douala.}
\item \eng{She passed her interview and signed the contract.}
\item \eng{He did an experiment of ten years in this company.}
\item \eng{We assisted to the recruitment meeting yesterday.}
\item \eng{My candidature was accepted by the manager.}
\item \eng{I have eighteen years and I am agree to work at weekends.}
\end{enumerate}
\end{exercice}

\begin{corrige}[Exercice : trouvez et corrigez les pièges]
1. \eng{Actually} (calque de « actuellement ») → \eng{At present, I am unemployed, but I am looking for a job in Douala.}
2. \eng{passed her interview} → \eng{She did well in her interview and signed the contract.}
3. \eng{an experiment} → \eng{He has ten years of experience in this company.}
4. \eng{assisted to} → \eng{We attended the recruitment meeting yesterday.}
5. \eng{candidature} → \eng{My application was accepted by the manager.}
6. Deux pièges : \eng{I have eighteen years} → \eng{I am eighteen years old} ; \eng{I am agree} → \eng{I agree to work at weekends.}
\end{corrige}

\begin{savaistu}[Pourquoi les faux amis existent-ils ?]
L'anglais et le français partagent des milliers de mots car, après 1066, le français est devenu la langue de la cour d'Angleterre pendant trois siècles. Mais les sens ont divergé avec le temps : \eng{assist} a gardé le sens ancien d'« aider », tandis que le français « assister (à) » a pris le sens d'« être présent ». Fait intéressant, les candidats bilingues camerounais ont un avantage lors des vrais entretiens : ils peuvent basculer entre les deux langues — mais seulement s'ils maintiennent fermement séparés les faux amis de chaque langue !
\end{savaistu}

\begin{aretenir}[Points clés]
\begin{itemize}
\item Un \cle{faux ami} ressemble au français mais a un sens différent en anglais.
\item « Réussir un entretien » = \eng{to do well in / to succeed in an interview} ; \eng{pass} = réussir un \emph{examen}.
\item « Formation » = \eng{training / education} ; « candidature » = \eng{application} ; « assister à » = \eng{attend} (sans préposition).
\item \eng{experience} = expérience professionnelle ; \eng{experiment} = expérience scientifique.
\item Calques de structure : \eng{I am twenty years old}, \eng{I agree}.
\item Stratégie anti-calque : penser la phrase directement en anglais et vérifier chaque verbe traduit.
\end{itemize}
\end{aretenir}

\coursec{Word families, pronunciation and spelling}

Dans la section précédente, nous avons débusqué les faux amis et les pièges de traduction qui guettent les francophones (\eng{to apply for}, \eng{a candidate}, \eng{a degree}...). Nous allons maintenant adopter une stratégie plus puissante encore : au lieu d'apprendre les mots un par un, nous allons les apprendre \cle{par familles}. En anglais comme en français, une même racine donne naissance à un verbe, à des noms et à des adjectifs. Connaître ces familles, c'est multiplier son vocabulaire par quatre ou cinq pour un même effort de mémorisation. Nous terminerons par la prononciation et l'orthographe, deux domaines où les mots de l'entretien d'embauche réservent de mauvaises surprises.

\courssub{Why learn words in families?}

Observez d'abord ce court extrait d'article, puis répondez mentalement à la question : combien de mots différents sont construits sur la même racine ?

\begin{textebox}[From a careers column — “The jobs market today”]
Every year, thousands of young Cameroonians \textbf{apply} for jobs in banks, NGOs and private companies. Some \textbf{applicants} send more than fifty \textbf{applications} before they receive a positive answer. When an \textbf{employer} finally invites a candidate to an \textbf{interview}, the \textbf{interviewee} is often nervous, while the \textbf{interviewer} tries to create a relaxed atmosphere. The questions usually concern the candidate's \textbf{qualifications} and previous \textbf{employment}. Young people who are \textbf{qualified} but \textbf{unemployed} sometimes lose hope, yet a well-prepared \textbf{application} can change everything.
\end{textebox}

\begin{definition}[Glossaire]
\begin{itemize}
\item \eng{an applicant} (n.) : un candidat, un postulant (la personne qui postule)
\item \eng{an application} (n.) : une candidature, un dossier de candidature (la chose)
\item \eng{an interviewee} (n.) : la personne interrogée, le candidat interrogé
\item \eng{employment} (n.) : l'emploi, le travail salarié
\item \eng{qualified} (adj.) : qualifié, diplômé, compétent
\end{itemize}
\end{definition}

Vous l'avez remarqué : \eng{apply}, \eng{applicant}, \eng{application} partagent la racine \emph{appl-} ; \eng{employer}, \eng{employment}, \eng{unemployed} partagent \emph{employ-} ; \eng{interview}, \eng{interviewer}, \eng{interviewee} partagent \emph{interview-}. C'est exactement le même phénomène qu'en français avec \emph{employer / employeur / employé / emploi}. La bonne nouvelle : les mécanismes de formation sont presque identiques dans les deux langues. La mauvaise nouvelle : les suffixes ne se correspondent pas toujours, et c'est là que naissent les erreurs.

\courssub{The APPLY family}

\begin{propriete}[La famille de APPLY]
\begin{center}
\begin{tabular}{|l|l|l|}
\hline
\rowcolor{popblueL} Mot anglais & Nature & Traduction française \\
\hline
\eng{to apply (for)} & verbe & postuler (à), faire une demande (de) \\
\hline
\eng{an applicant} & nom (personne) & un candidat, un postulant \\
\hline
\eng{an application} & nom (chose) & une candidature, une demande \\
\hline
\eng{applicable} & adjectif & applicable, valable \\
\hline
\end{tabular}
\end{center}
\end{propriete}

\begin{exemplebox}[La famille APPLY en contexte]
\eng{Over two hundred applicants applied for the post, but only ten applications were complete.} « Plus de deux cents candidats ont postulé au poste, mais seulement dix dossiers étaient complets. »

\eng{This rule is not applicable to part-time workers.} « Cette règle ne s'applique pas aux travailleurs à temps partiel. »
\end{exemplebox}

\begin{attention}[Le verbe se construit avec FOR]
On dit \eng{to apply \textbf{for} a job} (postuler à un emploi), jamais \eng{to apply a job}. En revanche, on dit \eng{to apply \textbf{to} a company} (s'adresser à une entreprise). Retenez : \cle{apply FOR the job, TO the company}.
\end{attention}

\courssub{The EMPLOY family}

C'est la famille la plus riche et la plus utile du thème. Observez l'arbre lexical ci-dessous : d'une seule racine, l'anglais tire le verbe, les deux personnes de la relation de travail, le nom abstrait et les adjectifs.

\begin{popfigure}
\begin{center}
\begin{tikzpicture}[
  grow=down,
  level distance=1.6cm,
  level 1/.style={sibling distance=3.4cm},
  every node/.style={draw, rounded corners, align=center, font=\small},
  edge from parent/.style={draw=popblue, thick, -{Stealth}}
]
\node[fill=popblue, text=white, font=\bfseries, align=center]{EMPLOY\\ \emph{to employ} (v.)}
  child { node[fill=poporangeL, align=center]{employer\\ \emph{l'employeur}\\ (celui qui emploie)} }
  child { node[fill=popgreenL, align=center]{employee\\ \emph{l'employé}\\ (celui qui est employé)} }
  child { node[fill=poppurpleL, align=center]{employment\\ \emph{l'emploi}\\ (nom abstrait)} }
  child { node[fill=poppinkL, align=center]{unemployed\\ \emph{au chômage}\\ (adjectif)} };
\end{tikzpicture}
\end{center}
\legende{Arbre lexical de la racine EMPLOY : un verbe, deux personnes, un nom, un adjectif.}
\end{popfigure}

\begin{propriete}[La famille de EMPLOY]
\begin{center}
\begin{tabular}{|l|l|l|}
\hline
\rowcolor{popblueL} Mot anglais & Nature & Traduction française \\
\hline
\eng{to employ} & verbe & employer, embaucher \\
\hline
\eng{an employer} & nom (personne) & un employeur, un patron \\
\hline
\eng{an employee} & nom (personne) & un employé, un salarié \\
\hline
\eng{employment} & nom (abstrait) & l'emploi, le travail \\
\hline
\eng{unemployment} & nom (abstrait) & le chômage \\
\hline
\eng{employed / unemployed} & adjectifs & employé / au chômage, sans emploi \\
\hline
\end{tabular}
\end{center}
\end{propriete}

\begin{exemplebox}[Exemples traduits]
\eng{The company employs fifty people in Douala.} « L'entreprise emploie cinquante personnes à Douala. »

\eng{Unemployment is high among graduates.} « Le chômage est élevé parmi les diplômés. »

\eng{She has been unemployed for six months, so she is looking for employment.} « Elle est au chômage depuis six mois, alors elle cherche du travail. »
\end{exemplebox}

\begin{propriete}[Les suffixes -er / -or et -ee : qui fait, qui subit]
En anglais, les suffixes \cle{-er} et \cle{-or} désignent \textbf{celui qui fait l'action} : \eng{interviewer} = celui qui interroge, \eng{employer} = celui qui emploie, \eng{trainer} = celui qui forme, \eng{director} = celui qui dirige.

Le suffixe \cle{-ee} désigne \textbf{celui qui subit l'action} : \eng{interviewee} = celui qu'on interroge, \eng{employee} = celui qu'on emploie, \eng{trainee} = le stagiaire (celui qu'on forme), \eng{payee} = le bénéficiaire (celui qu'on paie).

Même logique en français : \emph{employ\underline{eur}} / \emph{employ\underline{é}}, \emph{form\underline{ateur}} / \emph{form\underline{é}}. Le suffixe français \emph{-é} correspond exactement à l'anglais \emph{-ee} : c'est un allié, pas un piège !
\end{propriete}

\begin{attention}[Employer ou employee ? L'inversion classique]
Les francophones inversent très souvent ces deux mots, car \eng{employer} ressemble au verbe français \emph{employer} et \eng{employee} semble « bizarre ». Or :

\eng{My employer gave me a contract.} = « Mon employeur m'a donné un contrat. » (le patron)

\eng{The employees are on strike.} = « Les employés sont en grève. » (les salariés)

Astuce : \cle{-ee = subit = le salarié} (celui qui « subit » les ordres). Ne dites jamais \eng{I am an employer} pour dire « je suis employé » !
\end{attention}

\courssub{The QUALIFY and INTERVIEW families}

\begin{propriete}[La famille de QUALIFY]
\begin{center}
\begin{tabular}{|l|l|l|}
\hline
\rowcolor{popblueL} Mot anglais & Nature & Traduction française \\
\hline
\eng{to qualify} & verbe & obtenir une qualification, être qualifié pour \\
\hline
\eng{a qualification} & nom & un diplôme, une qualification \\
\hline
\eng{qualified} & adjectif & qualifié, diplômé, compétent \\
\hline
\end{tabular}
\end{center}
\end{propriete}

\begin{exemplebox}[Exemples traduits]
\eng{She qualified as a nurse in 2019.} « Elle a obtenu son diplôme d'infirmière en 2019. »

\eng{What qualifications do you need for this job?} « Quels diplômes faut-il pour ce poste ? »

\eng{He is well qualified for the position.} « Il est bien qualifié pour ce poste. » (attention : \eng{qualified FOR}, avec la préposition \emph{for})
\end{exemplebox}

\begin{propriete}[La famille de INTERVIEW]
\begin{center}
\begin{tabular}{|l|l|l|}
\hline
\rowcolor{popblueL} Mot anglais & Nature & Traduction française \\
\hline
\eng{to interview} & verbe & interviewer, faire passer un entretien à \\
\hline
\eng{an interview} & nom & un entretien (d'embauche), une interview \\
\hline
\eng{an interviewer} & nom (personne) & le recruteur, celui qui interroge \\
\hline
\eng{an interviewee} & nom (personne) & le candidat interrogé \\
\hline
\end{tabular}
\end{center}
\end{propriete}

\begin{exemplebox}[Exemples traduits]
\eng{The interviewer asked the interviewee about her qualifications.} « Le recruteur a interrogé la candidate sur ses diplômes. »

\eng{They interviewed five candidates for the job.} « Ils ont fait passer un entretien à cinq candidats pour ce poste. » (le verbe est \cle{transitif direct} : pas de préposition après \eng{to interview someone})
\end{exemplebox}

\begin{savaistu}[Pourquoi -ee se prononce « i » long ?]
Le suffixe \emph{-ee} vient du participe passé français \emph{-é} (employé, formé), arrivé en anglais avec les Normands après 1066. L'anglais a gardé l'idée de « personne qui subit » mais a transformé la prononciation en « i » long : \eng{employee} se dit « em-plo-i-i », \eng{interviewee} se dit « in-teur-viou-i ». Le français et l'anglais partagent donc la même logique depuis presque mille ans !
\end{savaistu}

\courssub{Pronunciation: say it like an anglophone}

Le vocabulaire de l'entretien contient des mots dont la prononciation ne se devine pas à partir de l'orthographe. Voici les cinq mots essentiels, notés en lettres françaises, avec la syllabe accentuée en majuscules.

\begin{center}
\begin{tabularx}{\linewidth}{|l|X|X|}
\hline
\rowcolor{popblueL} Mot & Prononciation (lettres françaises) & Syllabe accentuée \\
\hline
\eng{interview} & « IN-teur-viou » & \textbf{IN}terview (1\iere{} syllabe) \\
\hline
\eng{candidate} & « KAN-di-deit » & \textbf{CAN}didate (1\iere{} syllabe) \\
\hline
\eng{appointment} & « a-POINT-ment » & ap\textbf{POINT}ment (2\ieme{} syllabe) \\
\hline
\eng{nervous} & « NEUR-veuss » & \textbf{NER}vous (1\iere{} syllabe) \\
\hline
\eng{experience} & « iks-PI-rienss » & ex\textbf{PE}rience (2\ieme{} syllabe) \\
\hline
\end{tabularx}
\end{center}

\begin{methode}[La méthode des trois temps pour bien prononcer]
Pour chaque mot nouveau, suivez toujours les trois étapes :
\begin{enumerate}
\item \textbf{Lentement} : décomposez le mot syllabe par syllabe en frappant la syllabe accentuée : « IN... teur... viou ».
\item \textbf{Normalement} : dites le mot à vitesse naturelle, cinq fois de suite : « INterview, INterview, INterview... »
\item \textbf{Dans une phrase} : placez le mot dans une phrase complète et lisez-la à voix haute : \eng{I have an INterview on Monday.} « J'ai un entretien lundi. »
\end{enumerate}
Marquez toujours la syllabe accentuée dans votre cahier avec des majuscules : \cle{INterview}, \cle{CANdidate}, \cle{exPErience}. En anglais, un mot mal accentué peut devenir incompréhensible, contrairement au français où l'accent est presque plat.
\end{methode}

\begin{attention}[Pièges de prononciation]
\begin{itemize}
\item \eng{candidate} : la dernière syllabe se dit « deit » (comme \eng{date}), pas « datte ».
\item \eng{experience} : le \emph{x} se prononce « ks » (« iks-pi-rienss ») et l'accent tombe sur \emph{-PE-}, jamais sur la première syllabe comme en français \emph{expérience}.
\item \eng{nervous} : le \emph{-ous} final se dit « euss » très bref, pas « ouss ».
\end{itemize}
\end{attention}

\courssub{Spelling: the tricky words}

Quatre mots du thème posent des problèmes d'orthographe, y compris aux anglophones. Apprenez-les avec leurs astuces de mémorisation.

\begin{propriete}[Orthographe piégeuse]
\begin{center}
\begin{tabularx}{\linewidth}{|l|X|X|}
\hline
\rowcolor{popblueL} Mot & Piège & Astuce \\
\hline
\eng{interview} & double voyelle : \emph{ie} au début, \emph{ew} à la fin & « \textbf{ie} ... \textbf{ew} » : deux voyelles qui se regardent \\
\hline
\eng{necessary} & un seul \emph{c}, deux \emph{s} & « one \textbf{C}ollar, two \textbf{S}leeves » (un col, deux manches) \\
\hline
\eng{successful} & deux \emph{c}, deux \emph{s} & « \textbf{su\-\textbf{cc}\-ess\-ful} » : tout est double sauf le \emph{u} initial \\
\hline
\eng{experience} & se termine en \emph{-ence}, pas \emph{-ance} & comme en français \emph{expérience} : \emph{-ence} \\
\hline
\end{tabularx}
\end{center}
\end{propriete}

\begin{attention}[Erreurs fréquentes à la copie]
\begin{itemize}
\item \eng{intreview} ou \eng{intervew} : faux. Il faut \emph{inter-view} : \eng{inter} + \eng{view} (une « vue réciproque »).
\item \eng{neccessary} : faux. Un seul \emph{c}. En revanche \eng{successful} a bien deux \emph{c} et deux \emph{s}.
\item \eng{experiance} : faux. La fin est \emph{-ence}, comme en français.
\item \eng{apointment} : faux. Deux \emph{p} : \eng{appointment} (de \eng{appoint} + \eng{-ment}).
\end{itemize}
\end{attention}

\courssub{Practice: solved exercise and consolidation}

\begin{exoresolu}[Complétez avec le bon mot de la famille]
Complétez chaque phrase avec la forme correcte du mot entre parenthèses.

1. The \dots\ (employ) asked the \dots\ (interview) why she had left her previous job.

2. More than fifty \dots\ (apply) sent their \dots\ (apply) for one single post.

3. He is highly \dots\ (qualify) : he has two university \dots\ (qualify).

4. \dots\ (employ) among young graduates is a serious problem in many countries.

\tcblower
\textbf{Solution détaillée.}

1. \eng{The \textbf{employer} asked the \textbf{interviewee} why she had left her previous job.} — Le suffixe \emph{-er} désigne celui qui fait l'action (le patron), le suffixe \emph{-ee} celle qui la subit (la candidate interrogée). « L'employeur a demandé à la candidate pourquoi elle avait quitté son emploi précédent. »

2. \eng{More than fifty \textbf{applicants} sent their \textbf{applications} for one single post.} — \emph{Applicant} = la personne (au pluriel avec \emph{fifty}), \emph{application} = la chose, le dossier (au pluriel avec \emph{their}). « Plus de cinquante candidats ont envoyé leurs dossiers pour un seul poste. »

3. \eng{He is highly \textbf{qualified} : he has two university \textbf{qualifications}.} — L'adjectif \emph{qualified} suit le verbe \emph{is} ; le nom \emph{qualifications} suit le chiffre \emph{two}, donc pluriel. « Il est hautement qualifié : il a deux diplômes universitaires. »

4. \eng{\textbf{Unemployment} among young graduates is a serious problem.} — Le nom abstrait avec le préfixe négatif \emph{un-} ; le verbe \emph{is} au singulier confirme qu'il faut un nom singulier indénombrable. « Le chômage chez les jeunes diplômés est un problème grave. »
\end{exoresolu}

\begin{exercice}[À vous de jouer]
\textbf{A. Familles de mots.} Donnez la forme demandée :
\begin{enumerate}
\item le nom de la personne qui passe l'entretien (racine \emph{interview}) ;
\item l'adjectif contraire de \eng{employed} ;
\item le nom de la chose construit sur \emph{apply} ;
\item le verbe de la famille de \emph{qualification}.
\end{enumerate}
\textbf{B. Orthographe.} Corrigez si nécessaire : \eng{neccessary}, \eng{succesful}, \eng{experiance}, \eng{intervew}.
\textbf{C. Prononciation.} Soulignez la syllabe accentuée : \eng{candidate}, \eng{appointment}, \eng{experience}, \eng{interview}.
\end{exercice}

\begin{corrige}[Exercice « À vous de jouer »]
\textbf{A.} 1. \eng{the interviewee} (celui qui subit l'entretien). 2. \eng{unemployed}. 3. \eng{an application}. 4. \eng{to qualify}.

\textbf{B.} \eng{necessary} (un seul \emph{c}) ; \eng{successful} (deux \emph{c}, deux \emph{s}) ; \eng{experience} (\emph{-ence}) ; \eng{interview} (\emph{ie}... \emph{ew}).

\textbf{C.} \textbf{CAN}didate ; ap\textbf{POINT}ment ; ex\textbf{PE}rience ; \textbf{IN}terview.
\end{corrige}

\begin{aretenir}[L'essentiel de la section]
\begin{itemize}
\item Apprenez le vocabulaire par \cle{familles} : \eng{apply / applicant / application}, \eng{employ / employer / employee / employment / unemployed}, \eng{qualify / qualification / qualified}, \eng{interview / interviewer / interviewee}.
\item \cle{-er / -or} = celui qui fait l'action ; \cle{-ee} = celui qui la subit (comme le français \emph{-é}).
\item Prononciation : méthode des trois temps (lent, normal, phrase) et syllabe accentuée marquée en majuscules.
\item Orthographe : \eng{necessary} (1 \emph{c}, 2 \emph{s}), \eng{successful} (2 \emph{c}, 2 \emph{s}), \eng{experience} (\emph{-ence}), \eng{interview} (\emph{ie}... \emph{ew}).
\end{itemize}
\end{aretenir}

\coursec{The stages of a job interview}

Un entretien d'embauche (\eng{job interview}) en milieu anglophone suit une structure codifiée. La connaître permet d'anticiper les questions, de gérer son temps de parole et d'adopter le registre formel attendu. Au Cameroun, comme au Nigeria ou au Ghana, les entretiens de la fonction publique et des grandes entreprises se déroulent souvent devant un \eng{panel} (jury) de trois à cinq examinateurs.

\begin{definition}[Key terms: stages of the interview]
\begin{itemize}
\item \eng{application} (nom) : candidature, dossier de candidature.
\item \eng{to apply for a post} (verbe) : postuler à un poste.
\item \eng{shortlist} (verbe/nom) : présélectionner / liste restreinte.
\item \eng{shortlisted candidate} (nom) : candidat présélectionné.
\item \eng{panel} (nom) : jury, commission d'entretien.
\item \eng{interviewer} (nom) : l'examinateur, celui qui pose les questions.
\item \eng{interviewee} (nom) : le candidat, celui qui répond.
\item \eng{appointment} (nom) : nomination, rendez-vous.
\end{itemize}
\end{definition}

\begin{center}
\begin{tabularx}{\linewidth}{|l|X|X|}
\hline
\rowcolor{popblueL} \textbf{Étape} & \textbf{Anglais} & \textbf{Français} \\
\hline
1 & \eng{Application \& shortlisting} & Dépôt de candidature et présélection \\
\hline
2 & \eng{Invitation to interview} & Convocation à l'entretien \\
\hline
3 & \eng{Greeting \& introduction} & Accueil et présentations \\
\hline
4 & \eng{Core questioning} & Questions principales (compétences, motivation) \\
\hline
5 & \eng{Candidate's questions} & Questions du candidat \\
\hline
6 & \eng{Closing \& follow-up} & Clôture et suites données \\
\hline
\end{tabularx}
\end{center}

\begin{popfigure}
\begin{center}
\begin{tikzpicture}[
stg/.style={rectangle, rounded corners=4pt, draw=popblue, fill=popblueL, text width=2.8cm, align=center, minimum height=0.85cm, font=\small\bfseries},
arr/.style={-{Stealth[length=2.5mm]}, thick, popdark}
]
\node[stg] (a) at (0,0) {1. Application};
\node[stg, fill=popgreenL, draw=popgreen] (b) at (0,-1.3) {2. Invitation};
\node[stg, fill=poporangeL, draw=poporange] (c) at (0,-2.6) {3. Greeting};
\node[stg, fill=poppurpleL, draw=poppurple] (d) at (0,-3.9) {4. Core questions};
\node[stg, fill=poppinkL, draw=poppink] (e) at (0,-5.2) {5. Your questions};
\node[stg, fill=popgoldL, draw=popgold] (f) at (0,-6.5) {6. Closing};
\draw[arr] (a) -- (b);
\draw[arr] (b) -- (c);
\draw[arr] (c) -- (d);
\draw[arr] (d) -- (e);
\draw[arr] (e) -- (f);
\end{tikzpicture}
\end{center}
\legende{Les six étapes du processus de recrutement}
\end{popfigure}

\coursec{People and documents}

Cette section présente le vocabulaire des acteurs et des documents échangés avant, pendant et après l'entretien. Maîtriser ces termes est indispensable pour comprendre une convocation, remplir un formulaire ou rédiger un courriel de suivi.

\begin{definition}[People involved]
\begin{itemize}
\item \eng{applicant} / \eng{candidate} (nom) : candidat, postulant.
\item \eng{recruiter} / \eng{hiring manager} (nom) : recruteur, responsable du recrutement.
\item \eng{panel member} / \eng{interviewer} (nom) : membre du jury, examinateur.
\item \eng{HR officer} / \eng{HR manager} (nom) : responsable des ressources humaines.
\item \eng{referee} / \eng{reference} (nom) : personne de référence (qui donne une recommandation).
\item \eng{successful candidate} (nom) : candidat retenu, lauréat.
\item \eng{unsuccessful candidate} (nom) : candidat non retenu.
\end{itemize}
\end{definition}

\begin{definition}[Documents]
\begin{itemize}
\item \eng{CV} / \eng{curriculum vitae} (nom) : curriculum vitæ (au Royaume-Uni, au Cameroun, au Nigeria).
\item \eng{resumé} / \eng{resume} (nom) : CV (aux États-Unis, au Canada) — \textbf{attention à l'accent !}
\item \eng{cover letter} / \eng{letter of application} (nom) : lettre de motivation.
\item \eng{application form} (nom) : formulaire de candidature.
\item \eng{qualifications} (nom pl.) : diplômes, titres.
\item \eng{certificates} (nom pl.) : certificats, attestations.
\item \eng{transcript} (nom) : relevé de notes.
\item \eng{reference letter} / \eng{testimonial} (nom) : lettre de recommandation.
\item \eng{appointment letter} / \eng{offer letter} (nom) : lettre d'engagement, lettre d'offre.
\item \eng{rejection letter} (nom) : lettre de refus.
\end{itemize}
\end{definition}

\begin{attention}[CV vs Resumé — faux ami \eng{resumé}]
\begin{itemize}
\item \eng{CV} (British / Commonwealth) = \eng{curriculum vitae} complet, détaillé, 2 pages ou plus.
\item \eng{Resumé} (American) = résumé d'une page, très synthétique.
\item Au Cameroun (système britannique), on utilise \eng{CV}. N'écrivez jamais \eng{resumé} sans accent (ça s'écrit \eng{resume} en anglais US, mais le mot français \eng{résumé} prend un accent aigu).
\item \eng{Resumé} (avec accent) n'existe pas en anglais ; c'est un mot français.
\end{itemize}
\end{attention}

\begin{exemplebox}[Typical email invitation]
\eng{Subject: Invitation to Interview -- Accountant Position (Ref: ACC/2025/042)}

\eng{Dear Miss Amina,}

\eng{We are pleased to inform you that you have been shortlisted for the position of Accountant. The interview will take place on Tuesday, 15 July 2025 at 9:00 a.m. at our Buea Head Office, Molyko.}

\eng{Please bring the originals of your certificates, a valid ID card, and two passport-sized photographs.}

\eng{Kind regards,}

\eng{Human Resources Department}
\end{exemplebox}

\coursec{Collocations and key expressions}

Les \cle{collocations} (associations de mots figées) sont la marque d'un anglais naturel. Les apprendre par blocs évite les calques du français (\eng{*pass an interview} au lieu de \eng{attend an interview}, \eng{*make a question} au lieu de \eng{ask a question}).

\begin{propriete}[Collocations essentielles -- entretien d'embauche]
\begin{center}
\begin{tabularx}{\linewidth}{|X|X|X|}
\hline
\rowcolor{popblueL} \textbf{Verbe + Nom} & \textbf{Adjectif + Nom} & \textbf{Expressions figées} \\
\hline
\eng{apply for a post} & \eng{shortlisted candidate} & \eng{to be shortlisted} \\
\eng{attend an interview} & \eng{panel interview} & \eng{to think on one's feet} \\
\eng{conduct an interview} & \eng{dream job} & \eng{to make a good impression} \\
\eng{answer questions} & \eng{relevant experience} & \eng{to land a job} (informel) \\
\eng{ask questions} & \eng{key strengths} & \eng{to get the sack} (informel) \\
\eng{offer a post} & \eng{competitive salary} & \eng{to be on the shortlist} \\
\eng{accept / decline an offer} & \eng{previous employer} & \eng{to follow up} \\
\eng{sign a contract} & \eng{working conditions} & \eng{to hear back} \\
\hline
\end{tabularx}
\end{center}
\end{propriete}

\begin{exemplebox}[Collocations en contexte]
\eng{She \textbf{applied for} the post of secretary.} « Elle a postulé au poste de secrétaire. »

\eng{He \textbf{attended} three interviews last month.} « Il a passé trois entretiens le mois dernier. »

\eng{The panel \textbf{conducted} the interview professionally.} « Le jury a mené l'entretien professionnellement. »

\eng{I hope to \textbf{make a good impression}.} « J'espère faire bonne impression. »

\eng{They \textbf{offered me} the job on the spot.} « Ils m'ont offert le poste sur-le-champ. »

\eng{We will \textbf{follow up} with a formal letter.} « Nous ferons un suivi par lettre formelle. »
\end{exemplebox}

\begin{attention}[Calques français $\to$ anglais correct]
\begin{center}
\begin{tabularx}{\linewidth}{|X|X|X|}
\hline
\rowcolor{poporangeL} \textbf{Français (calque)} & \textbf{Anglais incorrect} & \textbf{Anglais correct} \\
\hline
passer un entretien & \eng{*pass an interview} & \eng{attend an interview} / \eng{have an interview} \\
faire bonne impression & \eng{*make a good impression} & \eng{make a good impression} (correct !) \\
poser une question & \eng{*make a question} & \eng{ask a question} \\
être présélectionné & \eng{*be pre-selected} & \eng{be shortlisted} \\
un poste vacant & \eng{*a vacant post} & \eng{a vacancy} / \eng{a vacant post} (les deux existent) \\
lettre de motivation & \eng{*motivation letter} & \eng{cover letter} / \eng{letter of application} \\
\hline
\end{tabularx}
\end{center}
\end{attention}

\coursec{False friends and French traps}

Les faux amis sont la première cause de perte de points au Baccalauréat. Cette liste cible ceux qui apparaissent spécifiquement dans le thème de l'entretien d'embauche.

\begin{definition}[Faux amis -- Spécial entretien]
\begin{center}
\begin{tabularx}{\linewidth}{|l|l|l|X|}
\hline
\rowcolor{popblueL} \textbf{Mot anglais} & \textbf{Sens anglais} & \textbf{Faux ami français} & \textbf{Traduction correcte du mot français} \\
\hline
\eng{opportunity} & occasion favorable, chance & opportunité (caractère opportun) & \eng{suitability} / \eng{relevance} \\
\eng{candidate} & candidat (à un poste, examen) & candidat (seul sens) & \eng{candidate} (vrai ami) \\
\eng{application} & candidature, dossier & application (mise en œuvre) & \eng{implementation} / \eng{application} (informatique) \\
\eng{resume} (US) & CV & résumé (synthèse) & \eng{summary} \\
\eng{interview} & entretien (embauche, presse) & interview (journalistique seulement) & \eng{interview} (les deux sens) \\
\eng{panel} & jury, commission & panneau (tableau) & \eng{board} / \eng{sign} \\
\eng{post} & poste, emploi & poste (police, frontière) & \eng{post} / \eng{station} / \eng{border post} \\
\eng{reference} & recommandation, référence & référence (renvoi) & \eng{reference} (vrai ami) \\
\eng{qualification} & diplôme, titre requis & qualification (aptitude) & \eng{qualification} / \eng{skill} \\
\eng{experience} & expérience professionnelle & expérience (vécu, essai) & \eng{experience} (vrai ami) \\
\eng{salary} & salaire (mensuel/annuel) & salaire (terme générique) & \eng{salary} / \eng{wage} (horaire) \\
\eng{benefits} & avantages sociaux & bénéfices (profits) & \eng{profits} / \eng{earnings} \\
\eng{contract} & contrat (écrit, juridique) & contrat (accord oral aussi) & \eng{contract} / \eng{agreement} \\
\hline
\end{tabularx}
\end{center}
\end{definition}

\begin{attention}[Pièges grammaticaux fréquents]
\begin{itemize}
\item \eng{*I am agree} $\to$ \eng{I agree} (verbe, pas adjectif).
\item \eng{*I have 25 years} $\to$ \eng{I am 25} / \eng{I am 25 years old}.
\item \eng{*I am working since two years} $\to$ \eng{I have been working for two years} (present perfect continuous).
\item \eng{*The salary is 500.000 francs} $\to$ \eng{The salary is 500,000 francs} (virgule pour les milliers en anglais, point pour les décimales).
\item \eng{*I am interesting in this post} $\to$ \eng{I am \textbf{interested} in this post} (\eng{-ed} = sentiment, \eng{-ing} = caractéristique).
\item \eng{*I look forward to hear from you} $\to$ \eng{I look forward to \textbf{hearing} from you} (gérondif après \eng{to} préposition).
\end{itemize}
\end{attention}

\coursec{Word families, pronunciation and spelling}

Travailler les familles de mots permet de passer du nom au verbe, à l'adjectif ou à l'adverbe sans hésiter. La prononciation (accent tonique) change souvent selon la catégorie grammaticale.

\begin{propriete}[Familles de mots -- Tableau de dérivation]
\begin{center}
\begin{tabular}{|l|l|l|l|l|}
\hline
\rowcolor{popblueL} \textbf{Verbe} & \textbf{Nom (chose/action)} & \textbf{Nom (personne)} & \textbf{Adjectif} & \textbf{Adverbe} \\
\hline
\eng{apply} & \eng{application} & \eng{applicant} & \eng{applicable} & \eng{applicably} \\
\eng{appoint} & \eng{appointment} & \eng{appointee} & \eng{appointed} & --- \\
\eng{interview} & \eng{interview} & \eng{interviewer} / \eng{interviewee} & \eng{interview} (adj.) & --- \\
\eng{recruit} & \eng{recruitment} & \eng{recruit} / \eng{recruiter} & \eng{recruited} & --- \\
\eng{employ} & \eng{employment} & \eng{employer} / \eng{employee} & \eng{employed} / \eng{employable} & \eng{employably} \\
\eng{qualify} & \eng{qualification} & \eng{qualifier} & \eng{qualified} / \eng{qualifying} & \eng{qualifiedly} \\
\eng{experience} (n) & \eng{experience} & --- & \eng{experienced} & \eng{experientially} \\
\eng{recommend} & \eng{recommendation} & \eng{referee} & \eng{recommended} & \eng{recommendedly} \\
\eng{succeed} & \eng{success} & \eng{successor} & \eng{successful} & \eng{successfully} \\
\eng{fail} & \eng{failure} & --- & \eng{failed} / \eng{failing} & \eng{failingly} \\
\hline
\end{tabular}
\end{center}
\end{propriete}

\begin{definition}[Règles d'accentuation (word stress)]
\begin{itemize}
\item Verbes 2 syllabes : accent sur la 2\textsuperscript{e} : \eng{ap\emph{ply}} (« a-plaï »), \eng{re\emph{cruit}} (« ri-kroït »).
\item Noms/Adjectifs 2 syllabes : accent sur la 1\textsuperscript{re} : \eng{\emph{ap}plicant} (« a-pli-kent »), \eng{\emph{post}} (« peust »).
\item Suffixe \eng{-tion} / \eng{-sion} : accent sur la syllabe avant le suffixe : \eng{appli\emph{ca}tion} (« a-pli-ké-chn »), \eng{qualifi\emph{ca}tion}.
\item Suffixe \eng{-ee} (personne qui subit l'action) : accent sur \eng{-ee} : \eng{interview\emph{ee}} (« in-tèr-vyou-ï »), \eng{appointe\emph{e}}.
\item Composés : accent sur le premier élément : \eng{\emph{short}list}, \eng{\emph{cover} letter}, \eng{\emph{job} interview}.
\end{itemize}
\end{definition}

\begin{exemplebox}[Prononciation approximative (lettres françaises)]
\begin{itemize}
\item \eng{application} : « a-pli-ké-chn »
\item \eng{candidate} : « kan-di-deut »
\item \eng{interview} : « in-tèr-vyou »
\item \eng{qualification} : « quo-li-fi-ké-chn »
\item \eng{experienced} : « ik-spi-ri-enst »
\item \eng{successful} : « seuk-ses-ful »
\item \eng{appointment} : « a-point-ment »
\item \eng{remuneration} : « ri-myoo-né-reu-chn »
\end{itemize}
\end{exemplebox}

\begin{attention}[Orthographe -- Pièges courants]
\begin{itemize}
\item \eng{curriculum vitae} (Latin, invariable au pluriel : \eng{curricula vitae} ou \eng{CVs}).
\item \eng{resumé} (US) s'écrit \eng{resume} sans accent en anglais ; le participe passé français \eng{résumé} prend un accent.
\item \eng{reference} (un seul \eng{r}, deux \eng{e}) — pas \eng{*referrence}.
\item \eng{successful} (double \eng{c}, double \eng{s}, un seul \eng{l} final) — pas \eng{*sucessful}, \eng{*successfull}.
\item \eng{interviewer} / \eng{interviewee} (double \eng{e} final pour \eng{ee}).
\item \eng{accommodation} (double \eng{c}, double \eng{m}) — souvent demandé dans les conditions de travail.
\end{itemize}
\end{attention}

\coursec{Vocabulary in context: a model dialogue}

Dans les sections précédentes, nous avons travaillé sur les familles de mots (\eng{apply}, \eng{application}, \eng{applicant}), la prononciation, l'orthographe, les collocations et les faux amis du lexique de l'entretien d'embauche. Il est temps maintenant de rassembler tout ce vocabulaire dans une situation de communication authentique : un \cle{dialogue complet d'entretien d'embauche}. C'est en contexte que les mots prennent tout leur sens — et c'est en contexte que l'examinateur du baccalauréat les évalue.

\courssub{Reading: a complete job interview dialogue}

Lisez le dialogue suivant. Il se déroule à Buea, dans les locaux d'une entreprise qui recrute un comptable. Observez bien la structure : salutations, présentation, questions du jury, clôture.

\begin{textebox}[A job interview in Buea]
\textbf{Panel:} Good morning, Miss Amina. Please take a seat.

\textbf{Amina:} Good morning, sir. Thank you for this opportunity.

\textbf{Panel:} You were shortlisted among fifty applicants. Tell us about yourself.

\textbf{Amina:} I hold a GCE Advanced Level and a diploma in accounting. I worked for two years as a cashier in Buea.

\textbf{Panel:} What are your strengths?

\textbf{Amina:} I am reliable, hard-working and I can think on my feet.

\textbf{Panel:} And what about your weaknesses?

\textbf{Amina:} I sometimes spend too much time checking details, but I am learning to work faster.

\textbf{Panel:} Why do you want this post?

\textbf{Amina:} Because your company has an excellent reputation, and I believe I can make a good impression and grow here.

\textbf{Panel:} Do you have any questions for us?

\textbf{Amina:} Yes, sir. When will the successful candidate start work?

\textbf{Panel:} On the first of next month. Thank you. We shall contact you next week.

\textbf{Amina:} Thank you very much, ladies and gentlemen. Have a nice day.
\end{textebox}

\begin{definition}[Glossaire du dialogue]
\begin{itemize}
\item \eng{to take a seat} (verbe) : s'asseoir — formule polie d'accueil.
\item \eng{cashier} (nom) : caissier, caissière.
\item \eng{reliable} (adjectif) : fiable, digne de confiance.
\item \eng{to think on one's feet} (expression) : réfléchir vite, avoir la répartie.
\item \eng{post} (nom) : poste, emploi.
\item \eng{successful candidate} (nom) : le candidat retenu.
\item \eng{shortlisted} (participe passé) : présélectionné (sur la liste restreinte).
\item \eng{ladies and gentlemen} (expression) : mesdames, messieurs (formule de politesse collective).
\end{itemize}
\end{definition}

\courssub{Anatomy of the dialogue: spotting the five stages}

Un entretien d'embauche en anglais suit presque toujours la même architecture. Sachez repérer chaque étape dans le dialogue :

\begin{center}
\begin{tabularx}{\linewidth}{|l|X|X|}
\hline
\rowcolor{popblueL} \textbf{Étape} & \textbf{Fonction} & \textbf{Exemple du dialogue} \\
\hline
1. Greeting & Saluer et installer le candidat & \eng{Good morning, Miss Amina. Please take a seat.} \\
\hline
2. Self-presentation & Le candidat se présente & \eng{I hold a GCE Advanced Level and a diploma in accounting.} \\
\hline
3. Panel's questions & Le jury interroge & \eng{What are your strengths?} \\
\hline
4. Candidate's questions & Le candidat pose une question & \eng{When will the successful candidate start work?} \\
\hline
5. Closing & Clôture et remerciements & \eng{We shall contact you next week.} \\
\hline
\end{tabularx}
\end{center}

\begin{popfigure}
\begin{center}
\begin{tikzpicture}[
stage/.style={rectangle, rounded corners=4pt, draw=popblue, fill=popblueL, text width=3.4cm, align=center, minimum height=0.9cm, font=\small\bfseries},
fleche/.style={-{Stealth[length=2.5mm]}, thick, popdark}
]
\node[stage] (s1) at (0,0) {1. Greeting};
\node[stage, fill=popgreenL, draw=popgreen] (s2) at (0,-1.4) {2. Self-presentation};
\node[stage, fill=poporangeL, draw=poporange] (s3) at (0,-2.8) {3. Panel's questions};
\node[stage, fill=poppurpleL, draw=poppurple] (s4) at (0,-4.2) {4. Candidate's questions};
\node[stage, fill=popgoldL, draw=popgold] (s5) at (0,-5.6) {5. Closing};
\draw[fleche] (s1) -- (s2);
\draw[fleche] (s2) -- (s3);
\draw[fleche] (s3) -- (s4);
\draw[fleche] (s4) -- (s5);
\end{tikzpicture}
\end{center}
\legende{Les cinq étapes d'un entretien d'embauche en anglais}
\end{popfigure}

\begin{methode}[How to succeed in an interview dialogue]
\begin{enumerate}
\item \textbf{Greet formally} : commencez par \eng{Good morning, sir / madam.} et remerciez : \eng{Thank you for this opportunity.}
\item \textbf{Present yourself briefly} : diplômes (\eng{I hold a diploma in...}) et expérience (\eng{I worked for two years as...}).
\item \textbf{Answer with examples} : ne donnez pas seulement un adjectif (\eng{reliable}), illustrez-le par un fait concret.
\item \textbf{Ask one question} : cela montre votre intérêt (\eng{When will the successful candidate start?}).
\item \textbf{Thank and close politely} : \eng{Thank you very much. Have a nice day.}
\end{enumerate}
\end{methode}

\begin{exemplebox}[Key sentences translated]
\eng{Thank you for this opportunity.} « Merci pour cette opportunité. »

\eng{You were shortlisted among fifty applicants.} « Vous avez été présélectionnée parmi cinquante candidats. »

\eng{I can think on my feet.} « Je sais réfléchir vite / j'ai de la répartie. »

\eng{We shall contact you next week.} « Nous vous contacterons la semaine prochaine. »
\end{exemplebox}

\begin{attention}[Stay formal from beginning to end]
Pendant tout l'entretien, évitez l'anglais familier : jamais \eng{yeah}, \eng{wanna}, \eng{gonna}, \eng{hi guys}. Préférez \eng{yes}, \eng{want to}, \eng{going to}, \eng{good morning}. De même, ne dites pas \eng{I want this job} (trop direct) mais \eng{I am very interested in this post.} Enfin, attention au faux ami : \eng{opportunity} signifie « occasion favorable », pas simplement « opportunité » au sens de « caractère opportun ».
\end{attention}

\begin{savaistu}[Interview panels in Anglophone Africa]
Au Nigeria et au Ghana, comme dans les régions anglophones du Cameroun, les entretiens de la fonction publique se déroulent souvent devant un \eng{panel} de plusieurs examinateurs assis derrière une longue table. Il est donc poli de s'adresser à tous : \eng{sirs, madams} ou \eng{members of the panel}. Regarder un seul examinateur pendant toute la réponse est considéré comme un manque d'assurance.
\end{savaistu}

\courssub{Comprehension check}

\begin{exoresolu}[True or false — justify with the text]
\textbf{Énoncé.} Say whether the following statements are true or false, and correct the false ones.
\begin{enumerate}
\item The candidate arrived late.
\item Amina worked as a cashier for two years.
\item The panel asks about weaknesses to embarrass the candidate.
\end{enumerate}
\tcblower
\textbf{Solution détaillée.}
\begin{enumerate}
\item \textbf{False.} Nothing in the dialogue says Amina was late; the panel welcomes her politely: \eng{Good morning, Miss Amina. Please take a seat.}
\item \textbf{True.} Amina says: \eng{I worked for two years as a cashier in Buea.}
\item \textbf{False.} The panel asks about weaknesses to test the candidate's \emph{honesty} and \emph{self-awareness}, not to embarrass her. A good candidate mentions a real weakness and shows how she is improving, as Amina does: \eng{I am learning to work faster.}
\end{enumerate}
\end{exoresolu}

\begin{exercice}[Multiple choice and short answers]
\textbf{A. Choose the correct answer.}
\begin{enumerate}
\item Why does the panel ask about weaknesses?
\begin{enumerate}
\item To refuse the candidate immediately.
\item To test the candidate's honesty and capacity to improve.
\item To compare her with the other fifty applicants.
\end{enumerate}
\item What post is Amina applying for?
\begin{enumerate}
\item A teaching post. \item An accounting post. \item A cashier post in a bank.
\end{enumerate}
\end{enumerate}
\textbf{B. Answer in one short sentence.}
\begin{enumerate}
\item How many applicants were shortlisted?
\item What qualification does Amina hold besides her GCE Advanced Level?
\item When will the panel contact Amina?
\end{enumerate}
\end{exercice}

\begin{corrige}[Exercice : Multiple choice and short answers]
\textbf{A.} 1. b — to test honesty and capacity to improve. 2. b — an accounting post (she holds a diploma in accounting).

\textbf{B.} 1. \eng{Fifty applicants were shortlisted.} 2. \eng{She holds a diploma in accounting.} 3. \eng{They will contact her next week.}
\end{corrige}

\courssub{Guided re-use: gap-fill dialogue}

\begin{exercice}[Complete the dialogue]
Complete the dialogue with the following words: \eng{shortlisted}, \eng{panel}, \eng{strengths}, \eng{impression}, \eng{opportunity}.

\textbf{Panel:} Good morning. Please take a seat. You were (1) \dots\ among thirty applicants.

\textbf{Candidate:} Good morning. Thank you for this (2) \dots

\textbf{Panel:} The (3) \dots\ would like to know your main (4) \dots

\textbf{Candidate:} I am hard-working and punctual, and I hope to make a good (5) \dots\ on your company.
\end{exercice}

\begin{corrige}[Exercice : Complete the dialogue]
(1) \eng{shortlisted} — (2) \eng{opportunity} — (3) \eng{panel} — (4) \eng{strengths} — (5) \eng{impression}.
\end{corrige}

\courssub{Speaking: role-play in pairs}

\begin{experience}[Re-enact the interview]
Travaillez en binôme. L'élève A joue le \eng{panel}, l'élève B joue le candidat. Rejouez le dialogue en changeant :
\begin{itemize}
\item le poste : \eng{a secretary}, \eng{a nurse}, \eng{a computer technician}, \eng{a sales assistant} ;
\item les diplômes et l'expérience du candidat ;
\item une question du jury : \eng{Why should we hire you?} / \eng{Can you work under pressure?}
\end{itemize}
Respectez les cinq étapes de l'organigramme et restez formels du début à la fin. Puis échangez les rôles.
\end{experience}

\courssub{Writing: a successful interview in the past}

\begin{methode}[Writing a short paragraph (8--10 lines)]
Rédigez un court paragraphe racontant un entretien réussi \textbf{au passé} (prétérit simple). Suivez le plan :
\begin{enumerate}
\item Situation : \eng{Last month, I attended a job interview in Douala.}
\item Déroulement : \eng{The panel greeted me warmly and asked about my strengths.}
\item Réponses : \eng{I explained that I was reliable and hard-working.}
\item Conclusion : \eng{A week later, they contacted me and offered me the post.}
\end{enumerate}
Attention aux prétérits irréguliers : \eng{go} $\to$ \eng{went}, \eng{take} $\to$ \eng{took}, \eng{think} $\to$ \eng{thought}, \eng{make} $\to$ \eng{made}.
\end{methode}

\begin{exemplebox}[Model paragraph]
\eng{Last month, I attended a job interview at a bank in Douala. I arrived early and greeted the panel politely. They asked me about my qualifications and my strengths. I told them that I was reliable and that I could think on my feet. I also mentioned my two years of experience as a cashier in Buea. At the end, I asked one question about the post and thanked the panel for the opportunity. A week later, they contacted me and offered me the job. I was the successful candidate!}

« Le mois dernier, j'ai passé un entretien d'embauche dans une banque à Douala. Je suis arrivé tôt et j'ai salué le jury poliment. Ils m'ont interrogé sur mes diplômes et mes points forts. Je leur ai dit que j'étais fiable et que je savais réfléchir vite. J'ai aussi mentionné mes deux années d'expérience comme caissier à Buea. À la fin, j'ai posé une question sur le poste et j'ai remercié le jury pour cette opportunité. Une semaine plus tard, ils m'ont contacté et m'ont offert le poste. J'étais le candidat retenu ! »
\end{exemplebox}

\begin{aretenir}[What to remember]
Un dialogue d'entretien suit cinq étapes : \cle{greeting} $\to$ \cle{self-presentation} $\to$ \cle{panel's questions} $\to$ \cle{candidate's questions} $\to$ \cle{closing}. Restez formel (\eng{yes}, \eng{want to}, \eng{going to}), répondez avec des exemples concrets, posez au moins une question et terminez par des remerciements. Pour raconter l'entretien à l'écrit, utilisez le prétérit simple et les marqueurs temporels \eng{last month}, \eng{a week later}.
\end{aretenir}

\newpage

\coursec{Activité d'intégration}

\coursec{Lesson 3 — Vocabulary: Words and expressions related to undergoing a job interview}

\courssub{Objectifs de l'activité}

À la fin de cette activité d'intégration, tu seras capable de :
\begin{itemize}
\item réemployer à l'oral et à l'écrit le \cle{vocabulaire} complet de l'entretien d'embauche : les \cle{étapes} de l'entretien, les \cle{personnes} et les \cle{documents} impliqués ;
\item utiliser correctement les \cle{collocations} et expressions clés de la leçon : \eng{to apply for a job}, \eng{to attend an interview}, \eng{to hire a candidate}, \eng{to meet the requirements} ;
\item éviter les \cle{faux amis} et les calques du français (\eng{actually} ne veut pas dire « actuellement », \eng{a formation} ne veut pas dire « une formation ») ;
\item employer les mots de la même \cle{famille} (\eng{employ, employer, employee, employment, unemployed}) dans des phrases correctes ;
\item tenir un rôle dans une \cle{simulation} d'entretien d'embauche avec aisance, politesse et langage soutenu.
\end{itemize}

\courssub{Situation}

\begin{textebox}[Job advertisement — The Cameroon Post, Douala]
\textbf{CAMEROON SUNRISE BANK IS RECRUITING!}

Cameroon Sunrise Bank, a leading financial institution with branches in Douala, Yaoundé, Buea and Bamenda, is looking for a \textbf{junior customer service officer} for its Douala branch.

\textbf{Requirements:}
The ideal candidate must hold at least a GCE Advanced Level certificate or a baccalauréat, speak fluent English and French, have good communication skills, and be ready to work under pressure. Previous experience is an advantage but not compulsory.

\textbf{How to apply:}
Interested candidates should send their application letter, an up-to-date CV and photocopies of their certificates to the Human Resources Manager before 30th November. Shortlisted candidates will be invited to sit a written test, and successful applicants will then attend an interview with a panel of three interviewers.

\emph{Cameroon Sunrise Bank — an equal opportunity employer.}
\end{textebox}

\textbf{Mise en situation.} Tu es élève en Terminale A au lycée de Douala. Avec tes camarades, tu as repéré cette annonce dans le journal. La classe se transforme aujourd'hui en salle de recrutement : certains élèves joueront les \eng{candidates} (candidats), d'autres les membres du \eng{panel} (jury d'entretien) dirigé par le \eng{Human Resources Manager}. Chaque candidat devra se présenter, répondre aux questions et poser une question à la fin de l'entretien, en utilisant le vocabulaire étudié dans la leçon.

\courssub{Consignes}

\begin{enumerate}
\item \textbf{Compréhension de l'annonce.} Lis attentivement l'annonce ci-dessus. Souligne les mots de vocabulaire vus dans la leçon (\eng{to recruit, requirements, candidate, to apply, application letter, CV, shortlisted, interview, panel, employer}). Puis réponds en anglais : What position is advertised? What are the requirements? What documents must candidates send?

\item \textbf{Préparation du lexique.} En binôme, complète le tableau des familles de mots et des collocations utiles pour l'entretien :
\begin{center}
\begin{tabularx}{\linewidth}{|X|X|X|}
\hline
\rowcolor{popblueL} Verbe & Nom (personne) & Nom (chose ou notion) \\
\hline
to employ & employer / employee & employment \\
\hline
to apply & applicant & application \\
\hline
to interview & interviewer / interviewee & interview \\
\hline
to qualify & --- & qualification \\
\hline
\end{tabularx}
\end{center}

\item \textbf{Préparation des rôles.} Répartissez les rôles : un \eng{HR Manager} (président du jury), deux \eng{interviewers}, et quatre à six \eng{candidates}. Chaque candidat prépare une courte fiche : nom, diplôme (\eng{qualifications}), qualités (\eng{strengths}), expérience éventuelle (\eng{previous experience}), et une question à poser au jury.

\item \textbf{Simulation.} Chaque entretien dure trois à quatre minutes et suit les étapes de la leçon : \eng{greeting and introduction}, \eng{questions about qualifications and experience}, \eng{questions about strengths and weaknesses}, \eng{the candidate's questions}, \eng{closing the interview}.

\item \textbf{Évaluation par les pairs.} Pendant chaque entretien, les autres élèves remplissent une grille : le candidat a-t-il utilisé au moins cinq mots du champ lexical ? A-t-il évité les faux amis ? A-t-il été poli (\eng{Good morning, Sir. Thank you for this opportunity.}) ?

\item \textbf{Production écrite de bilan.} Après la simulation, rédige un court paragraphe (huit à dix lignes) intitulé \eng{My job interview}, dans lequel tu racontes ton entretien au passé en réemployant au moins huit mots ou expressions de la leçon.
\end{enumerate}

\courssub{Ressources à mobiliser}

\begin{aretenir}[Boîte à outils de l'entretien]
\begin{itemize}
\item \textbf{Étapes :} \eng{to send an application, to be shortlisted, to attend an interview, to answer questions, to be hired / to be rejected}.
\item \textbf{Personnes et documents :} \eng{applicant, candidate, interviewer, panel, HR Manager, employer, CV, application letter, certificate, reference}.
\item \textbf{Collocations :} \eng{to apply for a post, to meet the requirements, to work under pressure, to have good communication skills, to gain experience, to offer somebody a job}.
\item \textbf{Expressions de politesse :} \eng{Good morning, Madam. It is a pleasure to meet you. Thank you for giving me this opportunity. I look forward to hearing from you.}
\item \textbf{Familles de mots :} \eng{employ / employer / employee / employment / unemployed} ; \eng{apply / applicant / application}.
\end{itemize}
\end{aretenir}

\begin{attention}[Pièges à éviter pendant la simulation]
\begin{itemize}
\item Ne dis pas \eng{*I have a formation in banking} ; dis \eng{I have training in banking} ou \eng{I was trained in...} (« une formation » = \eng{training}).
\item \eng{Actually} signifie « en fait », pas « actuellement » : \eng{I currently live in Douala} (« j'habite actuellement à Douala »).
\item Ne dis pas \eng{*I passed an interview} au sens de « j'ai subi un entretien » ; dis \eng{I attended an interview} ou \eng{I had an interview}. \eng{To pass} s'emploie pour un examen : \eng{I passed the written test}.
\item Attention à la prononciation : \eng{employee} se prononce « èm-ploï-i » (accent sur la dernière syllabe) ; \eng{employer} se prononce « èm-ploï-eur ».
\end{itemize}
\end{attention}

\courssub{Proposition de production et corrigé commenté}

\begin{exoresolu}[Modèle d'entretien — The Dream Job Interview]
Énoncé : produis un dialogue modèle entre le HR Manager (M) et une candidate, Nadège (C), puis rédige le paragraphe de bilan \eng{My job interview}.

\tcblower

\textbf{Solution — dialogue modèle :}

\eng{M: Good morning. Please come in and take a seat.}

\eng{C: Good morning, Sir. Thank you for inviting me to this interview.}

\eng{M: You are welcome. Could you introduce yourself briefly?}

\eng{C: My name is Nadège Ekane. I obtained my Baccalauréat last year at Lycée de New Bell, and I am currently taking a training course in customer relations.}

\eng{M: Why did you apply for this post?}

\eng{C: I believe I meet all the requirements: I speak English and French fluently, I have good communication skills, and I can work under pressure.}

\eng{M: Do you have any previous experience?}

\eng{C: Yes. I worked as a shop assistant in my uncle's store during the holidays, so I gained some experience in dealing with customers.}

\eng{M: What are your strengths and weaknesses?}

\eng{C: I am hardworking and punctual. My weakness is that I sometimes speak too fast when I am nervous, but I am working on it.}

\eng{M: Very well. Do you have any questions for us?}

\eng{C: Yes, Sir. When will the successful candidates be informed of the results?}

\eng{M: Within two weeks. Thank you, Miss Ekane. We will contact you soon.}

\eng{C: Thank you for this opportunity. I look forward to hearing from you. Goodbye.}

\textbf{Solution — paragraphe de bilan :}

\eng{Last Tuesday, I attended a job interview at Cameroon Sunrise Bank. I had sent my application letter and my CV two weeks earlier, and I was shortlisted for the interview. The panel was made up of three interviewers, including the Human Resources Manager. They asked me questions about my qualifications and my previous experience. I was nervous at first, but I answered politely and confidently. At the end, they told me they would contact the successful candidates within two weeks. I really hope to be hired, because this job is a great opportunity for me.}

\textbf{Commentaire des choix de langue :}
\begin{itemize}
\item \eng{Thank you for inviting me} : expression de politesse attendue à l'ouverture ; on évite le calque \eng{*thank you to invite me}.
\item \eng{I am currently taking a training course} : emploi correct de \eng{currently} (« actuellement ») et de \eng{training} (« formation »), deux faux amis neutralisés.
\item \eng{I meet all the requirements} : collocation exacte étudiée dans la leçon.
\item \eng{I gained some experience} : \eng{experience} est ici indénombrable, sans article ni pluriel ; on évite \eng{*experiences}.
\item \eng{I attended a job interview} : verbe correct pour « subir un entretien » ; \eng{pass} est réservé aux examens.
\item Le paragraphe de bilan conjugue les verbes au \cle{prétérit} (\eng{attended, sent, was shortlisted, asked, answered, told}) : c'est le temps du récit d'un événement passé.
\item Familles de mots réemployées : \eng{interview / interviewers}, \eng{apply / application}, \eng{employ / hired}.
\end{itemize}
\end{exoresolu}

\courssub{Bilan de l'activité}

Cette simulation t'a permis de passer du vocabulaire reconnu à la lecture (l'annonce) au vocabulaire réellement produit à l'oral et à l'écrit. Retiens trois acquis essentiels. D'abord, un entretien d'embauche suit des \cle{étapes} fixes — salutation, présentation, questions sur les diplômes et l'expérience, questions du candidat, clôture — et chaque étape possède ses expressions propres. Ensuite, les \cle{collocations} (\eng{apply for, meet the requirements, gain experience}) donnent à ton anglais un caractère naturel et soutenu, indispensable au niveau Terminale. Enfin, la vigilance face aux \cle{faux amis} (\eng{actually, formation, passer un entretien}) et aux pièges de prononciation distingue un candidat convaincant d'un candidat hésitant. À l'examen comme dans la vie réelle — à Douala, à Yaoundé ou ailleurs —, ce lexique te servira pour comprendre une annonce, rédiger une lettre de candidature et te présenter avec confiance devant un jury.

\newpage

\coursec{Méthodes et exercices résolus}

\courssub{The stages of a job interview}

\begin{definition}[Les trois phases du processus de recrutement]
Un entretien d'embauche (\eng{job interview}) s'inscrit dans un processus en trois temps. Maîtriser le vocabulaire spécifique à chaque phase permet de suivre une offre d'emploi de sa publication à la signature du contrat.
\end{definition}

\begin{textebox}[Vocabulaire par étapes : Before / During / After]
\begin{center}
\begin{tabularx}{\linewidth}{|X|X|X|}
\hline
\rowcolor{popblueL} \textbf{Before the interview (Avant)} & \textbf{During the interview (Pendant)} & \textbf{After the interview (Après)} \\
\hline
\eng{job advertisement / vacancy} (offre d'emploi) & \eng{interviewer} (recruteur) / \eng{panel} (jury) & \eng{to be hired / recruited} (être embauché) \\
\eng{to apply for a job / a post} (postuler) & \eng{candidate / interviewee} (candidat) & \eng{to be turned down / rejected} (être refusé) \\
\eng{application form} (formulaire de candidature) & \eng{to shake hands} (serrer la main) & \eng{to accept an offer} (accepter une offre) \\
\eng{CV / curriculum vitae} (CV) & \eng{to answer questions} (répondre aux questions) & \eng{probation period} (période d'essai) \\
\eng{covering letter} (GB) / \eng{cover letter} (US) (lettre de motivation) & \eng{strengths and weaknesses} (forces et faiblesses) & \eng{salary} (salaire mensuel/annuel) \\
\eng{shortlist} (liste restreinte) & \eng{to ask about experience} (demander l'expérience) & \eng{wages} (paie horaire/hebdomadaire) \\
\eng{to be shortlisted} (être présélectionné) & \eng{qualifications} (diplômes) & \eng{contract of employment} (contrat de travail) \\
\hline
\end{tabularx}
\end{center}
\end{textebox}

\begin{methode}[Construire son répertoire lexical par étapes]
\begin{enumerate}
\item \textbf{Classe chaque nouveau mot} dans l'une des trois colonnes ci-dessus.
\item \textbf{Apprends les collocations} (associations de mots fixes) : \eng{to apply \textbf{for} a job}, \eng{to fill \textbf{in} a form}, \eng{to attend an interview}, \eng{to gain experience}, \eng{to meet the requirements}, \eng{to offer a contract}.
\item \textbf{Contextualise} : imagine une candidature à la \eng{CDC} (Cameroon Development Corporation) à Buea, ou dans une banque à Douala, ou une ONG à Yaoundé.
\item \textbf{Note la prononciation} en lettres françaises : \eng{interview} « in-teur-viou », \eng{CV} « si-vi », \eng{shortlist} « chort-list », \eng{salary} « sa-leu-ri », \eng{probation} « pro-bé-cheune ».
\end{enumerate}
\end{methode}

\begin{exemplebox}[Exemples traduits]
\eng{She \textbf{applied for} the post of accountant.} — Elle a postulé au poste de comptable.\par
\eng{He was \textbf{shortlisted} among fifty candidates.} — Il a été présélectionné parmi cinquante candidats.\par
\eng{The \textbf{panel} asked tricky questions.} — Le jury a posé des questions pièges.\par
\eng{After a three-month \textbf{probation period}, he got a permanent contract.} — Après une période d'essai de trois mois, il a obtenu un CDI.
\end{exemplebox}

\courssub{People and documents}

\begin{definition}[Acteurs et documents clés]
L'entretien oppose deux parties : celui qui recrute (\eng{employer}, \eng{recruiter}, \eng{interviewer}) et celui qui postule (\eng{candidate}, \eng{applicant}, \eng{interviewee}). Les documents font foi : \eng{CV}, \eng{covering letter}, \eng{application form}, \eng{certificates}, \eng{references}.
\end{definition}

\begin{propriete}[Suffixes de personnes : \textbf{-er / -or} vs \textbf{-ee}]
\begin{center}
\begin{tabular}{|l|l|l|}
\hline
\rowcolor{popblueL} \textbf{Verbe} & \textbf{Agent (celui qui agit) : -er / -or} & \textbf{Patient (celui qui subit) : -ee} \\
\hline
\eng{employ} & \eng{employer} (employeur) & \eng{employee} (employé) \\
\eng{interview} & \eng{interviewer} (recruteur) & \eng{interviewee} (candidat interviewé) \\
\eng{recruit} & \eng{recruiter} (chasseur de têtes) & \eng{recruit} (recrue) \\
\eng{train} & \eng{trainer} (formateur) & \eng{trainee} (stagiaire) \\
\eng{apply} & \eng{applicant} (candidat) & — \\
\hline
\end{tabular}
\end{center}
\emph{Attention :} \eng{applicant} se forme sur \eng{apply} + \emph{-ant} (participe présent nominalisé), pas \eng{*applyee}.
\end{propriete}

\begin{textebox}[Documents de candidature : lexique précis]
\begin{itemize}
\item \eng{CV} (Curriculum Vitae) : utilisé au Royaume-Uni et au Cameroun ; aux USA on dit souvent \eng{resume} (pour un résumé d'une page).
\item \eng{Covering letter} (GB) = \eng{Cover letter} (US) : lettre de motivation. \textbf{Jamais} \eng{motivation letter} (calque).
\item \eng{Application form} : formulaire officiel à remplir (\eng{to fill in / out a form}).
\item \eng{References / Referees} : personnes qui donnent une recommandation (\eng{to provide references}).
\item \eng{Certificates / Diplomas} : originaux ou copies certifiées (\eng{certified copies}).
\end{itemize}
\end{textebox}

\begin{exoresolu}[Identifier les personnes et documents]
Énoncé — Match each definition to the correct word: \emph{interviewee, employer, covering letter, referee, applicant, employee}.\par
1. The person who offers a job. \quad 2. The person who answers the questions during an interview. \quad 3. A document sent with the CV to explain motivation. \quad 4. A person who writes a recommendation. \quad 5. Someone who has a job contract. \quad 6. Someone who sends a form to get a job.
\tcblower
Solution détaillée :
\begin{enumerate}
\item \textbf{employer} (celui qui embauche, agent de \eng{employ}).
\item \textbf{interviewee} (celui qui subit l'interview, suffixe \emph{-ee}).
\item \textbf{covering letter} (lettre de motivation accompagnant le CV).
\item \textbf{referee} (personne donnant une référence ; \eng{reference} = la lettre elle-même).
\item \textbf{employee} (salarié, suffixe \emph{-ee}).
\item \textbf{applicant} (candidat ayant postulé, de \eng{apply} + \emph{-ant}).
\end{enumerate}
\end{exoresolu}

\courssub{Collocations and key expressions}

\begin{propriete}[Collocations incontournables (mots qui vont ensemble)]
Apprendre un mot isolé ne suffit pas ; l'anglais fonctionne par \textbf{blocs lexicaux} (collocations). Une erreur de collocation (ex. \eng{*make an interview}) trahit l'influence du français.
\begin{center}
\begin{tabularx}{\linewidth}{|X|X|X|}
\hline
\rowcolor{popblueL} \textbf{Action} & \textbf{Collocation correcte (GB)} & \textbf{Traduction / Piège à éviter} \\
\hline
Postuler & \eng{apply \textbf{for} a job / a post} & \eng{*apply a job} (préposition manquante) \\
Remplir un formulaire & \eng{fill \textbf{in} a form} (GB) / \eng{fill \textbf{out} a form} (US) & \eng{*fill a form} \\
Passer un entretien & \eng{attend an interview} / \eng{have an interview} & \eng{*pass an interview} (\eng{pass} = réussir un examen) \\
Acquérir de l'expérience & \eng{gain experience} / \eng{acquire experience} & \eng{*win experience} \\
Répondre aux exigences & \eng{meet the requirements} / \eng{match the profile} & \eng{*respond the requirements} \\
Être embauché & \eng{be hired} / \eng{be recruited} / \eng{get the job} & \eng{*be engaged} (\eng{engaged} = fiancé) \\
Négocier le salaire & \eng{negotiate a salary} / \eng{discuss remuneration} & — \\
\hline
\end{tabularx}
\end{center}
\end{propriete}

\begin{textebox}[Expressions idiomatiques et formules de politesse]
\begin{itemize}
\item \eng{“Please have a seat.”} — « Asseyez-vous, je vous en prie. » (Formule standard d'accueil).
\item \eng{“Could you tell us about yourself?”} — Question d'ouverture classique (présentation en 1-2 min).
\item \eng{“What are your main strengths and weaknesses?”} — Forces / faiblesses. Réponse type : \eng{“I am a hard-working and reliable person. I sometimes struggle with public speaking, but I am taking a course to improve.”}
\item \eng{“Why should we hire you?”} — « Pourquoi vous et pas un autre ? » (Argumentaire : compétences + motivation).
\item \eng{“Do you have any questions for us?”} — Toujours en avoir une prête (ex. \eng{“What does a typical day look like?”}).
\item \eng{“We shall contact you shortly / by the end of the week.”} — Conclusion standard.
\item \eng{“I look forward to hearing from you.”} — Formule de politesse finale (écrit/oral) ; \eng{look forward to} + \textbf{-ing}.
\end{itemize}
\end{textebox}

\begin{methode}[Réemployer les collocations dans un dialogue type]
\begin{enumerate}
\item \textbf{Salutations} : \eng{Good morning. Please have a seat. / Thank you.}
\item \textbf{Présentation} : \eng{Could you tell us about yourself? / I hold a BTS in Accounting and have two years' experience at \ldots}
\item \textbf{Compétences} : \eng{What are your strengths? / I am punctual, honest and I work well under pressure.}
\item \textbf{Motivation} : \eng{Why do you want this job? / I admire your company's reputation and I want to contribute to its growth.}
\item \textbf{Conclusion} : \eng{Thank you for coming. We shall contact you next week. / Thank you. I look forward to hearing from you.}
\end{enumerate}
\end{methode}

\begin{exoresolu}[Transformer un récit en dialogue (style direct)]
Énoncé — Turn into a dialogue (6+ exchanges).\par
\eng{Mr Etoundi welcomed Miss Bih and asked her to sit down. He wanted to know why she had applied for the job of cashier. She explained that she had a diploma in accounting and two years' experience. He asked about her strengths. She said she was honest and punctual. He thanked her and promised to contact her soon.}
\tcblower
Solution détaillée — Passage du style indirect (récit au prétérit) au style direct (dialogue au présent). La question indirecte \eng{He wanted to know why she had applied} redevient \eng{Why did you apply?} (auxiliaire \eng{did} + sujet + base verbale).
\par
Réponse rédigée :
\eng{\textbf{Interviewer:} Good morning, Miss Bih. Please have a seat.\par
\textbf{Candidate:} Good morning, Sir. Thank you.\par
\textbf{Interviewer:} Why did you apply for the job of cashier?\par
\textbf{Candidate:} Because I have a diploma in accounting and two years' experience in a microfinance company in Yaoundé.\par
\textbf{Interviewer:} Very good. What are your main strengths?\par
\textbf{Candidate:} I am honest and punctual, and I work well under pressure.\par
\textbf{Interviewer:} Thank you very much, Miss Bih. We shall contact you soon.\par
\textbf{Candidate:} Thank you, Sir. I look forward to hearing from you.}
\end{exoresolu}

\courssub{False friends and French traps}

\begin{attention}[Les 7 faux amis et calques « tueurs » du thème]
\begin{enumerate}
\item \textbf{\eng{Formation} $\neq$ Formation professionnelle}. \eng{Formation} = formation géologique, mise en place. \rightarrow \eng{Training}, \eng{a training course}, \eng{vocational training}.
\item \textbf{\eng{Actually} $\neq$ Actuellement}. \eng{Actually} = « en fait, en réalité ». \rightarrow \eng{Currently}, \eng{at present}, \eng{nowadays}.
\item \textbf{\eng{To assist} $\neq$ Assister à}. \eng{To assist} = aider. \rightarrow \eng{To attend} (une réunion, un entretien), \eng{to be present at}.
\item \textbf{\eng{To pass an interview} $\neq$ Passer un entretien}. \eng{To pass} = réussir (un examen). \rightarrow \eng{To attend / have an interview}. On dit \eng{She passed the interview} seulement si elle a \textbf{réussi}.
\item \textbf{\eng{Engaged} $\neq$ Embauché}. \eng{Engaged} = fiancé. \rightarrow \eng{Hired}, \eng{recruited}, \eng{appointed}.
\item \textbf{\eng{Resume} (US) / \eng{CV} (GB) $\neq$ Résumé}. \eng{Resume} (verbe) = reprendre ; (nom US) = CV. \eng{Summary} = résumé.
\item \textbf{\eng{Stage} $\neq$ Stage (période en entreprise)}. \eng{Stage} = scène, étape. \rightarrow \eng{Internship} (stage étudiant), \eng{work placement} (GB), \eng{training period}.
\end{enumerate}
\end{attention}

\begin{propriete}[Calques structurels à éradiquer]
\begin{center}
\begin{tabularx}{\linewidth}{|X|X|X|}
\hline
\rowcolor{poporangeL} \textbf{Français (calque)} & \textbf{Anglais correct} & \textbf{Règle} \\
\hline
Je suis d'accord & \eng{I agree} & \eng{Agree} est un verbe, pas un adjectif. \\
J'ai 20 ans & \eng{I am twenty years old} & \eng{Be} + âge + \eng{years old}. \\
Il attend depuis 10h & \eng{He has been waiting since 10 o'clock} & \eng{Since} + action continue $\rightarrow$ Present Perfect Continuous. \\
Elle a passé un entretien & \eng{She attended / had an interview} & Voir faux ami \eng{pass}. \\
Il a suivi une formation & \eng{He took a training course} / \eng{He did a training course} & \eng{Follow a training} $\neq$ suivre une formation. \\
\hline
\end{tabularx}
\end{center}
\end{propriete}

\begin{exoresolu}[Correction ciblée : faux amis, prépositions, temps]
Énoncé — Correct the mistakes.\par
\eng{1. I am agree with the terms. \quad 2. She passed a job interview yesterday. \quad 3. He followed a formation in marketing. \quad 4. I am actually working in Douala. \quad 5. They wait the manager since 9 am. \quad 6. The employee was engaged last month. \quad 7. She is good in English.}
\tcblower
Solution détaillée :
\begin{enumerate}
\item \eng{I \textbf{agree} with the terms.} (Verbe \eng{agree}, pas \eng{be agree}).
\item \eng{She \textbf{attended} a job interview yesterday.} (ou \eng{had}). \eng{Passed} implique la réussite.
\item \eng{He took a \textbf{training course} in marketing.} (\eng{Formation} = faux ami ; \eng{follow} $\neq$ suivre un cours).
\item \eng{I am \textbf{currently} working in Douala.} (\eng{Actually} = en fait ; \eng{currently} = actuellement).
\item \eng{They \textbf{have been waiting for} the manager since 9 am.} (\eng{Wait for} + Present Perfect Continuous avec \eng{since}).
\item \eng{The employee was \textbf{hired} last month.} (\eng{Engaged} = fiancé).
\item \eng{She is good \textbf{at} English.} (\eng{Good at} + nom/-ing, pas \eng{in}).
\end{enumerate}
\end{exoresolu}

\courssub{Word families, pronunciation and spelling}

\begin{propriete}[Construction des familles de mots (dérivation)]
\begin{center}
\begin{tabular}{|l|l|l|l|l|}
\hline
\rowcolor{popblueL} \textbf{Verbe} & \textbf{Nom (action/état)} & \textbf{Nom (agent -er/-or)} & \textbf{Nom (patient -ee)} & \textbf{Adjectif} \\
\hline
\eng{apply} & \eng{application} & \eng{applicant} & — & \eng{applicable} \\
\eng{employ} & \eng{employment} & \eng{employer} & \eng{employee} & \eng{employable} \\
\eng{interview} & \eng{interview} & \eng{interviewer} & \eng{interviewee} & — \\
\eng{recruit} & \eng{recruitment} & \eng{recruiter} & \eng{recruit} & — \\
\eng{qualify} & \eng{qualification} & — & — & \eng{qualified} \\
\eng{appoint} & \eng{appointment} & — & \eng{appointee} & — \\
\eng{train} & \eng{training} & \eng{trainer} & \eng{trainee} & \eng{trained} \\
\hline
\end{tabular}
\end{center}
\emph{Remarques :} \eng{Unemployment} (chômage) = \emph{un-} + \eng{employment}. \eng{Punctual} $\rightarrow$ \eng{punctuality} (exception : \emph{-ity} au lieu de \emph{-ness}). \eng{Qualification} souvent au pluriel (\eng{qualifications}) pour « diplômes et titres ».
\end{propriete}

\begin{methode}[Règles d'orthographe : verbes en \emph{-y} et doublement]
\begin{enumerate}
\item \textbf{Verbes en \emph{-y} précédé d'une consonne} : le \emph{y} devient \emph{i} devant \emph{-ed} / \emph{-s} : \eng{apply} $\rightarrow$ \eng{applied}, \eng{applies} ; \eng{try} $\rightarrow$ \eng{tried}. Mais \emph{-ing} garde le \emph{y} : \eng{applying}.
\item \textbf{Verbes en \emph{-y} précédé d'une voyelle} : le \emph{y} ne change jamais : \eng{employ} $\rightarrow$ \eng{employed}, \eng{employing} ; \eng{enjoy} $\rightarrow$ \eng{enjoyed}.
\item \textbf{Doublement de la consonne finale} (accent sur dernière syllabe, structure CVC) : \eng{refer} $\rightarrow$ \eng{referred}, \eng{referring} ; \eng{occur} $\rightarrow$ \eng{occurred}. Pas de doublement si accent sur 1ère syllabe : \eng{offer} $\rightarrow$ \eng{offered}, \eng{interview} $\rightarrow$ \eng{interviewed}.
\end{enumerate}
\end{methode}

\begin{textebox}[Guide de prononciation (notation française approximative)]
\begin{itemize}
\item \eng{Applicant} : « a-pli-keunt » (accent sur 1ère syllabe).
\item \eng{Qualification} : « kwo-li-fi-ké-cheune » (son \eng{ti} = /ʃ/ « che »).
\item \eng{Experience} : « iks-pi-rieunce » (le \emph{x} = /ks/, \emph{ie} = /i:/).
\item \eng{Recruit} : « ri-kroute » (\eng{ui} = /u:/).
\item \eng{Interviewee} : « in-teur-viou-i » (accent final sur \emph{-ee}).
\item \eng{Salary} : « sa-leu-ri » (3 syllabes, pas « sa-la-ri »).
\item \eng{CV} : « si-vi » (épeler les lettres).
\item \eng{Shortlisted} : « chort-lis-tide » (\emph{ed} = /ɪd/ après \emph{t/d}).
\end{itemize}
\end{textebox}

\begin{exoresolu}[Mots de la même famille : forme correcte]
Énoncé — Complete with the right form of the word in brackets.\par
\eng{1. The \ldots\ (employ) interviewed ten \ldots\ (apply). \quad 2. Her \ldots\ (qualify) are impressive. \quad 3. He received a letter of \ldots\ (appoint). \quad 4. \ldots\ (employ) rises among youth. \quad 5. The \ldots\ (interview) was nervous.}
\tcblower
Solution détaillée :
\begin{enumerate}
\item \eng{The \textbf{employer} interviewed ten \textbf{applicants}.} (Agent \emph{-er} ; candidat = \eng{applicant}).
\item \eng{Her \textbf{qualifications} are impressive.} (Pluriel usuel pour « diplômes »).
\item \eng{a letter of \textbf{appointment}} (\eng{appoint} + \emph{-ment} = nomination/embauche).
\item \eng{\textbf{Unemployment} rises among youth.} (Négation \emph{un-} + \eng{employment} = chômage).
\item \eng{The \textbf{interviewee} was nervous.} (Patient \emph{-ee} = celui qui passe l'entretien).
\end{enumerate}
\end{exoresolu}

\courssub{Vocabulary in context: a model dialogue}

\begin{textebox}[Dialogue modèle : Entretien pour un poste de Commercial (Douala)]
\eng{\textbf{Interviewer:} Good morning. You are Mr Fonkeng, aren't you? Please have a seat.\par
\textbf{Candidate:} Good morning. Yes, that's right. Thank you.\par
\textbf{Interviewer:} First of all, could you tell us why you applied for this position?\par
\textbf{Candidate:} Certainly. I saw your advertisement in \textit{The Cameroon Tribune}. I have a Higher Professional Diploma in Marketing and three years' experience as a sales representative for a brewery in Douala. I know the local market well.\par
\textbf{Interviewer:} That sounds promising. What would you say are your main strengths?\par
\textbf{Candidate:} I am dynamic, persistent and I have excellent communication skills in both English and French. I also have a clean driving licence.\par
\textbf{Interviewer:} And your weaknesses?\par
\textbf{Candidate:} I sometimes find it hard to delegate, but I am learning to trust my team more.\par
\textbf{Interviewer:} Good. What are your salary expectations?\par
\textbf{Candidate:} I am looking for a salary in the range of 250,000 to 300,000 FCFA per month, plus commissions, in line with the market.\par
\textbf{Interviewer:} We shall discuss that later. Do you have any questions for us?\par
\textbf{Candidate:} Yes. What does the probation period involve?\par
\textbf{Interviewer:} It is a three-month period with monthly evaluations. We shall contact you by Friday.\par
\textbf{Candidate:} Thank you for your time. I look forward to hearing from you.\par
\textbf{Interviewer:} Thank you, Mr Fonkeng. Goodbye.}
\end{textebox}

\begin{experience}[Jeu de rôle : Simulation d'entretien (binômes)]
\begin{enumerate}
\item \textbf{Rôles} : A = Recruteur (Directeur RH d'une entreprise de cacao à Douala / ONG à Yaoundé / Hôtel à Kribi). B = Candidat (BTS / Licence / Master, avec ou sans expérience).
\item \textbf{Préparation (5 min)} : A prépare 5 questions types ; B prépare son « pitch » (présentation 1 min) et 3 forces/faiblesses.
\item \textbf{Déroulement (5 min)} : Salutations $\rightarrow$ Questions/Réponses $\rightarrow$ Salaire/Période d'essai $\rightarrow$ Conclusion.
\item \textbf{Contraintes linguistiques} : Utiliser au moins 5 collocations de la leçon (\eng{apply for}, \eng{attend an interview}, \eng{strengths/weaknesses}, \eng{salary expectations}, \eng{probation period}, \eng{look forward to}). Interdiction des faux amis (\eng{formation}, \eng{actually}, \eng{pass an interview}, \eng{I am agree}).
\item \textbf{Retour} : La classe note le vocabulaire utilisé et les erreurs entendues. Correction collective.
\end{enumerate}
\end{experience}

\begin{exercice}[Entraînement complet : Vocabulaire de l'entretien]
\textbf{A. Complétez avec le mot juste (choix multiple).}
1. She sent her CV and a \ldots\ letter. \quad a) motivation \quad b) covering \quad c) presentation \quad d) application\par
2. He was \ldots\ for the final interview with three other candidates. \quad a) listed \quad b) shortlisted \quad c) outlined \quad d) enrolled\par
3. The \ldots\ asked difficult questions about my previous job. \quad a) interviewer \quad b) interviewee \quad c) employer \quad d) employee\par
4. I \ldots\ with your analysis of the market. \quad a) am agree \quad b) agree \quad c) am agreed \quad d) have agree\par
5. He has been \ldots\ for a reply since Monday. \quad a) waiting \quad b) waiting for \quad c) waited \quad d) waited for\par
\textbf{B. Mettez le mot entre parenthèses à la forme correcte.}
6. The \ldots\ (recruit) process took two months.\par
7. Youth \ldots\ (employ) is a national challenge.\par
8. She is a very \ldots\ (qualify) engineer.\par
9. We need your \ldots\ (apply) form by Friday.\par
10. The \ldots\ (interview) answered confidently.\par
\textbf{C. Traduisez en anglais (attention aux pièges).}
11. Il a suivi une formation en informatique l'an dernier.\par
12. Actuellement, je cherche du travail.\par
13. Elle a passé un entretien hier matin.\par
14. Nous attendons le directeur depuis une heure.\par
15. Je suis d'accord pour signer le contrat.
\end{exercice}


\begin{corrige}[Exercice Entraînement complet]
\textbf{A. Choix multiple}
1. \textbf{b) covering} (\eng{covering letter} = lettre de motivation ; \eng{motivation letter} est un calque).\par
2. \textbf{b) shortlisted} (terme technique : liste restreinte).\par
3. \textbf{a) interviewer} (celui qui pose les questions = agent \emph{-er} ; le candidat est l'\eng{interviewee}).\par
4. \textbf{b) agree} (verbe, pas \eng{be agree}).\par
5. \textbf{b) waiting for} (\eng{wait} est intransitif, il faut \eng{for} ; présent perfect continuous \eng{has been waiting for} car \eng{since Monday}).
\textbf{B. Dérivation}
6. \textbf{recruitment} (nom d'action).\par
7. \textbf{unemployment} (chômage = négation de l'emploi).\par
8. \textbf{qualified} (adjectif : \eng{qualified engineer}).\par
9. \textbf{application} (nom : \eng{application form}).\par
10. \textbf{interviewee} (patient \emph{-ee} = celui qui subit l'interview).
\textbf{C. Traduction (pièges)}
11. \eng{He took a training course in computer science last year.} (\eng{formation} $\rightarrow$ \eng{training} ; \eng{follow} $\rightarrow$ \eng{take/do}).\par
12. \eng{I am currently looking for a job.} (\eng{actuellement} $\rightarrow$ \eng{currently} ; \eng{look for} = chercher).\par
13. \eng{She attended a job interview yesterday morning.} (\eng{passer un entretien} $\rightarrow$ \eng{attend/have} ; pas \eng{pass}).\par
14. \eng{We have been waiting for the manager for an hour.} (\eng{since/for} + durée $\rightarrow$ Present Perfect Continuous ; \eng{wait for}).\par
15. \eng{I agree to sign the contract.} (\eng{agree} verbe ; \eng{agree to} + infinitif).
\end{corrige}

\begin{savaistu}[Le marché de l'emploi au Cameroun : contexte bilingue]
Le Cameroun est officiellement bilingue (français/anglais). Dans les entreprises anglophones (ex. \eng{CDC}, \eng{PAMOL}, \eng{SONARA} à ses débuts, hôtels de la côte, ONG internationales), les entretiens se déroulent souvent en anglais. Le \eng{CV} suit le format britannique (2 pages max, pas de photo, pas d'âge, pas d'état civil détaillé). La \eng{covering letter} est cruciale : elle doit être adressée à une personne nommée (\eng{Dear Mr/Ms [Name]}, \eng{Yours sincerely}) ou \eng{Dear Sir/Madam}, \eng{Yours faithfully}. Les salaires sont négociés en FCFA (souvent net/mois). La \eng{probation period} légale est de 3 mois (renouvelable une fois) selon le Code du Travail. Maîtriser ce lexique, c'est un atout professionnel réel pour un élève de Terminale A.
\end{savaistu}

\newpage

\coursec{Exercices}

\courssub{Je vérifie mes connaissances}

\begin{exercice}[QCM : la bonne collocation]
Pour chaque phrase, choisis la bonne réponse (a, b ou c).
\begin{enumerate}
\item To \cle{\dots} an interview means to be present at it. \\
a) assist \quad b) attend \quad c) assist to
\item The company has \cle{\dots} five candidates for the post. \\
a) shortlisted \quad b) shortlistened \quad c) shortly listed
\item She sent her \cle{\dots} letter along with her CV. \\
a) covering \quad b) cover \quad c) recovering
\item During the interview, the \cle{\dots} asked him many questions. \\
a) pannel \quad b) panel \quad c) panal
\item You should arrive \cle{\dots} time for your interview. \\
a) in \quad b) at \quad c) on
\end{enumerate}
\end{exercice}

\begin{exercice}[Vrai ou faux ?]
Lis le texte suivant, puis dis si les affirmations sont \emph{true} ou \emph{false}. Justifie chaque réponse par une ligne du texte.
\begin{textebox}[A young graduate's first interview]
Marius, a young graduate from the University of Yaoundé I, was invited for an interview at a bank in Douala. He arrived twenty minutes early and greeted the panel politely. When he was asked about his strengths, he spoke confidently about his computer skills. However, he had forgotten to bring copies of his certificates, which made a bad impression on the interviewers. A week later, he received a phone call: he had not been selected, but the manager encouraged him to apply again next year.
\end{textebox}
\begin{enumerate}
\item Marius was late for his interview.
\item He answered the questions with confidence.
\item He brought all his documents with him.
\item The bank offered him the job immediately.
\end{enumerate}
\end{exercice}

\begin{exercice}[Vocabulaire : trouve le mot]
Complète chaque définition par le mot anglais correspondant, choisi dans la leçon.
\begin{enumerate}
\item A document that summarises your education and work experience: \dots
\item The person who applies for a job: a \dots
\item The people who ask questions during an interview: the \dots
\item A meeting in which an employer asks questions to choose the best candidate: a job \dots
\item Qualities you are good at: your \dots
\end{enumerate}
\end{exercice}

\begin{exercice}[Conjugaison et formes à compléter]
Complète les phrases avec la forme correcte du verbe entre parenthèses (prétérit ou participe passé).
\begin{enumerate}
\item She \dots (send) her application letter two weeks ago.
\item The candidates \dots (be) interviewed by a panel of three managers.
\item He \dots (speak) about his experience during the interview.
\item They \dots (choose) the best applicant yesterday.
\item I \dots (write) my CV before applying for the job.
\end{enumerate}
\end{exercice}

\courssub{Je m'entraîne}

\begin{exercice}[Traduction]
Traduis les phrases suivantes en anglais, en utilisant le lexique de la leçon.
\begin{enumerate}
\item Elle a été présélectionnée pour un poste dans une banque à Douala.
\item Il faut envoyer une lettre de motivation avec ton CV.
\item Le jury lui a posé des questions sur ses qualifications.
\item Il a fait bonne impression pendant l'entretien.
\item Arrive à l'heure et habille-toi correctement.
\end{enumerate}
\end{exercice}

\begin{exercice}[Corrige les faux amis]
Chaque phrase contient un calque du français. Corrige-la.
\begin{enumerate}
\item He \emph{passed} the interview and got the job.
\item She \emph{made a formation} in computer science.
\item I want to \emph{postulate} for this position.
\item The \emph{formation} of the candidates was impressive.
\item He is very \emph{actually} in his way of dressing.
\end{enumerate}
\end{exercice}

\begin{exercice}[Familles de mots]
Complète le tableau avec les mots de la même famille.
\begin{center}
\begin{tabularx}{\linewidth}{|X|X|X|X|}
\hline
\rowcolor{popblueL} Verbe & Personne & Chose / notion & Adjectif \\
\hline
to apply & \dots & \dots & \dots \\
\hline
to employ & \dots & \dots & \dots \\
\hline
to qualify & \dots & \dots & \dots \\
\hline
to interview & \dots & \dots & \dots \\
\hline
\end{tabularx}
\end{center}
\end{exercice}

\begin{exercice}[Appariement : expressions idiomatiques]
Relie chaque expression anglaise à sa traduction française.
\begin{center}
\begin{tabularx}{\linewidth}{|X|X|}
\hline
\rowcolor{popblueL} English & Français \\
\hline
1. to sell yourself & a. réfléchir vite, improviser \\
\hline
2. to break the ice & b. être présélectionné \\
\hline
3. to think on your feet & c. se mettre en valeur \\
\hline
4. to make a good impression & d. une lettre de motivation \\
\hline
5. to be shortlisted & e. briser la glace \\
\hline
6. a covering letter & f. faire bonne impression \\
\hline
\end{tabularx}
\end{center}
\end{exercice}

\courssub{Je me prépare au Bac}

\begin{exercice}[Compréhension écrite — type Bac]
Lis le texte suivant, puis réponds aux questions en anglais, par des phrases complètes.
\begin{textebox}[An interview in Buea]
Last month, Ngong Delphine, a final-year student from Buea, was shortlisted for the post of secretary in a cocoa export company. The interview was held before a panel of four people. Although she was nervous at first, she managed to break the ice by smiling and greeting everyone warmly. The interviewers asked her about her qualifications, her strengths and her weaknesses. She sold herself well, explaining that she was hardworking, punctual and fluent in both English and French. When she did not understand a question, she politely asked the interviewer to repeat it instead of guessing. Two days later, the manager phoned her: she had got the job. “Your confidence and honesty made a very good impression on us,” he said. Delphine advises all young job seekers to prepare their documents carefully and to practise answering questions with a friend before the big day.
\end{textebox}
\textbf{A. Comprehension questions}
\begin{enumerate}
\item What post did Delphine apply for?
\item How many people interviewed her?
\item How did she break the ice at the beginning of the interview?
\item What did she do when she did not understand a question?
\item Why did the company choose her, according to the manager?
\end{enumerate}
\textbf{B. Language work}
\begin{enumerate}
\item Find in the text the English equivalents of: \emph{présélectionnée}, \emph{qualifications}, \emph{faire bonne impression}.
\item Give the noun corresponding to the verb \emph{to apply} and the adjective corresponding to \emph{confidence}.
\item Put into the passive voice: “The interviewers asked her many questions.”
\end{enumerate}
\end{exercice}

\begin{exercice}[Dialogue à trous — type Bac]
Complète le dialogue d'entretien suivant avec les mots de la liste : \emph{shortlisted — panel — strengths — opportunity — impression — qualifications}.
\begin{textebox}[At the interview]
Interviewer: Good morning, Mr Etoundi. Please sit down. You were \dots\ (1) among more than fifty candidates, so congratulations.\\
Candidate: Thank you very much, Madam. It is a great \dots\ (2) for me.\\
Interviewer: The \dots\ (3) would like to know more about your \dots\ (4). Do you have a degree in accounting?\\
Candidate: Yes, I do. I graduated from the University of Douala last year.\\
Interviewer: Very good. Now, what are your main \dots\ (5)?\\
Candidate: I am hardworking, honest and I can work under pressure.\\
Interviewer: Thank you. You have made a very good \dots\ (6) on us. We shall contact you next week.
\end{textebox}
\end{exercice}

\begin{exercice}[Expression écrite — type Bac]
\textbf{Topic:} Your friend has a job interview next week. Write a letter giving him or her five pieces of advice. Use the vocabulary of the lesson (interview, panel, CV, covering letter, strengths, to make a good impression, punctual, etc.). Write between \textbf{80 and 100 words}. Do not write your real name; sign as “Manga”.
\end{exercice}

\newpage

\coursec{Corrigés des exercices}

\begin{corrige}[Exercice 1 — QCM : la bonne collocation]
\begin{enumerate}
\item \textbf{b) attend}. \eng{To attend an interview} signifie « assister à / se présenter à un entretien ». \eng{Assist} signifie « aider » ; \eng{assist to} n'existe pas (on dirait \eng{to assist in} ou \eng{to help with}).
\item \textbf{a) shortlisted}. Le participe passé de \eng{to shortlist} est \emph{shortlisted} (verbe régulier). \eng{Shortlistened} est une forme inventée par attraction de \eng{listen} ; \eng{shortly listed} est un calque de « présélectionné ».
\item \textbf{a) covering}. La collocation correcte est \eng{a covering letter} (britannique) ou \eng{a cover letter} (américain). \eng{Recovering} signifie « se remettre (d'une maladie) ».
\item \textbf{b) panel}. Orthographe : un seul \emph{n} — \eng{panel} (« jury, commission »). Prononciation : « pa-neul ».
\item \textbf{c) on}. La collocation est \eng{on time} (« à l'heure »). Attention : \eng{in time} signifie « à temps » (assez tôt pour faire quelque chose).
\end{enumerate}
\textbf{Points de vigilance :} mémoriser les collocations comme des blocs : \eng{to attend an interview}, \eng{to be shortlisted}, \eng{a covering letter}, \eng{to arrive on time}.
\end{corrige}

\begin{corrige}[Exercice 2 — Vrai ou faux ?]
\begin{enumerate}
\item \textbf{False.} Justification : \eng{“He arrived twenty minutes early”} — il est arrivé en avance, pas en retard.
\item \textbf{True.} Justification : \eng{“he spoke confidently about his computer skills”} — il a répondu avec assurance.
\item \textbf{False.} Justification : \eng{“he had forgotten to bring copies of his certificates”} — il avait oublié ses diplômes.
\item \textbf{False.} Justification : \eng{“he had not been selected”} — il n'a pas été retenu ; le directeur l'a seulement encouragé à postuler de nouveau.
\end{enumerate}
\textbf{Méthode :} toujours citer la ligne exacte du texte qui justifie la réponse ; ne jamais se contenter de déduire ou d'interpréter.
\end{corrige}

\begin{corrige}[Exercice 3 — Vocabulaire : trouve le mot]
\begin{enumerate}
\item a \cle{CV} (curriculum vitae) / a \cle{r\'esum\'e} (américain) — « un CV ».
\item an \cle{applicant} (ou a \cle{candidate}) — « un candidat ».
\item the \cle{interviewers} (ou the \cle{panel}) — « les membres du jury ».
\item a job \cle{interview} — « un entretien d'embauche ».
\item your \cle{strengths} — « tes points forts » (contraire : \eng{weaknesses}, « points faibles »).
\end{enumerate}
\textbf{Point de vigilance :} \eng{strengths} se prononce « strenkths » et s'emploie presque toujours au pluriel dans ce contexte.
\end{corrige}

\begin{corrige}[Exercice 4 — Conjugaison et formes à compléter]
\begin{enumerate}
\item She \textbf{sent} her application letter two weeks ago. (prétérit irrégulier de \eng{send} : send — sent — sent ; le marqueur \eng{two weeks ago} impose le prétérit.)
\item The candidates \textbf{were interviewed} by a panel of three managers. (voix passive au prétérit : were + participe passé ; \eng{interview} est régulier.)
\item He \textbf{spoke} about his experience during the interview. (speak — spoke — spoken.)
\item They \textbf{chose} the best applicant yesterday. (choose — chose — chosen.)
\item I \textbf{wrote} my CV before applying for the job. (write — wrote — written.)
\end{enumerate}
\textbf{Points de vigilance :} ne jamais écrire \eng{speaked}, \eng{choosed} ou \eng{writed} : ces verbes sont irréguliers. Après \eng{before}, on emploie la forme en \emph{-ing} : \eng{before applying}.
\end{corrige}

\begin{corrige}[Exercice 5 — Traduction]
\begin{enumerate}
\item \eng{She was shortlisted for a post in a bank in Douala.} (voix passive : was + participe passé ; \eng{post} ou \eng{position}, jamais \eng{place}.)
\item \eng{You must send a covering letter with your CV.} (ou \eng{You should send\ldots} ; « lettre de motivation » = \eng{covering letter}, jamais \eng{motivation letter} en anglais britannique.)
\item \eng{The panel asked him questions about his qualifications.} (\eng{to ask somebody questions} ; pas de préposition après \eng{ask} + personne.)
\item \eng{He made a good impression during the interview.} (collocation : \eng{to make a good impression}, et non \eng{to do}.)
\item \eng{Arrive on time and dress properly.} (impératif ; \eng{on time} = « à l'heure » ; \eng{properly} = « correctement ».)
\end{enumerate}
\end{corrige}

\begin{corrige}[Exercice 6 — Corrige les faux amis]
\begin{enumerate}
\item \eng{He \textbf{took} the interview and got the job.} (ou \eng{He was successful at the interview}.) En anglais, \eng{to pass} s'emploie pour un \emph{examen} (\eng{to pass an exam}), pas pour un entretien.
\item \eng{She \textbf{did a training course} in computer science.} (ou \eng{She was trained in computer science}.) \eng{Formation} au sens de « formation professionnelle » se traduit par \eng{training}.
\item \eng{I want to \textbf{apply} for this position.} \eng{Postulate} existe mais signifie « postuler » au sens philosophique (poser un principe) ; « postuler à un emploi » = \eng{to apply for a job}.
\item \eng{The \textbf{training} of the candidates was impressive.} (ou \eng{The candidates' qualifications were impressive.}) \eng{Formation} en anglais signifie « formation » au sens de création/forme.
\item \eng{He is very \textbf{up to date} in his way of dressing.} (ou \eng{He dresses very fashionably}.) \eng{Actually} signifie « en fait, réellement », jamais « actuel » ; « actuel » = \eng{current, present}.
\end{enumerate}
\textbf{Point de vigilance :} ces calques (\eng{pass an interview}, \eng{formation}, \eng{postulate}, \eng{actually}) sont parmi les erreurs les plus fréquentes des candidats francophones au Bac.
\end{corrige}

\begin{corrige}[Exercice 7 — Familles de mots]
\begin{center}
\begin{tabularx}{\linewidth}{|X|X|X|X|}
\hline
\rowcolor{popblueL} Verbe & Personne & Chose / notion & Adjectif \\
\hline
to apply & applicant & application & applicable \\
\hline
to employ & employer / employee & employment & employed / employable \\
\hline
to qualify & qualifier & qualification & qualified \\
\hline
to interview & interviewer / interviewee & interview & interviewed \\
\hline
\end{tabularx}
\end{center}
\textbf{Points de vigilance :}
\begin{itemize}
\item \eng{Employer} = « l'employeur » (celui qui embauche) ; \eng{employee} = « l'employé » (celui qui est embauché). Le suffixe \emph{-ee} désigne la personne qui subit l'action.
\item \eng{Interviewer} pose les questions ; \eng{interviewee} y répond.
\item Prononciation : \eng{applicant} « a-pli-kente », \eng{qualification} « kouo-li-fi-ké-cheune ».
\end{itemize}
\end{corrige}

\begin{corrige}[Exercice 8 — Appariement : expressions idiomatiques]
\begin{center}
\begin{tabularx}{\linewidth}{|X|X|}
\hline
\rowcolor{popblueL} English & Français \\
\hline
1. to sell yourself & c. se mettre en valeur \\
\hline
2. to break the ice & e. briser la glace \\
\hline
3. to think on your feet & a. réfléchir vite, improviser \\
\hline
4. to make a good impression & f. faire bonne impression \\
\hline
5. to be shortlisted & b. être présélectionné \\
\hline
6. a covering letter & d. une lettre de motivation \\
\hline
\end{tabularx}
\end{center}
\textbf{Exemple d'emploi :} \eng{During the interview, she had to think on her feet.} « Pendant l'entretien, elle a dû réfléchir vite. »
\end{corrige}

\begin{corrige}[Exercice 9 — Compréhension écrite — type Bac]
\textbf{Barème indicatif :} A. 5 points (1 pt par réponse correcte et complète) ; B. 3 points (1 pt par item). Total : 8 points.

\textbf{A. Comprehension questions — réponses modèles :}
\begin{enumerate}
\item \eng{She applied for the post of secretary in a cocoa export company.}
\item \eng{She was interviewed by a panel of four people.}
\item \eng{She broke the ice by smiling and greeting everyone warmly.}
\item \eng{When she did not understand a question, she politely asked the interviewer to repeat it.}
\item \eng{According to the manager, the company chose her because her confidence and honesty made a very good impression on them.}
\end{enumerate}

\textbf{B. Language work :}
\begin{enumerate}
\item \emph{présélectionnée} = \eng{shortlisted} ; \emph{qualifications} = \eng{qualifications} ; \emph{faire bonne impression} = \eng{to make a good impression}.
\item \eng{to apply} $\rightarrow$ le nom : \eng{application} (ou \eng{applicant} pour la personne) ; \eng{confidence} $\rightarrow$ l'adjectif : \eng{confident}.
\item Voix passive : \eng{She was asked many questions by the interviewers.} (ou \eng{Many questions were asked to her by the interviewers}, moins élégant). Prétérit de \eng{ask} : \emph{asked} ; auxiliaire \eng{was} + participe passé.
\end{enumerate}

\textbf{Points de vigilance :} répondre par des phrases complètes qui reprennent les termes de la question ; à la voix passive, le complément d'agent est introduit par \eng{by} ; ne pas confondre \eng{confident} (adjectif) et \eng{confidence} (nom).
\end{corrige}

\begin{corrige}[Exercice 10 — Dialogue à trous — type Bac]
\textbf{Barème indicatif :} 6 points (1 pt par mot correctement placé).
\begin{enumerate}
\item \textbf{shortlisted} — \eng{You were shortlisted among more than fifty candidates.} (« Tu as été présélectionné parmi plus de cinquante candidats. »)
\item \textbf{opportunity} — \eng{It is a great opportunity for me.} (« C'est une grande opportunité pour moi. »)
\item \textbf{panel} — \eng{The panel would like to know more\ldots} (« Le jury aimerait en savoir davantage\ldots »)
\item \textbf{qualifications} — \eng{\ldots about your qualifications.} (« \ldots sur tes qualifications. »)
\item \textbf{strengths} — \eng{What are your main strengths?} (« Quels sont tes principaux points forts ? »)
\item \textbf{impression} — \eng{You have made a very good impression on us.} (« Tu as fait une très bonne impression sur nous. »)
\end{enumerate}
\textbf{Point de vigilance :} vérifier la cohérence grammaticale : \eng{were} + participe passé (1), article indéfini \eng{a great} + nom singulier (2), \eng{your} + nom pluriel attendu (5).
\end{corrige}

\begin{corrige}[Exercice 11 — Expression écrite — type Bac]
\textbf{Barème indicatif (sur 10) :} respect de la tâche et de la longueur (80–100 mots) : 2 pts ; richesse et exactitude du lexique de la leçon : 3 pts ; correction grammaticale : 3 pts ; organisation et formes de la lettre (salutation, formule finale, signature « Manga ») : 2 pts.

\textbf{Proposition de rédaction modèle :}
\begin{textebox}[Model letter]
Dear Alain,\\
I heard that you have a job interview next week. Congratulations on being shortlisted! Here is my advice. First, prepare your CV, your covering letter and copies of your certificates. Second, learn about the company before the big day. Third, arrive on time and dress properly to make a good impression on the panel. Fourth, speak clearly about your strengths and sell yourself. Finally, if you do not understand a question, ask the interviewer to repeat it. Stay calm and confident. Good luck!\\
Yours,\\
Manga
\end{textebox}
(96 mots.)

\textbf{Points de vigilance :}
\begin{itemize}
\item Utiliser l'impératif pour les conseils : \eng{arrive}, \eng{dress}, \eng{speak}, \eng{stay}.
\item Employer les collocations exactes : \eng{to be shortlisted}, \eng{to make a good impression}, \eng{to sell yourself}, \eng{on time}.
\item Ne pas écrire son vrai nom : signer « Manga ».
\item Éviter les calques : pas de \eng{pass the interview} ni de \eng{motivation letter}.
\item Compter les mots : rester entre 80 et 100.
\end{itemize}
\end{corrige}

\newpage

\coursec{Fiche bilan}

\begin{popfigure}
\begin{tikzpicture}[mindmap, grow cyclic, every node/.style={concept, align=center, font=\small},
  root concept/.append style={concept color=popblue, font=\bfseries, text=white},
  level 1/.append style={level distance=3.4cm, sibling angle=60},
  level 2/.append style={level distance=2.2cm, font=\scriptsize}]
\node[root concept, align=center]{Job\\Interview\\Vocabulary}
  child[concept color=poporange]{ node[align=center]{Stages\\of the interview}
    child { node[concept color=poporangeL, align=center]{apply, shortlist,\\interview, hire} }
  }
  child[concept color=popgreen]{ node[align=center]{People \&\\documents}
    child { node[concept color=popgreenL, align=center]{applicant,\\interviewer, CV,\\cover letter} }
  }
  child[concept color=poppurple]{ node[align=center]{Collocations\\\& expressions}
    child { node[concept color=poppurpleL, align=center]{attend an interview,\\make a good\\impression} }
  }
  child[concept color=poppink]{ node[align=center]{False\\friends}
    child { node[concept color=poppinkL, align=center]{formation $\neq$\\training,\\candidat = candidate} }
  }
  child[concept color=popgold]{ node[align=center]{Word families\\pronunciation}
    child { node[concept color=popgoldL, align=center]{employ, employer,\\employee,\\employment} }
  }
  child[concept color=popblue!70]{ node[align=center]{Vocabulary\\in context}
    child { node[concept color=popblueL, align=center]{model dialogue,\\role play} }
  };
\end{tikzpicture}
\legende{Carte mentale : le vocabulaire de l'entretien d'embauche}
\end{popfigure}

\begin{bilan}
\textbf{1. Lexique clé (English -- Français)}
\begin{itemize}
\item \eng{a job interview} : un entretien d'embauche ; \eng{to undergo / to attend an interview} : passer un entretien
\item \eng{an applicant / a candidate} : un candidat ; \eng{the interviewer / the panel} : le recruteur / le jury
\item \eng{a CV (curriculum vitae) / a résumé} : un CV ; \eng{a cover letter} : une lettre de motivation
\item \eng{a vacancy / an opening} : un poste vacant ; \eng{to apply for a job} : postuler à un emploi
\item \eng{to be shortlisted} : être présélectionné ; \eng{to be hired / to get the job} : être embauché
\item \eng{qualifications} : diplômes ; \eng{work experience} : expérience professionnelle
\item \eng{strengths and weaknesses} : points forts et points faibles ; \eng{a reference} : une référence
\item \eng{salary / wages} : salaire ; \eng{full-time / part-time} : à temps plein / à temps partiel
\end{itemize}

\textbf{2. Collocations et expressions à connaître par cœur}
\begin{center}
\begin{tabularx}{\linewidth}{|X|X|}
\hline
\rowcolor{popblueL} English & Français \\
\hline
to make a good (first) impression & faire bonne (première) impression \\
\hline
to dress smartly / formally & s'habiller de façon soignée \\
\hline
to answer questions confidently & répondre avec assurance \\
\hline
to shake hands firmly & serrer la main avec fermeté \\
\hline
to be suitable for the post & convenir au poste \\
\hline
“We will get back to you soon.” & « Nous vous recontacterons bientôt. » \\
\hline
\end{tabularx}
\end{center}

\textbf{3. Faux amis et pièges}
\begin{itemize}
\item \eng{formation} (anglais) = formation (géologique, militaire) ; « formation professionnelle » se dit \eng{training}.
\item « candidat » = \eng{candidate / applicant}, jamais \eng{candidat}.
\item \eng{actually} = en fait (pas « actuellement » = \eng{currently}).
\item \eng{a délai} = \eng{a deadline / a time limit} ; \eng{delay} = retard.
\item « passer un entretien » = \eng{to attend / to have an interview} (pas \eng{to pass an interview} : \eng{pass} = réussir).
\end{itemize}

\textbf{4. Familles de mots et prononciation}
\begin{itemize}
\item \eng{employ} (verbe) $\rightarrow$ \eng{employer} (employeur), \eng{employee} (employé), \eng{employment / unemployment} (nom).
\item \eng{apply} $\rightarrow$ \eng{application}, \eng{applicant} ; \eng{qualify} $\rightarrow$ \eng{qualification}.
\item Prononciation : \eng{interview} « ine-teur-viou », \eng{CV} « si-vi », \eng{employee} « im-plo-yi », \eng{suitable} « sou-teu-beul ».
\end{itemize}

\textbf{5. Méthode express pour un jeu de rôle d'entretien}
\begin{enumerate}
\item Saluer et se présenter : \eng{Good morning, my name is...}
\item Parler de ses diplômes et de son expérience : \eng{I hold a degree in... / I have worked as...}
\item Donner ses qualités avec des exemples concrets.
\item Poser une question sur le poste, remercier et prendre congé : \eng{Thank you for your time.}
\end{enumerate}
\end{bilan}

\begin{aretenir}[Checklist avant le Bac]
\begin{itemize}
\item[$\square$] Je connais les étapes d'un recrutement : \eng{apply -- shortlist -- interview -- hire}.
\item[$\square$] Je distingue \eng{employer} (employeur) et \eng{employee} (employé).
\item[$\square$] Je sais que « lettre de motivation » se dit \eng{cover letter}.
\item[$\square$] Je dis \eng{to apply for a job} (avec \eng{for}) et non \eng{to apply a job}.
\item[$\square$] Je traduis « passer un entretien » par \eng{to attend an interview}, jamais \eng{to pass}.
\item[$\square$] J'évite le faux ami \eng{formation} : je dis \eng{training}.
\item[$\square$] Je sais utiliser \eng{candidate} et \eng{applicant} pour « candidat ».
\item[$\square$] Je connais les collocations : \eng{make a good impression}, \eng{dress smartly}, \eng{shake hands}.
\item[$\square$] Je sais parler de mes \eng{strengths and weaknesses} en anglais.
\item[$\square$] Je prononce correctement \eng{interview}, \eng{CV} et \eng{employee}.
\item[$\square$] Je sais terminer un entretien poliment : \eng{Thank you for your time. I look forward to hearing from you.}
\item[$\square$] Je peux tenir un dialogue complet d'entretien d'embauche à l'oral et à l'écrit.
\end{itemize}
\end{aretenir}

\findecours

\end{document}
